% !TeX encoding = UTF-8
% !TeX spellcheck = <none>
%%%%%%%%%%%%%%%%%%%%%%%%%%%%%%%%%%%%%%%%%%%%%%%%%%%%%%%%%%%
%  _____           _____ 
% |_   _|   /\    / ____|
%   | |    /  \  | (___  
%   | |   / /\ \  \___ \ 
%  _| |_ / ____ \ ____) |
% |_____/_/    \_\_____/ 
%                        
%%%%%%%%%%%%%%%%%%%%%%%%%%%%%%%%%%%%%%%%%%%%%%%%%%%%%%%%%%%%
% This the main document of the diploma thesis template
% Created by Josua Arndt 29.04.2016
% TEMPLATE VERSION 1.5
% TEMPLATE DATE 30.11.2017
%
% - Do not change anything in this document without the permission of your supervisor
% - Compile the document for further information
% 
% For questions and comments contact johannes.bastl@ias.rwth-aachen.de
%
%%%%%%%%%%%%%%%%%%%%%%%%%%%%%%%%%%%%%%%%%%%%%%%%%%%%%%%%%%%%


%\RequirePackage[log, major-version]{snapshot} % This will produce a dependency list in a file \jobname.dep.
% checks if all packages are available and are in the same major-version as listet 
%\RequireVersions[]{ 
%	% here should be the input of the original_A_Template_BA_MA.dep_x
%	}
% nag check if you use some deprecated packages
\RequirePackage[l2tabu, orthodox]{nag}
%%%%%%%%%%%%%%%%%%%%%%%%%%%%%%%%%%%%%%%%%%%%%%%%%%%%%%%%%%%%
% SUPPRESSED WARNINGS %
%%%%%%%%%%%%%%%%%%%%%%%%%%%%%%%%%%%%%%%%%%%%%%%%%%%%%%%%%%%%
% Pdf files from MS Office Products with misleading page group
% pdfTeX warning: pdflatex.exe (file <FILENAME>.pdf): PDF inclusion: multiple pdfs with page group included in a single page>]
\pdfsuppresswarningpagegroup=1  

% Use KOMA-Skript scrbook for documentclass produces nice PDFs
\documentclass[
% draft, % set all packages to draft mode for fast compiling
fontsize=12pt, 		% Fontsize
open=right, 		% the Chapter begins on right side, a Blank page will be included if neccesary  
toc=bibliography, 	% Include Bibliography to Table Of Contents (TOC)
toc=listof, 		% Include list o figures to Table Of Contents
headsepline, 		% A Line under the Headers
%headinclude=false, 	% Headers are calculated as part of the Textarea 
%footinclude=false, 	% Footers are calculated as part of the Textarea 
cleardoublepage=empty, % With this option, you may define the page style of the interleaf pages created by the \cleardoublepage to break until the wanted page. empty adds Empty Pages before Next chapter if needed. So a Blank Page, no Header, No Footer, No Page Number. \cleardoublepage is called automatically after each chapter to flush pending floating Objects.
numbers=noendperiod, % No Period after the Number oin the Header
chapterprefix = false, % don´t Use Chpater X before the Chapter Titel
% scrbook uses "flushbottom " because in double-sided documents, it is preferred to have the same visual baseline, opposite is "raggedbottom"
parskip=half, % Makes some withspaces between Paragrafs and no ident
usegeometry, % use package geometry to change page layout
]{scrbook}
% change linespacing 
\usepackage[onehalfspacing]{setspace}
\usepackage{enumitem}
\setcounter{secnumdepth}{3} % Number subsubsection depth

% ##################################################################################
% # KOMA Hack to implement some newer feature to Packages wich does't use them     #
% # In this Document until now the warning of listings that doesn't use 'tocbasic' #
% ##################################################################################
\usepackage{scrhack}

% #######################################
% # Header and Footer                   #
% #######################################
% Komascript substitute for Fancyheadings Paket
\usepackage[markcase=used,automark,autooneside=false,headsepline]{scrlayer-scrpage} %\\

% ##################################################
% # Input encoding, Output encoding, File Encoding #
% ##################################################
%
% !!! Attention !!! Make sure your file Encoding is UTF8
%
% The microtype package improves the spacing between words and letters. It does a lot more and most % people won`t notice the difference. But still, the resulting document will be easier to read and % looks better when microtype is loaded. Load this package after fonts, if any, as the package  
% behavior is dependent on this font.
\AtBeginDocument{%
	\IfFileExists{type1ec.sty}
	{%
		\usepackage{type1ec} % required to use micotype
		\RequirePackage{microtype}%
	}
	{%
		\PackageWarning{Josua Arndt}{%
			\MessageBreak Package cm-super not available, therefore microtype not used.\MessageBreak
			Install cm-super to get better text layout.\MessageBreak
		}%
	}%
}

% The two packages are not connected, though it is best to call fontenc first and then inputenc
% 1. fontenc is oriented to output, that is, what fonts to use for printing characters.
\usepackage[T1]{fontenc} % Character Coding set to Type1
% 2. Inputenc for Input encoding, allows the user to input accented, as ä,ö,ü,é  characters directly from the keyboard
% Josua old = latin = ISO-8859-1. will make Problems 
% change to utf8 = utf8 encoding for inputed charactes from keyboard.
\usepackage[utf8]{inputenc} % 
%\usepackage[latin1]{inputenc}


% #######################################
% # Conditionals                   		#
% #######################################
%% Do not change the Conditionals here, 
% change them in the Settings.tex file
\usepackage{etoolbox}
\newbool{PrintVersion}
\newbool{NoLogo}
\newbool{IsEnglish} 
\newbool{ConfidentialVersion}
\newbool{secondSupervisor}

\newcommand{\confidential}[2]{ \ifbool{ConfidentialVersion}{#1}{#2}} 

%%%%%%%%%% MODIFY THE SETTINGS IN THIS FILE %%%%%%%%%%%%%%%%%%%%%%%%%%%%%%%%%%%
%TODO SETTINGS change the settings in the settings file
% !TeX spellcheck = en_US
% !TeX encoding = UTF-8
%% Change the settings of the Template in this file

% \setboolean{DraftMode}{false} % no Effect
% To set Packages to draft mode go to A_Template_BA_MA.tex 
% and set document options to draft. Makes Compiling much faster. 

% change Links and IAS Logo to Black and white Printable output.
% also set binding correction 
%TODO when printing change this
\setbool{PrintVersion}{false} 
% adds CONFIDANTIAL Print to Bottom line
\setbool{ConfidentialVersion}{false}
% Changes TOC, ABR, LOF and most parts of the document to english
\setbool{IsEnglish}{true} 

% Changes First Page to use No Logo
\setbool{NoLogo}{false}

\newcommand{\Qualifikation}{Qualifikation}
\newcommand{\IASAuthor}{Josua Arndt}
\newcommand{\AuthorAus}{Worms}
\newcommand{\ThesisTitleGerman}{{Deutscher Titel, kann auch lang sein um zu testen ob die nächste Zeile gut aussieht! Auch wenn es mehr als zwei Zeilen sind. Oder Drei, möglicherweise auch Vier oder sogar Fünf. Sollten SI Einheiten notweindig sein nutze textopdfstring sonst gibt es ein Fehler \texorpdfstring{\SI{455}{\kilo\hertz}}{455kHz} $\mathbf{\Sigma}$  \boldmath$\Sigma$  }}
\newcommand{\ThesisTitleEnglish}{English Title, could also be very long to test if the next line looks good! Also if there are more then two lines. If you need SI units in in your Title use texorpdfstring else you get an error \texorpdfstring{\SI{455}{\kilo\hertz}}{455kHz} }

\newcommand{\IASSubject}{Subject of Work} % Bachelorarbeit, Studienarbeit, Master Thesis and so on.

\newcommand{\IASMatNr}{123456} % Matrikel Number
\newcommand{\IASSupervisor}{BetreuerVorname BetreuerNachname, M.Sc.} % Supervisor " Family Name, First Name Qualifikation][M.Sc.]"
\setbool{secondSupervisor}{true} 
\newcommand{\IASSupervisortwo}{BetreuerVorname BetreuerNachname, M.Sc.}
\newcommand{\IASNumber}{D123} % Number of the Work, given by your Supervisor and the secretary.

\newcommand{\FirstPrufer}{First Prufer}
\newcommand{\SecondPrufer}{Second Prufer}
\newcommand{\IASSubmissionDate}{dd.mm.yyyy}

% Path to your presentations will be added to the Presentation slides appendix
% TODO change to your ppt presentations
\newcommand{\firstPresentation} {./PowerPoint/IAS_HowToMakeNicePlots_Latex&PowerPoint} 
% TODO remove comment for master thesis second presentation
%\newcommand{\secondPresentation} {./PowerPoint/IAS_HowToMakeNicePlots_Latex&PowerPoint} 
\newcommand{\finalPresentation} {./PowerPoint/IAS_HowToMakeNicePlots_Latex&PowerPoint} 

% You have to load your Bibliography file after the Biblatex packet



%\pdfminorversion=3			%Could help in case the printery is unable to print
%\pdfobjcompresslevel=0		%Could help in case the printery is unable to print
%%%%%%%%%%%%%%%%%%%%%%%%%%%%%%%%%%%%%%%%%%%%%%%%%%%%%%%%%%%%%%%%%%%%%%%%%%%%%%%
%TODO COMMENT OUT the next line. It forces the compilation to not work and shows if you worked throhgt the TODOs and the README_Template PDF.
%\usepackage{iHaveNotWorkThroughTheTodos}

% ##################################################################
% # Configuration of Language 									   #
% ##################################################################
\ifbool{IsEnglish}
{
	\usepackage[english]{babel}
	\selectlanguage{english}
}
{
	\usepackage[ngerman]{babel}
	\selectlanguage{ngerman}
}


%This package provides advanced facilities for inline and display quotations. It is designed for a wide range of tasks ranging from the most simple applications to the more complex demands of formal quotations. The facilities include commands, environments, and userdefinable ‘smart quotes’ which dynamically adjust to their context. Quotation marks are switched automatically if quotations are nested and they can be adjusted to the current language if the babel package is available. There are additional facilities designed to cope with the more specific demands of academic writing, especially in the humanities and the social sciences. All quote styles as well as the optional active quotes are freely configurable.
\usepackage{csquotes} 

% ################################################################
% # Configuration of Colors and Colors of Links for data in PDF  #
% ################################################################
%COLORS 68 dvips colors, see https://en.wikibooks.org/wiki/LaTeX/Colors
\usepackage[pdftex,table,dvipsnames]{xcolor}
% Changed color of links for the online version

% Colors for the links
\definecolor{blueLink}{cmyk}{0.965,0.601,0,0.322}
\definecolor{lightBlueLink}{cmyk}{0.727,0.455,0,0.267}
\definecolor{darkBlueLink}{cmyk}{0.922,0.922,0,0.498}
\definecolor{lightGray}{cmyk}{0,0,0,0.627}


%COLORS use this
% Some Definitions for plot colors used in the 
% Thesis of IAS Students so they look the same
\definecolor{IASblue}{cmyk}{7,0.493,0,0.412}
\definecolor{IASred}{cmyk}{0,0.705,0.75,0.137}
\definecolor{IASgreen}{cmyk}{0.333,0,0.467,0.412}
\definecolor{IASorange}{cmyk}{0,0.451,1,0}
\definecolor{IASLogoBlue}{cmyk}{1,0.485,0,0.353}

%%%%%%%%%%%%%%%%%%%%%%%%%%%%%%%%%%%%%%%%%%%%%%%%%%
% COLORS RWTH Corporate Design 
% http://www9.rwth-aachen.de/go/id/bnxz the Link "Farben" download the PDF
\definecolor{rwthci-blue-cmyk}{cmyk}{1,0.5,0,0}
\definecolor{rwthci-petrol-cmyk}{cmyk}{1,0.3,0.5,0.3}
\definecolor{rwthci-turk-cmyk}{cmyk}{1,0,0.4,0}
\definecolor{rwthci-green-cmyk}{cmyk}{.7,0,1,0}
\definecolor{rwthci-maygreen-cmyk}{cmyk}{.35,0,1,0}
\definecolor{rwthci-orange-cmyk}{cmyk}{0,0.4,1,0}
\definecolor{rwthci-red-cmyk}{cmyk}{0.15,1,1,0}
\definecolor{rwthci-bordeaux-cmyk}{cmyk}{0.25,1,0.7,0.2}
\definecolor{rwthci-violet-cmyk}{cmyk}{0.7,1,0.35,0.15}
\definecolor{rwthci-purple-cmyk}{cmyk}{0.6,0.6,0,0}


% ################################################################
% # Graphics, Captions, Subfigures and sucaptions
% ################################################################
% used to include Graphics
\usepackage[pdftex]{graphicx} 
%TODO add ALL your graphics paths
% all path where graphics are located have to be inserted here so latex can find them.
% also the Satnalone files because they are included with includepdf
\graphicspath{ {Bilder/} {Bilder/SVG/}{Standalone/} } 
%test svn was ist da los!!
% allows to include eps graphix with includegraphics
\usepackage{epstopdf}  

% subfig colides with caption and produces errors in the lof and lot.
% only use supcaption instead of subfigure or subfig
\usepackage[%
%font={small,sf},
%labelfont=bf,
%format=hang,    
%format=plain,
%margin=0pt,
%width=0.8\textwidth,
tableposition=top
]{caption}
\usepackage[list=true,labelformat=simple]{subcaption}
\captionsetup[table]{position=top}
\renewcommand\thesubfigure{(\alph{subfigure})}
\renewcommand\thesubtable{(\alph{subtable})}
% ##########################################################################
% # Package for nice Graphs 											   #
% ##########################################################################

\usepackage{pgfplots}% loads tikz which loads then loads xcolor and graphicx, to avoid option clash xcolor and graphicx have to be loaded prevously
\pgfplotsset{compat=newest} % use newest PGF plots Package. can be set to depending Version.
%compat=newest is not always the best practice since a newer version can break the functionality of %the code. Instead, using the version number that the code is created with, e.g., compat=1.8 should %make things smoother for debugging.

% This document has following sizes
% textwidth = 160mm        = 455.24945pt
% textheight= 207.89667 mm = 591.53027pt
\newcommand{\originalTextWidth}{455.24945pt}
\newcommand{\originalTextHeight}{591.53027pt}

% to place multiple plots side by side and stacked
\usepackage{calc}
\usepgfplotslibrary{groupplots}
\usetikzlibrary{calc,matrix,positioning}
% Option that should not be change by the user, this are set here so it is obvious and makes the code more readable


%\pgfplotsset{
%	every axis/.append style={
%		width=0.75*\originalTextWidth, 
%		height=0.1*\originalTextHeight,
%		scale only axis, % the width and height specifies only axis not the labels etc. to avoid missplacing of the different plots when Labels are used or not.
%		xlabel style={yshift=-.0em, inner ysep=0pt} % used to properly crop off the plot so the distance between x-label and Caption is not to big.
%	},
%}
%\tikzset{every picture/.style={trim axis left, trim axis right}}% used to avoid missplacing of different plots when the numbers on the x-axis stands out left or right.


\ifbool{PrintVersion} %
{% Print Version - Black and white

	\pgfplotscreateplotcyclelist{RWTH_CI_cyclelistlist}{%
		solid, every mark/.append style={solid, fill=gray}, mark=*\\%
		dotted, every mark/.append style={solid, fill=gray}, mark=square*\\%
		densely dotted, every mark/.append style={solid, fill=gray}, mark=otimes*\\%
		loosely dotted, every mark/.append style={solid, fill=gray}, mark=triangle*\\%
		dashed, every mark/.append style={solid, fill=gray},mark=diamond*\\%
		loosely dashed, every mark/.append style={solid, fill=gray},mark=*\\%
		densely dashed, every mark/.append style={solid, fill=gray},mark=square*\\%
		dashdotted, every mark/.append style={solid, fill=gray},mark=otimes*\\%
		dasdotdotted, every mark/.append style={solid},mark=star\\%
		densely dashdotted,every mark/.append style={solid, fill=gray},mark=diamond*\\%
	}% can now be used throughout the document

}%
{%  Online Version - Colors
	% there are also predefined cyclelists
	\pgfplotscreateplotcyclelist{RWTH_CI_cyclelistlist}{
		{color=rwthci-blue-cmyk, solid},
		{color=rwthci-petrol-cmyk, dashdotted},
		{color=rwthci-turk-cmyk, solid},
		{color=rwthci-green-cmyk, dashdotted},
		{color=rwthci-maygreen-cmyk, solid},
		{color=rwthci-orange-cmyk, dashdotted},
		{color=rwthci-red-cmyk, solid},
		{color=rwthci-bordeaux-cmyk, dashdotted},
		{color=rwthci-violet-cmyk, solid},
		{color=rwthci-purple-cmyk,	dashdotted}% error in older version of pgf if last comma is not set it will not work if last comma is set one line will be black !!!!!
	} % can now be used throughout the document
}


% ##########################################################################
% Some necessary Packages
% ##########################################################################
\usepackage{pdfpages} % To include whole PDF-Pages e.g in the appendix
\usepackage{url} % for URLs in the Bibliography


\usepackage{textcomp}% To avoid messing up anything else in a document, siunitx is set up to avoid loading font packages and the like. As such, it uses its 'own' version of the \textminus and \textmu, and sets up a minimal amount of support for that. However, that confuses microtype. Loading textcomp 'fixes' this as siunitx then uses the textcomp version of the symbols, and microtype knows what is going on with that. I do recommend loading textcomp with siunitx, but it's not absolutely required so it's not enforced.

\usepackage[ % To include SI_unit Support
	binary-units=true % units of bits and bytes
	% use abbreviations \fg for femtogram is displayes fg or microgram \ug μg
]{siunitx} 

\ifbool{IsEnglish}
{\sisetup{output-decimal-marker = {.},per-mode = symbol, detect-all}}
{\sisetup{output-decimal-marker = {,},  per-mode = symbol,  detect-all}}

%ToDO define new Units
\DeclareSIUnit\decibelmW{dBm}

% ##################################################################
% # Configuration of Bibliography and source file for Bibliography #
% ##################################################################
% bibliography
% biblatex used with biber backend
% You have tu run Build, then Biber under Tools/Command/Biber, then again Build to genrerate the Bibliography
%
% There is a Profile for TexStudio in the folder which configures the user Macros in Tools/User/makeGlsAcr to build both.
% The "Build & View" Button is defined to run userCommand0 which makes both everytime you compile.
% The  "Build & View" Button is also configured to run biber everytime you compile.
% You have to load the profil under Options/Load Profile.
\usepackage[
	backend=biber, 		% use Biber as Backend
	%citestyle=authoryear,
	citestyle=numeric,	% numerik, list references as number
	%style=alphabetic,	% style = alphabetic , eg Jones 1997, [Jo97], stlye = numeric eg : [1]
	% natbib=true,		% natbib = true or false , for compatibility with mcite comands
	% sorting=nty,		% sorting = nyt , sorting according to name, title and year and so on.
	sorting = none , 	% list as they apear in the document
	bibencoding=UTF8,   % set encoding to UTF8, make sure all your files are encoded with utf8
	url=false, 
	doi=true,
	eprint=false,
	]{biblatex}

% the following lines are only for this Template, make your one .bib file
% You have to load your Bibliography file after the Biblatex packet
% use \addbibresource{\jobname.bib} , \jobname calls the Bib which is named like the Main Tex file.
% but you can change \jobname to every name. eg myBib.bib
% do not change name to generic "\jobname" because it will then not be listed in the folder three at the left side
\addbibresource{references.bib}



% #############################################################################
% # Following different packages are loaded which can be used in the Document #
% #############################################################################
% Tabularx for table with 
\usepackage{tabularx, tabulary, longtable, ltxtable, booktabs}
% solve problem with all tabular packages which read its contents twice for calculating the width but does not reset the line counter.
\let\Tabular\tabular
\def\tabular{\global\rownum=0\relax\Tabular}

\usepackage{xfrac}
%This will allow you to change the row spacing to your liking by invoking the ra{} command. I found that 1.3 was sufficient.
%\newcommand{ra}[1]{renewcommand{arraystretch}{#1}}
%\renewcommand{\arraystretch}{1.3}



% For Equations
\usepackage{amsmath}
% needeed for equation in IEEEeqnarray
\usepackage{IEEEtrantools}

% ##########################################################################
% # Configurationthe Standalone Packet and setting some widths and heights #
% ##########################################################################
% use Package standalone to exclude tikzPictures, tables, etc. in seperated files
% and avoid to rebuild them every time.
% all Packages which are used in the SubDucuments have to be included in this header also.
\usepackage{standalone}
% if you change the Setting for A4 to A5 or vice versa, set theconfiguration for the mode of standalone to "build"
% To use the variabl tikzAxisWidth use tex for standolen mode, so you can change the width of the axis with a global variable. 

\standaloneconfig{mode=buildnew}
% tex = use sourcefile, default
% image = use existing image file produced by the source file
% build = build image from source, then use it
% buildmissing= only build image if it does not exist
% buildnew = only build image if source file is newer 


%TODO COMMENT OUT the entire following block if SVG is not used.
% ###################################################
% #				SVG Image conversion				#
% ###################################################
\usepackage{svg}
%TODO SVG set path to SVG Images
\setsvg{svgpath={"Bilder/SVG/"}}

%TODO SVG IF you compile an your own PC, set the Path to inkscape to the environment path or change the inkscape call in the next macros.
% new command for extractiong pdfs out of svg without converting the Text to Latex text.
\newcommand{\executeiffilenewer}[3]{%
	\ifnum\pdfstrcmp{\pdffilemoddate{#1}}%
	{\pdffilemoddate{#2}}>0%
	{\immediate\write18{#3}}\fi%
}
\newcommand{\includesvgNoLatex}[2][1]{%
	\executeiffilenewer{#2.svg}{#2.pdf}{%
		inkscape.exe -z -C --file="#2.svg" --export-pdf="#2.pdf" %--export-latex
	}%
	\ifx\relax#1\relax% if empty rela=relax 
	{\includegraphics[ page=1]{#2.pdf}}
	\else{\includegraphics[ #1, page=1]{#2.pdf}}\fi%
	
}%

% new command for extractiong pdfs out of svg without converting the Text to Latex text.
\newcommand{\convertPdfToPpm}[1]{%
	\executeiffilenewer{#1.pdf}{#1.ppm}{%
		pdftoppm -r "600" -singlefile -aa "yes" -aaVector "yes" "#1.pdf"  "#1"%
	}%
}%
\newcommand{\convertPpmToPng}[1]{%
	\executeiffilenewer{#1.ppm}{#1.png}{%
		magick "#1.ppm" -background white -flatten -alpha remove "#1.png"%
	}%
}%
\newcommand{\includePdfToPng}[2][1]{%
	\convertPdfToPpm{#2}%
	\convertPpmToPng{#2}%
	
	\ifx\relax#1\relax% if empty rela=relax 
	{\includegraphics[page=1]{#2.pdf}}
	\else{\includegraphics[ #1, page=1]{#2.png}}\fi%
}%
%% to add subfolder and avoid errors like
%%Package pdftex.def Error: File `arial_svg.pdf' not found. ...[width=\unitlength,page=1]{arial_svg.pdf}}
%\usepackage{xparse}
% ###################################################
% #				SVG Image conversion end			#
% ###################################################



% ##########################################################################
%Listings
% ##########################################################################
\usepackage{listings} % To show code in the Text
\renewcommand\lstlistingname{Code}
\renewcommand{\lstlistlistingname}{List of Source \lstlistingname}

\makeatletter
\let\old@lstKV@SwitchCases\lstKV@SwitchCases
\def\lstKV@SwitchCases#1#2#3{}
\makeatother
\usepackage{lstlinebgrd}
\makeatletter
\let\lstKV@SwitchCases\old@lstKV@SwitchCases

\lst@Key{numbers}{none}{%
    \def\lst@PlaceNumber{\lst@linebgrd}%
    \lstKV@SwitchCases{#1}%
    {none:\\%
     left:\def\lst@PlaceNumber{\llap{\normalfont
                \lst@numberstyle{\thelstnumber}\kern\lst@numbersep}\lst@linebgrd}\\%
     right:\def\lst@PlaceNumber{\rlap{\normalfont
                \kern\linewidth \kern\lst@numbersep
                \lst@numberstyle{\thelstnumber}}\lst@linebgrd}%
    }{\PackageError{Listings}{Numbers #1 unknown}\@ehc}}
\makeatother

\lstdefinelanguage{SKILL}{
	morekeywords=[0]{% SKILL Keywords from Finder
		DATA, ERC, IDC, IF, IS, IT, LVS, MP, NG, OP, OPT, SIMULATOR, VAR, VDC, VF, VN, VN2, VNP, VNPP, VPD, VS, VT, abs, abs, absAbgen, absAbstract, absAddPin, absAdvancedCreateAbstractSubblock, absAttachTechLib, absCopyBinOptions, absDeleteBin, absDeleteBinMoveCellsTo, absDeselectAllBins, absDeselectBin, absDeselectBinFrom, absDeselectCell, absDeselectCellFrom, absDeselectCells, absDisableUpdate, absDistributeCells, absEnableUpdate, absExecDelPin, absExecLayoutPrep, absExit, absExportLEF, absExportOptions, absExportReport, absExtract, absGetBinOption, absGetCellProp, absGetLibrary, absGetOption, absGetTerminalProp, absImportCTLF, absImportDEF, absImportGDS, absImportLEF, absImportLogical, absImportVerilog, absIsPointInsideFig, absModifyPinProp, absModifyTerminalProp, absMoveSelectedCellsToBin, absNewBin, absOpenAdvancedAbstract, absPins, absRemoveFloatTermsInCV, absRemoveFloatTermsInLib, absRenameBin, absSelect, absSelectAllBins, absSelectBin, absSelectBinFrom, absSelectCell, absSelectCellFrom, absSelectCells, absSelectCellsInList, absSetBinOption, absSetCellList, absSetCellProp, absSetCellProp, absSetGroundSignal, absSetLayerMinPinSize, absSetLibrary, absSetOption, absSetPowerSignal, absSetSnapToGrid, absSetTerminalProp, absSetTextMap, absSlice, absSort, absTextCharRemove, absTextTrunc, absVerify, absVersion, ac, acos, acos, acvMeasTim, acvMeasTimpwr, add1, add1, addCheck, addDefstructClass, addPoint, addSubwindow, addSubwindowTitle, addTitle, addWaveLabel, addWindowLabel, adpAddInst, adpAddTerm, adpAreaEst, adpBlockAnalysis, adpBlockAnalysisReport, adpCheckRgnOverlap, adpChipAnalysis, adpChipAnalysisReport, adpCmdAlignBlocks, adpCmdAnalyze, adpCmdChop, adpCmdDelete, adpCmdMemberAdd, adpCmdMemberFit, adpCmdMemberSub, adpCmdMemberUnfit, adpCmdMove, adpCmdMoveOptions, adpCmdSnapModeAdd, adpComputeDelaysMC, adpConfigEdit, adpConfigRead, adpConfigReadCB, adpConfigWrite, adpConnect, adpCreateLSWWidgets, adpCreateSearchWidget, adpDelayReportMC, adpDeleteInst, adpDesignBoundary, adpDisconnect, adpECOMC, adpEcoCellView, adpEditSelSet, adpEnvInitEnv, adpEstAllowedAreaUtil, adpEstAllowedGateUtil, adpExtractRCMC, adpFlatten, adpGGrid, adpGroupAnalysis, adpGroupAnalysisReport, adpGroupAreaEst, adpGroupSimpleAreaEst, adpHiEditDisplayOptions, adpHiPDEFRead, adpHiPDEFWrite, adpHierGen, adpIOPinConstraint, adpIOPinDef, adpIOPinStatus, adpIORandomPlaceGE, adpIOSideConstraint, adpInfoFormCB, adpInitDisplay, adpInitFloorplan, adpInitProc, adpMGrid, adpMacroConstraint, adpMacroStatus, adpModifyInst, adpNetConstraint, adpPGrid, adpPadConstraint, adpPadStatus, adpPartition, adpPartitionConstrMC, adpRGrid, adpRegion, adpRegionAnalysis, adpRegionAnalysisReport, adpReportFormatReadCB, adpReportFormatWriteCB, adpRgnSimpleAreaEst, adpRow, adpSDFWriteDelaysCB, adpShpToSymPin, adpSimpleAreaEst, adpSoftGroup, adpStdCellConstraint, adpStdCellStatus, adpbrDgnBrMC, adpbrEstimateAreaForAutoAbstract, adpbrFitMC, adpbrPropMastMC, adpbrRegisterAreaEstimation, adpbrSetMasterPropForSoftBlk, adpbrUnRegisterAreaEstimation, adpbrZoomInMC, adpfpAnalyzeCERegionMC, adpfpCheckInstOverlap, adpfpCheckRegionMC, adpfpCheckRgnOverlap, adpfpClearErrorMC, adpfpCreateGroupMC, adpfpCreateRgnMC, adpfpInitFPMC, adpfpLoadFile, adpfpLoadFileMC, adpfpMoveCellToRow, adpfpSaveFile, adpfpSaveFileMC, adpfpSoftPinOptimizeMC, adphmBreakBlock, adphmBreakBlockMC, adphmCreateBlkAbsMC, adphmCreateBlock, adphmCreateBlockMC, adphmCreateUpdateAutoAbstract, adphmCreateUpdateAutoAbstractMC, adphmGrpToBlock, adphmInstToBlock, adphmRgnToBlock, adphmTuneParamCB, adphmUpdateAutoLayoutMC, adpnlSetEnv, adpnlSymToAbst, adppaEstNetDelay, adppaEstPathDelay, adppaGenRCReport, adppaGenRSPF, adppaGenSDF, adppaLoadSDF, aedCopyCDF, aedDeleteCDF, aedEditCDF, aelCheckRange, aelDisplayOPParam, aelEngNotation, aelEnvCompile, aelEnvCreate, aelEnvExecute, aelEnvFreeCompExpr, aelEnvGetErrStr, aelEnvGetGlobal, aelEnvInterpret, aelEnvListDeferredFuncs, aelEnvListDeferredGlobals, aelEnvListExprFuncs, aelEnvListExprGlobals, aelEnvListFuncs, aelEnvListGlobals, aelEnvListGlobalsValues, aelEnvName, aelEnvSetGlobalList, aelEnvSetGlobals, aelGetSignifDigits, aelNumber, aelPopSignifDigits, aelPushSignifDigits, aelSetLineage, aelSignum, aelStrDblNotation, aelSuffixNotation, aelSuffixWithUnits, aelSumOPParam, ahdlUpdateViewInfo, ahdlUpdateViewInfo, alias, almGetModuleName, almGetNamePrefix, almGetNetlistProcedure, almGetNetlistType, almGetOpPointParamMap, almGetOtherParameterList, almGetParameterList, almGetPropMappingList, almGetStringParameterList, almGetTerminalList, almGetTerminalMap, almGetViewInfoNameList, almHasViewInformation, almSetModuleName, almSetNamePrefix, almSetNetlistProcedure, almSetOpPointParamMap, almSetOtherParameterList, almSetParameterList, almSetPropMappingList, almSetStringParameterList, almSetTerminalList, almSetTerminalMap, alphaNumCmp, alphalessp, amcGetNumBins, amcGetSamplingMethod, amcSetNumBins, amcSetSamplingMethod, amsCheckCV, amsError, amsGetInstanceName, amsGetNetlister, amsGetPortExpr, amsGetUniqueName, amsInfo, amsIsPresent, amsMapName, amsMtlinePrintParams, amsMtlineTermOrder, amsNetlist, amsNportTermOrder, amsPrint, amsPrintAlias, amsPrintAliases, amsPrintAttribute, amsPrintAttributes, amsPrintIO, amsPrintIOs, amsPrintInstance, amsPrintInstanceMasterName, amsPrintInstanceParameter, amsPrintInstanceParameters, amsPrintInstancePorts, amsPrintParameter, amsPrintParameters, amsPrintPort, amsPrintPorts, amsPrintWire, amsPrintWires, amsProcessCellViews, amsSpectreToVams, amsUIOptionsForm, amsUIRunNetlisterForm, amsWarning, analysis, ancAdjustNameCase, ancGetSimInstName, and, ansAnalogRegCDFsimInfo, ansCdlCompPrim, ansCdlGetMultiplicity, ansCdlGetSegmentConnections, ansCdlGetSegmentInfo, ansCdlGetSegmentInstParams, ansCdlGetSimPropValue, ansCdlHnlPrintInst, ansCdlPrintConnections, ansCdlPrintInheritedParams, ansCdlPrintInstName, ansCdlPrintInstParams, ansCdlPrintInstProps, ansCdlPrintModelName, ansCdlPrintModuleName, ansCdlPrintString, apaStop, apaStopCB, appTrigger, append, append1, appendPath, appendWaves, apply, applyEnterFun, argc, arglist, argv, armGetBrowser, armSetCalc, arrayp, arrayref, artBlankString, artClose, artCloseAllFiles, artCurrentInstSimName, artFlush, artFprintf, artGenerateHierSymbolCDF, artGetCdfTargetCV, artGetCellViewDesignVarList, artInvalidateAmap, artListOpenFiles, artListToWaveform, artMakeString, artOutfile, asBlockPlowDown, asBlockPlowLeft, asBlockPlowRight, asBlockPlowUp, asBlockPush, asMoveInst, asMoveSideSoftBlock, asiAddAnalysis, asiAddAnalysisField, asiAddAnalysisOption, asiAddDesignVarList, asiAddEnvOption, asiAddFlowchartLink, asiAddFlowchartStep, asiAddKeepOption, asiAddOceanAlias, asiAddSimOption, asiAddVerilogArgs, asiAnalogAutoloadProc, asiChangeAnalysis, asiChangeAnalysisField, asiChangeAnalysisOption, asiChangeAnalysisOptionFormProperties, asiChangeEnvOption, asiChangeEnvOptionFormProperties, asiChangeFlowchartStep, asiChangeKeepOption, asiChangeKeepOptionFormProperties, asiChangeSimOption, asiChangeSimOptionFormProperties, asiCheck, asiCheckAcEnabledWhenNoiseEnabled, asiCheckAnalysis, asiCheckBlank, asiCheckBlankInstExists, asiCheckBlankNetExists, asiCheckBlankNumeric, asiCheckBlankNumericGeq, asiCheckBlankNumericGreater, asiCheckBlankNumericLeq, asiCheckBlankNumericNequal, asiCheckDesignVariable, asiCheckDigitalSimulationSuccess, asiCheckExpression, asiCheckExpressionGreater, asiCheckMultipleGreater, asiCheckSimulationSuccess, asiCloseOceanScript, asiConstructDigitalNetlist, asiCornerSimCB, asiCreateAnalysisField, asiCreateAnalysisOption, asiCreateCdsenvFile, asiCreateFlowchart, asiCreateFormatter, asiCreateLogFileVerilog, asiDcStore, asiDefineDataAccessFunction, asiDefineDataMappingFunction, asiDeleteAnalysis, asiDeleteAnalysisField, asiDeleteAnalysisOption, asiDeleteEnvOption, asiDeleteFlowchartLink, asiDeleteFlowchartStep, asiDeleteKeepOption, asiDeleteSimOption, asiDisableAnalysis, asiDisplayAnalysis, asiDisplayAnalysisField, asiDisplayAnalysisOption, asiDisplayAnalysisOptionFormProperties, asiDisplayEnvOption, asiDisplayEnvOptionFormProperties, asiDisplayFlowchart, asiDisplayKeepOption, asiDisplayKeepOptionFormProperties, asiDisplayNetlistOption, asiDisplaySimOption, asiDisplaySimOptionFormProperties, asiEditDigitalStimulus, asiEnableAnalysis, asiExecuteFlowchart, asiFinalNetlist, asiFormatAnalysis, asiFormatAnalysisList, asiFormatControlStmts, asiFormatDefinitionFileList, asiFormatDesignVarList, asiFormatGraphicalStimuli, asiFormatInitCond, asiFormatKeepList, asiFormatModelLibSelectionList, asiFormatNodeSet, asiFormatSimulatorOptions, asiFormatTextStimulusFileList, asiGetAnalogKeepList, asiGetAnalogRunDir, asiGetAnalogSimulator, asiGetAnalysis, asiGetAnalysisField, asiGetAnalysisFieldChoices, asiGetAnalysisFieldList, asiGetAnalysisFieldVal, asiGetAnalysisFormFieldChoices, asiGetAnalysisFormFieldVal, asiGetAnalysisName, asiGetAnalysisNameList, asiGetAnalysisOptionChoices, asiGetAnalysisOptionList, asiGetAnalysisOptionSendMethod, asiGetAnalysisOptionVal, asiGetAnalysisParamNameList, asiGetAnalysisSigList, asiGetAnalysisSimFieldList, asiGetAnalysisType, asiGetCalcResultsDir, asiGetCommandFooter, asiGetCommandLineOption, asiGetCornerDesignVarList, asiGetCornerList, asiGetCornerModelSelectionList, asiGetCornerModelingStyles, asiGetCornerName, asiGetCornerResultsDir, asiGetCornerTemperature, asiGetCurrentSession, asiGetDefinitionFileList, asiGetDesignCellName, asiGetDesignLibName, asiGetDesignVarList, asiGetDesignViewName, asiGetDigitalCommandLineOption, asiGetDigitalKeepList, asiGetDigitalNetlistFileName, asiGetDigitalRunDir, asiGetDigitalSimExecName, asiGetDigitalSimulator, asiGetDigitalSimulatorLogFileName, asiGetDigitalStimulusFileName, asiGetDrlData, asiGetEnabledAnalysisList, asiGetEnvOptionChoices, asiGetEnvOptionVal, asiGetFlowchart, asiGetFormFieldChoices, asiGetFormFieldVal, asiGetFormattedVal, asiGetFormatter, asiGetId, asiGetKeepList, asiGetKeepOptionChoices, asiGetKeepOptionVal, asiGetLogFileList, asiGetMarchList, asiGetModelLibFile, asiGetModelLibSection, asiGetModelLibSelectionList, asiGetNetlistDir, asiGetNetlistOption, asiGetNetworkId, asiGetOutputList, asiGetPlotList, asiGetPredefinedCommandLineOption, asiGetPsfDir, asiGetReservedWordList, asiGetSelObjName, asiGetSelObjType, asiGetSelObjValue, asiGetSession, asiGetSimExecName, asiGetSimInputFileList, asiGetSimInputFileName, asiGetSimInputFileSuffix, asiGetSimName, asiGetSimOptionChoices, asiGetSimOptionNameList, asiGetSimOptionSendMethod, asiGetSimOptionVal, asiGetSimulationRunCommand, asiGetSimulatorList, asiGetStimulusGlobals, asiGetStimulusInputs, asiGetTextStimulusFileList, asiGetTool, asiGetTopCellView, asiGetVerilogCommandLineOption, asiGetVerilogHost, asiGetVerilogHostMode, asiInit, asiInit, asiInit, asiInit, asiInit, asiInit, asiInit, asiInit, asiInitAnalysis, asiInitDataAccessFunction, asiInitEnvOption, asiInitMixedKeepOption, asiInitSimOption, asiInitVerilog, asiInitVerilogEnvOption, asiInitVerilogFNLEnvOption, asiInitVerilogHNLEnvOption, asiInitVerilogSimOption, asiInitialize, asiInitializeNetlisterMixed, asiInterruptSim, asiInvalidateControlStmts, asiInvalidateFlowchartStep, asiIsAnalysisEnabled, asiIsCaseSensitive, asiIsConfigDesign, asiMapInstanceName, asiMapNetName, asiMapOutputName, asiMapTerminalName, asiNeedSuffixEvaluation, asiNetlist, asiNetlistMixed, asiOpenOceanScript, asiPartitionDesign, asiPrepareDigitalSimulation, asiQuitSimulator, asiRawNetlist, asiRegisterTool, asiRunSimulation, asiSelectAnalysisCompParam, asiSelectAnalysisInst, asiSelectAnalysisNet, asiSelectAnalysisSource, asiSelectInst, asiSelectNet, asiSelectSourceInst, asiSelectTerm, asiSelectTermNet, asiSendAnalysis, asiSendControlStmts, asiSendDesignVars, asiSendInitCond, asiSendInitFile, asiSendKeepList, asiSendModelPath, asiSendNetlist, asiSendNodeSets, asiSendOptions, asiSendRestore, asiSendSim, asiSendUpdateFile, asiSetAnalogSimulator, asiSetAnalysisFieldChoices, asiSetAnalysisFieldVal, asiSetAnalysisFormFieldChoices, asiSetAnalysisFormFieldVal, asiSetAnalysisFormWidth, asiSetAnalysisOptionChoices, asiSetAnalysisOptionFormProperties, asiSetAnalysisOptionVal, asiSetDesignVarList, asiSetDigitalSimulator, asiSetEnvOptionChoices, asiSetEnvOptionVal, asiSetFormFieldChoices, asiSetFormFieldVal, asiSetKeepList, asiSetKeepOptionChoices, asiSetKeepOptionVal, asiSetMarchList, asiSetNetlistFormatterClass, asiSetNetlistOption, asiSetPlotList, asiSetSimOptionChoices, asiSetSimOptionVal, asiSetSyncFlag, asiSetValid, asiSetVerilogHost, asiSetVerilogHostMode, asiTransientStore, asiVerilogNetlistMoreCB, asiWriteOceanScript, asin, asin, assoc, atan, atan, atan2, atof, atoi, atom, auCdl, auHiUltraPCell, auLvs, auLvsGetLabelSuffix, auProbeAddDevsForNet, avSwitch, average, awvAddSubwindow, awvAnalog2Digital, awvAppendExpression, awvAppendList, awvAppendWaveform, awvAxesOptionMenuCB, awvClearPlotWindow, awvClearSubwindowHistory, awvClearWindowHistory, awvCloseWindowMenuCB, awvCreateBusFromWaveList, awvCreatePlotWindow, awvCrossHairDeleteList, awvCrossHairMenuACB, awvCrossHairMenuBCB, awvCursorMenuCB, awvCurvesMenuCB, awvDeleteAllWaveforms, awvDeleteMenuCB, awvDeleteSubwindow, awvDeleteWaveform, awvDigital2Analog, awvDisableRedraw, awvDisplayDate, awvDisplayGrid, awvDisplaySubwindowTitle, awvDisplayTitle, awvEraseWindowMenuCB, awvExitWindowFunctionAdd, awvExitWindowFunctionDel, awvExitWindowFunctionGet, awvExpandBusMenuCB, awvFastZoomInMenuCB, awvFastZoomOutMenuCB, awvFitMenuCB, awvGetCurrentSubwindow, awvGetCurrentWindow, awvGetDisplayMode, awvGetDrawStatus, awvGetInitializationTimeout, awvGetOnSubwindowList, awvGetPlotStyle, awvGetSavePromptNeeded, awvGetSmithModeType, awvGetSubwindowList, awvGetUnusedEntityList, awvGetWaveNameList, awvGetWaveformToolName, awvGetWindowList, awvGetXAxisLabel, awvGetXMarkerNames, awvGetYAxisLabel, awvGetYMarkerNames, awvHardCopyMenuCB, awvHorizontalMarkerMenuCB, awvInitWindowFunctionAdd, awvInitWindowFunctionDel, awvInitWindowFunctionGet, awvIsPlotWindow, awvLabelMenuCB, awvLoadMenuCB, awvLoadWindow, awvLogXAxis, awvLogYAxis, awvMakeBusMenuCB, awvMakeWindowActiveMenuCB, awvMoveMenuCB, awvPanMenuCB, awvPlaceWaveformLabel, awvPlaceWindowLabel, awvPlaceXMarker, awvPlaceXMarker, awvPlaceYMarker, awvPlaceYMarker, awvPlotExpression, awvPlotList, awvPlotStyleMenuCB, awvPlotWaveform, awvPlotWaveformOption, awvPrintWaveform, awvRedisplaySubwindow, awvRedisplayWindow, awvRedrawWindowMenuCB, awvRemoveDate, awvRemoveLabel, awvRemoveSubwindowTitle, awvRemoveTitle, awvResetAllWindows, awvResetWindow, awvResetWindowMenuCB, awvRotateStripsMenuCB, awvSaveMenuCB, awvSaveWindow, awvSetCurrentSubwindow, awvSetCurrentWindow, awvSetCursorPrompts, awvSetDisplayMode, awvSetDisplayStatus, awvSetInitializationTimeout, awvSetOptionDefault, awvSetOptionValue, awvSetOrigin, awvSetPlotStyle, awvSetSavePromptNeeded, awvSetSmithModeType, awvSetSmithXLimit, awvSetSmithYLimit, awvSetUpdateStatus, awvSetWaveformDisplayStatus, awvSetWaveformToolName, awvSetXAxisLabel, awvSetXLimit, awvSetXScale, awvSetYAxisLabel, awvSetYLimit, awvSetYRange, awvSimplePlotExpression, awvSmithAxisMenuCB, awvSubwindowMenuCB, awvSubwindowTitleMenuCB, awvSwapStripsMenuCB, awvTitleMenuCB, awvToCompositeMenuCB, awvToSmithAdmittanceMenuCB, awvToSmithImpedanceMenuCB, awvToSmithPolarMenuCB, awvToStripMenuCB, awvUndoMenuCB, awvUpdateAllWindows, awvUpdateWindow, awvUpdateWindowMenuCB, awvVerticalMarkerMenuCB, awvXAxisMenuCB, awvYAxesMenuCB, awvZoomInMenuCB, awvZoomOutMenuCB, awviEditMenuCB, awviMakeActiveMenuCB, awviPLoadMenuCB, awviPSaveMenuCB, awviPUpdateMenuCB, awviShowOutputMenuCB, b1f, band, bandwidth, bcdp, begin, begin, bitfield, bitfield1, blankstrp, bnand, bnor, bnot, booleanp, bor, bottomEdge, boundp, break, breakpt, buildString, bxnor, bxor, cPwrContour, cReflContour, caar, calCalcInput, calCalculatorFormCB, calCloseCalculatorCB, calCreateSpecialFunction, calCreateSpecialFunctionsForm, calGetBuffer, calRegisterSpecialFunction, calSpecialFunctionInput, caliMemoriesFileFormCB, caliMemoryNameFormCB, caliModeToggle, caliRestoreDefaultWindowSize, caliShowMemoriesCB, caliWaveObjectFormCB, callAs, callInitProc, callNextMethod, cancelEnterFun, car, case, case, cat, ccpCopy, ccpCopyConfig, ccpCopyDesign, ccpCopyExactDesign, ccpExpand, ccpExpandConfig, ccpExpandDesign, ccpExpandExactDesign, ccpRegMonitor, ccpRegTrigger, ccpRemoveTrigger, ccpRename, ccpRenameReferenceLib, cdcUICbk, cdfCopyCDF, cdfCopyParam, cdfCreateBaseCellCDF, cdfCreateBaseLibCDF, cdfCreateParam, cdfCreateUserCellCDF, cdfCreateUserLibCDF, cdfDeleteCDF, cdfDeleteParam, cdfDump, cdfDumpAll, cdfEditScaleFactors, cdfFindParamByName, cdfFormatFloatString, cdfGetBaseCellCDF, cdfGetBaseLibCDF, cdfGetCellCDF, cdfGetInstCDF, cdfGetLibCDF, cdfGetUnitScaleFactor, cdfGetUserCellCDF, cdfGetUserLibCDF, cdfParseFloatString, cdfRefreshCDF, cdfSaveCDF, cdfSetUnitScaleFactor, cdfSyncInstParamValue, cdfUpdateInstParam, cdfUpdateInstSingleParam, cdlinGuiDisplay, cdr, cdsGetInstPath, cdsGetNetlistMode, cdsMosDevice, cdsName, cdsNetExpr, cdsParam, cdsSetNetlistMode, cdsTerm, cdsViaDevice, ceiling, centerBox, changeEnterFun, changeNextEnterFun, changeWorkingDir, charToInt, characterizationRules, ciwAddTranslatorItem, ciwCreateMenu, ciwCreatePulldown, ciwHiExit, ciwMenuInit, ciwMenuInstall, ciwMenuLoadFile, ciwRemoveTranslatorItem, className, classOf, classp, clear, clearAll, clearAllEnterFunctions, clearExitProcs, clearSubwindow, clip, close, close, cmdCtrlOption, cmdOption, cmdShiftOption, cmxfAsc, cmxfAxisCreate, cmxfAxisDelete, cmxfAxisGetCV, cmxfAxisGetDirection, cmxfAxisGetId, cmxfAxisGetName, cmxfAxisGetX, cmxfAxisGetY, cmxfAxisIsActive, cmxfAxisIsConstrained, cmxfAxisIsHidden, cmxfAxisIsId, cmxfAxisIsInherited, cmxfAxisIsReadOnly, cmxfAxisIsValid, cmxfAxisNameGen, cmxfAxisSetActive, cmxfAxisSetDirection, cmxfAxisSetX, cmxfAxisSetY, cmxfCatAttrGetDefVal, cmxfCatAttrGetType, cmxfCatGetId, cmxfCatGetName, cmxfCatGetNumAttrs, cmxfCatGetWeight, cmxfCatIsAttr, cmxfCatIsId, cmxfClstrChngMems, cmxfClstrCreate, cmxfClstrDelete, cmxfClstrGetBBox, cmxfClstrGetCV, cmxfClstrGetCenter, cmxfClstrGetId, cmxfClstrGetName, cmxfClstrGetSize, cmxfClstrHasMems, cmxfClstrIsActive, cmxfClstrIsConstrained, cmxfClstrIsHidden, cmxfClstrIsId, cmxfClstrIsInherited, cmxfClstrIsMember, cmxfClstrIsReadOnly, cmxfClstrIsValid, cmxfClstrNameGen, cmxfClstrSetActive, cmxfConChngMems, cmxfConCreate, cmxfConDelete, cmxfConGetAttrVal, cmxfConGetCV, cmxfConGetCatId, cmxfConGetId, cmxfConGetName, cmxfConGetRefFig, cmxfConGetWeight, cmxfConHasMems, cmxfConIsActive, cmxfConIsHidden, cmxfConIsId, cmxfConIsInherited, cmxfConIsOverCon, cmxfConIsReadOnly, cmxfConIsSatisfied, cmxfConIsValid, cmxfConModify, cmxfConNameGen, cmxfConSetActive, cmxfConSetAttrVal, cmxfConSetWeight, cmxfConVerify, cmxfFigIsConstrained, cmxfGenAxis, cmxfGenAxisToCon, cmxfGenCatToCon, cmxfGenClstr, cmxfGenClstrToCon, cmxfGenClstrToFig, cmxfGenClstrToMemName, cmxfGenCon, cmxfGenConToFig, cmxfGenConToMemName, cmxfGenConToNet, cmxfGenFigToCon, cmxfGenIsId, cmxfGenNC, cmxfGenNCToCon, cmxfGenNCToMemName, cmxfGenNCToNet, cmxfGenNetToCon, cmxfGenSuperCat, cmxfGenSuperCatToCat, cmxfGuiReportGen, cmxfGuiStartUp, cmxfNCChngMems, cmxfNCCreate, cmxfNCDelete, cmxfNCGetCV, cmxfNCGetId, cmxfNCGetName, cmxfNCGetSize, cmxfNCHasMem, cmxfNCIsActive, cmxfNCIsConstrained, cmxfNCIsHidden, cmxfNCIsId, cmxfNCIsInherited, cmxfNCIsMember, cmxfNCIsReadOnly, cmxfNCIsValid, cmxfNCNameGen, cmxfNCSetActive, cmxfNetIsConstrained, cmxfSCatGetId, cmxfSCatGetName, cmxfSCatIsId, cmxfShapeNameGen, cmxfStartGenAxis, cmxfStartGenAxisToCon, cmxfStartGenCatToCon, cmxfStartGenClstr, cmxfStartGenClstrToCon, cmxfStartGenClstrToFig, cmxfStartGenClstrToMemName, cmxfStartGenCon, cmxfStartGenConToFig, cmxfStartGenConToMemName, cmxfStartGenConToNet, cmxfStartGenFigToCon, cmxfStartGenNC, cmxfStartGenNCToCon, cmxfStartGenNCToMemName, cmxfStartGenNCToNet, cmxfStartGenNetToCon, cmxfStartGenSuperCat, cmxfStartGenSuperCatToCat, cmxfStopGen, colorIndex, comment, compactorLayers, compactorRules, compareTime, complex, complexp, compress, compression, compressionVRI, compressionVRICurves, concat, cond, cond, conjugate, conn2Sch, conn2SchImpHdlDisplay, conn2SchStartUp, connectRules, cons, constar, cont, controls, converge, convolve, copy, copyDefstructDeep, copy_, corAddCorner, corAddDesignVar, corAddGroupAndVariantChoices, corAddMeas, corAddModelFileAndSectionChoices, corAddProcess, corAddProcessVar, corCopyCorner, corDeleteAllDesignVars, corDeleteAllMeas, corDeleteAllProcesses, corGetModelStyleChoices, corSetCornerGroupVariant, corSetCornerModelFileSection, corSetCornerNomTempVal, corSetCornerRunTempVal, corSetCornerVarVal, corSetDesignVarVal, corSetMeasExpression, corSetMeasGraphicalOn, corSetMeasLower, corSetMeasTarget, corSetMeasTextualOn, corSetMeasUpper, corSetModelFile, corSetProcessVarVal, cornerDesVar, cornerMeas, cornerRun, cornerRunTemp, correlationTable, cos, cos, count, cputime, createDir, createFinalNetlist, createNetlist, cross, csh, currentSubwindow, currentWindow, customViaDefs, dagAddActionToObject, dagCreateClass, dagCreateNode, dagDeleteActionFromObject, dagDestroyNode, dagDisplayTool, dagGetCurrentObject, dagGetCurrentTool, dagLinkParentToChild, dagNumToTool, dagOpenTool, dagPopTool, dagPushTool, dagRefreshObject, dagSetActionStatus, dagSetCurrentTool, dagSetExpandedActionStatus, dagSetFont, dagUnlinkParentFromChild, dataFilter, dataTrigger, dataTypes, db10, db20, dbAddFigToNet, dbAddObjectToGroup, dbAddObjectsToMarker, dbAllCellViews, dbCVHasInheritedSig, dbCVHasInheritedTerm, dbCVHasNetSetProp, dbChangeGroupType, dbCheck, dbCheckRecursion, dbClose, dbCloseBag, dbCompactCellView, dbCompressPointArray, dbCompressionPlot, dbComputeBBox, dbComputeBBoxNoNLP, dbConcatTransform, dbConnectObject, dbConvertArcToLine, dbConvertDonutToPolygon, dbConvertEllipseToPolygon, dbConvertLineToPath, dbConvertPathToPolygon, dbCopyCellView, dbCopyFig, dbCopyProp, dbCopySingleProp, dbCreateArc, dbCreateBlockage, dbCreateCluster, dbCreateClusterBoundary, dbCreateConn, dbCreateConnByName, dbCreateConnByNewName, dbCreateDonut, dbCreateDot, dbCreateEllipse, dbCreateEnumProp, dbCreateGroup, dbCreateHierProp, dbCreateInst, dbCreateInstByMasterName, dbCreateInstPin, dbCreateInstPinByName, dbCreateInstTerm, dbCreateLabel, dbCreateLib, dbCreateLine, dbCreateMarker, dbCreateNet, dbCreatePRBoundary, dbCreateParamInst, dbCreateParamInstByMasterName, dbCreateParamSimpleMosaic, dbCreateParamSimpleMosaicByMasterName, dbCreatePath, dbCreatePin, dbCreatePolygon, dbCreateProp, dbCreateRangeProp, dbCreateRect, dbCreateRow, dbCreateSigNetExpr, dbCreateSimpleMosaic, dbCreateSimpleMosaicByMasterName, dbCreateSnapBoundary, dbCreateStrongGroup, dbCreateTerm, dbCreateTermNetExpr, dbCreateTextDisplay, dbCreateTrackPattern, dbCreateUniqueNamedNet, dbCreateViaShapeArray, dbCreateXformPCell, dbDebugStats2, dbDefineProc, dbDeleteAllNet, dbDeleteGroupByName, dbDeleteGroupMemberFromGroup, dbDeleteObject, dbDeleteObjectFromGroup, dbDeletePropByName, dbDeleteSigNetExpr, dbDeleteTermNetExpr, dbDisableCellViewCounter, dbDisassembleRoute, dbDisconnectObject, dbDumpBag, dbDumpBagWithLib, dbDumpPcDefinePcell, dbDumpPcell, dbDumpPcellWithLib, dbEnableCellViewCounter, dbEvalParamCell, dbExpandToMemNameExpr, dbExternallyConnectPins, dbFindAnyInstByName, dbFindClusterBoundaryByName, dbFindClusterByName, dbFindMemInstByName, dbFindMosaicByName, dbFindNetByName, dbFindOpenBag, dbFindOpenCellView, dbFindProp, dbFindRouteSpecHeaderByName, dbFindRowByName, dbFindRowHeader, dbFindSigByAlias, dbFindSigByName, dbFindTermByName, dbFindViaHeaderByName, dbFix, dbFlattenInst, dbFullLibPath, dbFullPath, dbGetAnyInstSwitchMaster, dbGetAssociateTextDisplay, dbGetBag, dbGetBagDdId, dbGetBagOwner, dbGetBagTimeStamp, dbGetBuildInfo, dbGetCellViewCounter, dbGetCellViewDdId, dbGetCellViewInheritedSig, dbGetCellViewInheritedTerm, dbGetCellViewSiteDef, dbGetCellViewSiteDefName, dbGetDatabaseType, dbGetFileDatabaseType, dbGetFileTimeStamp, dbGetInstByName, dbGetInstNameNumInst, dbGetInstanceByName, dbGetMaxHierDepth, dbGetMaxNumBit, dbGetMemName, dbGetMemNetSig, dbGetMemNetSigName, dbGetNameNumBit, dbGetNeighbor, dbGetNeighborList, dbGetNetTerms, dbGetOpenBags, dbGetOpenCellViews, dbGetOverlaps, dbGetPinConnection, dbGetPropByName, dbGetSigNameMemNets, dbGetSigNetExpr, dbGetStrongGroups, dbGetTermGroundSensitivity, dbGetTermNetExpr, dbGetTermSupplySensitivity, dbGetTrueOverlaps, dbGetVersion, dbGetq, dbHasExternalConnections, dbHasMustConnections, dbHasShortedTerminals, dbIsAnyInstBaseNameUsed, dbIsBagId, dbIsBagReadable, dbIsBagWritable, dbIsCellViewModified, dbIsCellViewStdViaMaster, dbIsConnCurrent, dbIsId, dbIsNetNameDescend, dbIsObjectBound, dbIsObjectConnected, dbIsSameMaster, dbIsSigInherited, dbIsSigNameGlobal, dbIsTermInherited, dbIsValidSigNetExpr, dbIsValidTermNetExpr, dbLayerAnd, dbLayerAndNot, dbLayerEnclose, dbLayerHoles, dbLayerInside, dbLayerNoHoles, dbLayerOr, dbLayerOutside, dbLayerSize, dbLayerStraddle, dbLayerTile, dbLayerXor, dbMakeBagCurrent, dbMakeNet, dbMergeNet, dbMergeSingleProp, dbMergeStrongGroups, dbMoveFig, dbMoveGroupMember, dbMustConnectPins, dbNetIsSig, dbObjectHasConnection, dbOpenAutoSavedCellView, dbOpenBag, dbOpenBag, dbOpenCellViewByType, dbOpenHier, dbOpenPanicCellView, dbParseSigNetExpr, dbParseTermNetExpr, dbProduceMemName, dbProduceOverlap, dbProduceOverlapInst, dbPurge, dbPurgeBag, dbReOpenBag, dbRefreshBag, dbRefreshCellView, dbRegPostSaveTrigger, dbRegPurgeTrigger, dbRemoveObjectsFromMarker, dbReopen, dbReplaceEnumProp, dbReplaceHierProp, dbReplaceParam, dbReplaceProp, dbReplacePropList, dbReplaceRangeProp, dbReplaceSigNetExpr, dbReplaceTermNetExpr, dbSave, dbSaveBag, dbSearchCDF, dbSearchPropByName, dbSetAutoSave, dbSetBoundaryEdge, dbSetCellViewSiteDef, dbSetConnCurrent, dbSetInstHeaderMasterName, dbSetInstHeadersCounter, dbSetNIMsgMode, dbSetNetNameDescend, dbSetTermGroundSensitivity, dbSetTermSupplySensitivity, dbSetTextDisplayNameValueVisible, dbSetq, dbStronglyConnectPins, dbSubFigFromNet, dbTransformBBox, dbTransformCellView, dbTransformPoint, dbUndefineProc, dbUnregPostSaveTrigger, dbUnregPurgeTrigger, dbValidP, dbWeaklyConnectPins, dbWriteSkill, dbWriteSkillWithLib, dbm, dbobjectp, dc, dcmatchSummary, ddAmUsingDM, ddAutoCheckin, ddAutoCheckout, ddAutoCtlGetVars, ddAutoCtlSetVars, ddCatAddItem, ddCatClose, ddCatFindCat, ddCatGetCatMembers, ddCatGetCatName, ddCatGetCellsNotInCat, ddCatGetLibCats, ddCatIsObjInCat, ddCatOpen, ddCatOpenEx, ddCatRemove, ddCatSave, ddCatSubItem, ddCheckin, ddCheckout, ddClearLibUnmanaged, ddCreateLib, ddDeleteLocal, ddDeleteObj, ddDumpLibList, ddGetForcedLib, ddGetLibList, ddGetObj, ddGetObjAccess, ddGetObjAccessEx, ddGetObjChildren, ddGetObjChildrenEx, ddGetObjDMSys, ddGetObjFiles, ddGetObjFilesEx, ddGetObjGroup, ddGetObjGroupEx, ddGetObjLastModify, ddGetObjLastModifyEx, ddGetObjLib, ddGetObjName, ddGetObjOwner, ddGetObjOwnerEx, ddGetObjParent, ddGetObjReadPath, ddGetObjReadPathEx, ddGetObjSize, ddGetObjSizeEx, ddGetObjType, ddGetObjWritePath, ddGetObjWritePathEx, ddGetPathDMSys, ddGetStartup, ddGetTmpdirWarn, ddGetTriggerList, ddGetUpdatedLib, ddIsId, ddIsObjLibMarked, ddIsObjReadable, ddIsObjReadableEx, ddIsObjWritable, ddIsObjWritableEx, ddIsPathManageable, ddIsPathManaged, ddLibHasTmpDir, ddLockDescribe, ddLockFree, ddLockGetStatus, ddLockPath, ddLockSet, ddMapGetDataTypeFileName, ddMapGetDataTypeList, ddMapGetFileDataType, ddMapGetFileTool, ddMapGetFileViewType, ddMapGetToolNameDataType, ddMapGetToolNameList, ddMapGetToolNameViewType, ddMapGetViewTypeFileName, ddMapGetViewTypeList, ddMapIsFileNameKnown, ddMapIsToolLegal, ddMapIsToolNameKnown, ddMarkObjMaster, ddNeedCheckout, ddObjHasNamedObj, ddObjHasNamedObjEx, ddObjIsInMaster, ddObjIsInTmp, ddProduceObjRelNames, ddProduceObjRelNamesEx, ddRegTrigger, ddReleaseObj, ddSetForcedLib, ddSetLibUnmanaged, ddSetTmpdirWarn, ddUnregTrigger, ddUpdateLibList, ddsAboutDFII, ddsBrowserAction, ddsCheckSessionName, ddsCheckinPreferences, ddsCheckoutPreferences, ddsCloseLibManager, ddsCvtAMSTranslateCell, ddsCvtAMSTranslateLib, ddsCvtToolBoxAMS, ddsEditProp, ddsEndSyncWithForm, ddsGetCurrentLib, ddsGetCustomLibSelect, ddsGetServiceName, ddsHiCloseData, ddsHiCreateLibrary, ddsHiDefragCellView, ddsHiDefragLib, ddsHiMakeReadOnly, ddsHiOpenCdsLibEditor, ddsHiPreferences, ddsHiRefresh, ddsLEWhatsNewAtStartup, ddsOnLineHelp, ddsOpenConvertToolBox, ddsOpenLibManager, ddsProgramName, ddsRegCustomLibSelect, ddsServSelect, ddsSyncWithForm, ddsUnregCustomLibSelect, ddsUpdateSyncWithForm, ddsViewDDProp, ddsWhatsNew, deCopy, deDisableApp, deEditInPlace, deEnableApp, deGetAllRootTypes, deGetAllViewTypes, deGetAppInfo, deGetCellView, deGetConfigId, deGetContext, deGetEIP, deGetEditCellView, deGetEditContext, deGetEditLevel, deGetEditRootType, deGetEditViewType, deGetMaxEditLevel, deGetMaxInContextLevel, deGetProbeContext, deGetProbeTrigger, deGetPushed, deGetRootType, deGetVector, deGetViewType, deHiSetDefaultApplication, deInheritEnabling, deInheritMenus, deInstallApp, deIsEditInContext, deNew, deNewCellView, deOpen, deOpenCellView, dePush, deRegApp, deRegProbeTrigger, deRegUserTriggers, deReturn, deSetCellView, deSetConfigId, deSetContext, deSetProbeContext, deSwap, deUnInheritEnabling, deUnInheritMenus, deUnRegApp, deUnRegProbeTrigger, deUnRegUserTriggers, deViewTypeIsRegistered, debugQuit, debugStatus, declare, declareLambda, declareNLambda, declareSQNLambda, defCapDepends, defCapPrefixes, defInitProc, defMathConstants, defUserInitProc, defclass, defgeneric, define, definitionFile, defmacro, defmethod, defprop, defstruct, defstructp, defun, defvar, delay, delete, deleteChecks, deleteDir, deleteFile, deleteJob, deletePoint, deleteSubwindow, deleteWaveform, deriv, desVar, design, devices, dft, dftbb, difference, digitalHostMode, digitalHostName, disableAllChecks, disableChecks, disableDeviceChecking, discipline, display, displayChecks, displayMode, displayNetlist, dlAddArc, dlAddBox, dlAddCircle, dlAddDonut, dlAddEventObject, dlAddPath, dlAddPoint, dlAddPolygon, dlAddRasterText, dlAddSegment, dlAddSkillObject, dlAddStrokeText, dlAttachDlistToWidget, dlAttachDlistToWindow, dlClearDisplayList, dlCloseWidget, dlConfigureButton, dlDetachDlistFromWidget, dlDetachDlistFromWindow, dlDisplay, dlDlistToIcon, dlDrawArc, dlDrawBox, dlDrawCircle, dlDrawDonut, dlDrawPath, dlDrawPoint, dlDrawPolygon, dlDrawRasterText, dlDrawSegment, dlDrawStrokeText, dlEnableItem, dlFitDlistOnDraw, dlMakeDisplayList, dlMakeDlistButton, dlMakePenTable, dlMakeStipple, dlMakeWidget, dlMapWidget, dlMoveButton, dlResizeButton, dlSaveDlist, dlSetClearOnDraw, dlSetCurrentPen, dlSetDlistPosition, dlSetDlistScale, dlSetPenColor, dlSetPenFillStyle, dlSetPenFilled, dlSetPenStipple, dlSetPenTable, dlSetWidgetName, dlUnMapWidget, dmAddCellToCategory, dmCreateCellCategory, dmRemoveCellCategory, dmRemoveCellFromCategory, dmbGetCurrentLib, dmbLibBrowserFormDone, dmbOpenLibBrowser, dmbSetBrowserInteractions, dmbSyncWithCellForm, dmbSyncWithForm, dmbUpdateBrowser, do, dplp, drAddElem, drCreateEmptyWaveform, drCreateVec, drCreateWaveform, drDefineColor, drDefineDisplay, drDefineLineStyle, drDefinePacket, drDefinePacketAlias, drDefineStipple, drDeleteColor, drDeleteDisplay, drDeleteLineStyle, drDeletePacket, drDeleteStipple, drDumpDrf, drFindPacket, drGetColor, drGetDisplay, drGetDisplayIdList, drGetDisplayName, drGetDisplayNameList, drGetElem, drGetLineStyle, drGetLineStyleIndexByName, drGetPacket, drGetPacketAlias, drGetPacketFillStyle, drGetPacketList, drGetStipple, drGetStippleIndexByName, drGetWaveformXType, drGetWaveformXVec, drGetWaveformYType, drGetWaveformYVec, drIsDataVector, drIsParamWave, drIsWaveform, drLoadDrf, drPutWaveformXVec, drPutWaveformYVec, drSetElem, drSetPacket, drType, drVectorLength, drain, drawEnterFun, drbBrowseFormCB, dreInvokeDre, dtpr, dump, ecoCmpCB, ecoCompare, ed, edi, ediFinishStatus, edifDisplay, edifLayerNumMap, edifinDisplay, edifinMakeRenameString, edifoutAddCellInfo, edifoutAddInstInfo, edifoutAddInterfaceInfo, edifoutAddLibraryInfo, edifoutAddNetInfo, edifoutAddPortInfo, edifoutAddViewInfo, edifoutEditCellProperty, edifoutEditInstProperty, edifoutEditLibProperty, edifoutEditNetProperty, edifoutEditPortProperty, edifoutEditProperty, edifoutEditViewProperty, edifoutMakeRenameString, edit, edl, electricalRules, enableAllChecks, enableChecks, enableDeviceChecking, enableTrigger, encrypt, enterArc, enterBox, enterCircle, enterDonut, enterEllipse, enterLine, enterMultiRep, enterNumber, enterPath, enterPoint, enterPoints, enterPolygon, enterScreenBox, enterSegment, enterString, envCyclicStringToIndex, envGetAvailableTools, envGetDefVal, envGetLoadedTools, envGetModifiedTools, envGetVal, envGetVarType, envIsToolModified, envIsVal, envLoadFile, envLoadVals, envOption, envRegLoadDumpTrigger, envRegSetTrigger, envSetToolCurrValToDefault, envSetToolDefaultToCurrVal, envSetVal, envSetVarCurrValToDefault, envSetVarDefaultToCurrVal, envStoreEnv, envobj, eq, equal, equivalentLayers, eqv, err, error, errset, errsetstring, eval, evalstring, evcdFile, evcdInfoFile, evenp, evmQpsk, exists, exit, exp, exp, expandMacro, expr, expt, eyeDiagram, famAddValue, famCreateFamily, famEval, famGetSweepName, famGetSweepValues, famIsFamily, famMap, famValue, fboundp, fileLength, fileSeek, fileTell, fileTimeModified, findClass, finishEnterFun, fix, fixp, flip, float, floatp, floor, fnlAbortNetlist, fnlCurrentCell, fnlCurrentCellCdsName, fnlCurrentInst, fnlCurrentInstCdsName, fnlCurrentIteration, fnlCurrentModelExtName, fnlCurrentSig, fnlCurrentSigPathName, fnlGetGlobalSigNames, fnlInstCdsNameExtName, fnlPathList, fnlPrint, fnlSearchPropString, fnlSigCdsNameExtName, fnlTermCdsNameExtName, fnlTermExtName, fnlTopCell, for, for, forall, forcenode, foreach, foreach, fourEval, fprintf, freq, frequency, fscanf, fscanf, funcall, funobj, ga, gac, gainBwProd, gainMargin, gc, gciDisplayMarkers, gciForceNetSpecial, gciWriteCellViewLEFDEF, gcsummary, gdmAddSpecToSpecList, gdmCreateSpec, gdmCreateSpecFromDDID, gdmCreateSpecList, gdmInspectSpec, gdmIsSpecId, gdmNextFromSpecList, gdmResetSpecList, gdmSpecListp, gdmSpecType, gdmSpecp, gdmcancel, gdmci, gdmco, gdmdelete, gdmexport, gdmhistory, gdmsetdefver, gdmsetname, gdmstatus, gdmsubmit, gdmupdate, geAddHilightArc, geAddHilightCircle, geAddHilightDonut, geAddHilightDot, geAddHilightEllipse, geAddHilightFig, geAddHilightLabel, geAddHilightLine, geAddHilightPath, geAddHilightPolygon, geAddHilightRectangle, geAddInstProbe, geAddNetProbe, geAddPathProbe, geAddSelectBox, geAddSelectPoint, geAddTermProbe, geAdjustPoint, geAdvWireSelection, geAskOverwriteMode, geCellViewToDlist, geCenter, geChangeCellView, geChangeEditMode, geClosePolygon, geCompressPointArray, geComputeQueryBbox, geCoordp, geCopySelSet, geCreateBend, geCreateHilightSet, geCreateMarker, geCreateMarkerByBBox, geCreateMarkerByPoints, geCreateTaper, geCreateTrl, geCycleSelectNext, geDeleteAllAreaViewLevel, geDeleteAllHilightSet, geDeleteAllInstViewLevel, geDeleteAllMarker, geDeleteAllProbe, geDeleteAreaViewLevel, geDeleteHilightObject, geDeleteHilightSet, geDeleteInstProbe, geDeleteInstViewLevel, geDeleteNetProbe, geDeletePathProbe, geDeleteProbe, geDeleteSelSet, geDeleteTermProbe, geDeselectAll, geDeselectAllFig, geDeselectAllObject, geDeselectArea, geDeselectFig, geDeselectFigNoFilter, geDeselectFigPoint, geDeselectObject, geDeselectObjectNoFilter, geDeselectObjectPoint, geDeselectPoint, geDiscardEdits, geDrawHilightSet, geEditInPlace, geEditSelSet, geEditSingleObject, geEditToWindowPoint, geEquivWindow, geEquivWindowSet, geExplainProbe, geFindPathToCellViewInst, geGetAdjustedPath, geGetAllAreaViewLevel, geGetAllInstViewLevel, geGetAllProbe, geGetAreaViewLevel, geGetCellViewWindow, geGetCrossProbeValue, geGetCurrentHilightSet, geGetDisabledMenus, geGetEditCellView, geGetEditCellViewWindow, geGetEditRep, geGetEnablePredicate, geGetFigHilightObjectId, geGetHierMemInst, geGetInstHier, geGetInstHierPath, geGetInstTransform, geGetInstViewLevel, geGetMarkerDisplayText, geGetMarkerOwner, geGetMarkerReason, geGetMarkerType, geGetObjectSelectedSet, geGetPotentialSelection, geGetProbeNetStopLevel, geGetProbeUserLPP, geGetScrollPercent, geGetSelSet, geGetSelSetCount, geGetSelSetFigPoint, geGetSelectedSet, geGetTopCellViewWindow, geGetTopLevelCellView, geGetWindowBox, geGetWindowCellView, geGetWindowRep, geGetWindowStatus, geHasMarker, geHiCommonFindMarker, geHiDeleteAllMarker, geHiDeleteMarker, geHiDisplaySelectionForm, geHiDragFig, geHiExplainMarker, geHiFindMarker, geInstIdToFileSpec, geIsFigAllSelected, geIsFigPartiallySelected, geIsFigSelected, geIsMarkerShape, geIsObjectAllSelected, geIsObjectPartiallySelected, geIsObjectSelected, geIsPathReentrant, geIsPolygonReentrant, geIsValidHilightObject, geIsValidHilightSet, geLoadIcon, geLoadNextProbe, geLoadProbe, geMakeProbeWithColor, geMakeProbeWithLabel, geMakeProbeWithOnColor, geNewWindow, geOpen, gePointQuery, gePointToFig, gePopHilightStack, gePrintAllProbe, geProbeP, gePush, gePushHilightStack, geQuerySelSet, geRefresh, geRefreshCellView, geRefreshCellViews, geRefreshWindow, geRefreshWorldView, geRegCrossProbeTrigger, geResetCellView, geRestoreWinView, geReturn, geReturnToLevel, geSave, geSaveAs, geSaveIcon, geSaveProbe, geSaveWinView, geScroll, geSelectAll, geSelectAllFig, geSelectArea, geSelectFig, geSelectFigNoFilter, geSelectFigPoint, geSelectFigs, geSelectObject, geSelectObjectNoFilter, geSelectPoint, geSetAreaViewLevel, geSetCrossProbeValue, geSetEnablePredicate, geSetEnterFunctionSelectionMode, geSetInstViewLevel, geSetProbeNetStopLevel, geSetProbeUserLPP, geSetScrollPercent, geSetWindowStatus, geShift, geSingleSelectBox, geSingleSelectPoint, geSubSelectBox, geSubSelectPoint, geSwitch, geToggleAreaSelectOption, geToggleFilterSize, geTogglePartialSelect, geTransformUserBBox, geTransformUserPoint, geUnRegCrossProbeTrigger, geUpdateWindowMenus, geValidateWindow, geWindowToEditPoint, geWriteCellView, geZoomToGrid, geiLoadProbe, geiSaveProbe, gensym, geoBoxCenterPt, geoPtAbovep, geoPtPtDistance, geoPtRightp, geqp, get, getAllLoadedFiles, getAsciiWave, getCallingFunction, getCurrentTime, getCurrentWindow, getData, getDirFiles, getFnWriteProtect, getFunType, getFunctions, getInstallPath, getLogin, getMaxScreenCoords, getMethodSpecializers, getPrompts, getResult, getShellEnvVar, getSkillPath, getSkillVersion, getTempDir, getVarWriteProtect, getVersion, getWarn, getWorkingDir, get_filename, get_pname, get_string, getc, getchar, getd, getq, getqq, gets, gets, globalSigAlias, globalSignal, gmax, gmin, gmsg, gmux, go, gp, gpc, graphicsOff, graphicsOn, greaterp, groupDelay, gt, hardCopy, hardCopyOptions, harmonic, harmonicFreqList, harmonicList, hdbAddConfigToBag, hdbBind, hdbClose, hdbCloseConfigsInBag, hdbCreateConfigBag, hdbCreatePathVector, hdbCreatePathVectorStack, hdbDestroyPathVector, hdbExplainBinding, hdbExplainSignature, hdbExplainStop, hdbGetCellName, hdbGetLibName, hdbGetPathStackStr, hdbGetPathStr, hdbGetPathVectorStackTop, hdbGetTopCellName, hdbGetTopLibName, hdbGetTopViewName, hdbGetViewName, hdbIsAtStopPoint, hdbIsConfig, hdbIsOpenConfig, hdbOpen, hdbPopCell, hdbPopPathVect, hdbPushCell, hdbPushPathVect, hdbSaveAs, hdpCreatePinMC, hdpMkPartitionForm, heHiEditConfig, heHiSetInstBinding, heHiShowViewsFound, heHiUpdate, help, hiAbout, hiAbsolutePan, hiActivateBatch, hiActivateMessages, hiAddCyclicChoice, hiAddField, hiAddFields, hiAddFixedMenu, hiAddMenuItem, hiAddTextWordDelimiter, hiAppendInputCmd, hiBatchProcess, hiBeginProcess, hiBoxCenter, hiCancelProgressBox, hiChangeBannerLabel, hiChangeFormCallback, hiChangeFormTitle, hiCheckAbort, hiCloseChild, hiCloseWindow, hiCloseWorldView, hiCollapseTreeItem, hiContChild, hiCreate2DMenu, hiCreateAppForm, hiCreateBBoxField, hiCreateBooleanButton, hiCreateButton, hiCreateButtonBoxField, hiCreateColorArray, hiCreateComboField, hiCreateCyclicField, hiCreateFloatField, hiCreateFloatSpinBox, hiCreateForm, hiCreateFormButton, hiCreateFormLabel, hiCreateFrameField, hiCreateHorizontalFixedMenu, hiCreateHypertextField, hiCreateIntField, hiCreateLabel, hiCreateLayerCyclicField, hiCreateListBoxField, hiCreateListField, hiCreateMLTextField, hiCreateMenu, hiCreateMenuItem, hiCreateOptionsForm, hiCreateOutputStringField, hiCreatePointField, hiCreatePointListField, hiCreatePulldownMenu, hiCreateRadioField, hiCreateReportField, hiCreateScaleField, hiCreateScrollRegion, hiCreateSeparatorField, hiCreateSeparatorMenuItem, hiCreateSimpleHypertextField, hiCreateSimpleMenu, hiCreateSliderMenuItem, hiCreateSpinBox, hiCreateStringField, hiCreateTabField, hiCreateToggleField, hiCreateTree, hiCreateTreeItem, hiCreateTreeTable, hiCreateVerticalFixedMenu, hiCreateWindow, hiDBoxCancel, hiDBoxOK, hiDeiconifyWindow, hiDeleteBannerLabel, hiDeleteBannerMenu, hiDeleteBannerMenus, hiDeleteField, hiDeleteFields, hiDeleteForm, hiDeleteMenu, hiDeleteMenuItem, hiDeleteStatusBanner, hiDeltaPan, hiDisableMenuItem, hiDisableTailViewfile, hiDisplayAppDBox, hiDisplayBlockingDBox, hiDisplayEdgeMenu, hiDisplayFixedMenu, hiDisplayForm, hiDisplayListBox, hiDisplayMenu, hiDisplayModalDBox, hiDisplayModelessDBox, hiDisplayNonBlockingDBox, hiDisplayProgressBox, hiDisplaySysModalDBox, hiDisplayUserDBox, hiDisplayWindow, hiDisplayWindowMenu, hiEdgeFixedMenuDone, hiEditPropList, hiEditfile, hiEnableBlink, hiEnableMenuItem, hiEnableTailViewfile, hiEncap, hiEndLog, hiEscapeStringChars, hiExpandTreeItem, hiFixedMenuDown, hiFlush, hiFocusToCIW, hiFocusToCursor, hiFocusToEncap, hiFormApply, hiFormCancel, hiFormClose, hiFormDefaults, hiFormDone, hiFormFinish, hiFormList, hiFormUnmap, hiGenTextIndex, hiGetAbsWindowScreenBBox, hiGetAppType, hiGetAttention, hiGetBBoxResource, hiGetBannerMenus, hiGetBeepVolume, hiGetBindKey, hiGetBindKeyInheritAlias, hiGetBindKeyInheritRoot, hiGetBindKeyPrefixList, hiGetBindKeys, hiGetCIWindow, hiGetCommandPoint, hiGetCurrentCmd, hiGetCurrentField, hiGetCurrentForm, hiGetCurrentIndex, hiGetCurrentWindow, hiGetCursor, hiGetDBoxDefaultLocation, hiGetDbuPoint, hiGetDisplayName, hiGetDrawThruDelta, hiGetEncapSkillCmd, hiGetExitStatus, hiGetFieldInfo, hiGetFieldOverlaps, hiGetFont, hiGetGeometryResource, hiGetHelp, hiGetIconName, hiGetInsertionPosition, hiGetLayerCyclicValue, hiGetListBoxFieldFit, hiGetListBoxValue, hiGetLogFileName, hiGetMaxScreenCoords, hiGetMenuItems, hiGetMultiClickTime, hiGetNumMenus, hiGetNumVisibleItems, hiGetPid, hiGetPoint, hiGetPriority, hiGetProgress, hiGetProgressTotalSteps, hiGetScreenPoint, hiGetScreenSize, hiGetScrollBarInfo, hiGetStringResource, hiGetStrokeBBox, hiGetStrokeFirstPt, hiGetStrokeLastPt, hiGetTextCharAtLoc, hiGetTextClass, hiGetTextDispLoc, hiGetTextFieldFit, hiGetTextIndexLoc, hiGetTextLineColumn, hiGetTextSelByLoc, hiGetTextSelection, hiGetTextSourceLength, hiGetTextWidth, hiGetTextWordDelimiter, hiGetTopListItem, hiGetTree, hiGetTreeItemDescription, hiGetTreeItemParent, hiGetTreeItems, hiGetTreeParent, hiGetUndoLimit, hiGetUserAbort, hiGetViewBBox, hiGetWMOffsets, hiGetWidgetType, hiGetWindowFixedMenu, hiGetWindowIconifyState, hiGetWindowList, hiGetWindowMenu, hiGetWindowName, hiGetWindowState, hiGraphicMode, hiHelp, hiHighlightField, hiHypertextReadin, hiHypertextReplace, hiHypertextScroll, hiIconifyWindow, hiIgnoreProp, hiInEnterFun, hiInFormApply, hiInheritBindKey, hiInsertBannerMenu, hiInsertMenuItem, hiInstantiateForm, hiIs2DMenu, hiIsActiveChild, hiIsAliveChild, hiIsBlinkEnabled, hiIsForm, hiIsFormDisplayed, hiIsIcon, hiIsInFieldCancel, hiIsInReplay, hiIsInstantiated, hiIsMenu, hiIsMenuItemEnabled, hiIsMenuSlotFilled, hiIsProgressBoxCancelled, hiIsScrollRegion, hiIsWidgetType, hiIsWindowSpecifier, hiItemInsertTree, hiItemRemoveTree, hiKillAllProcs, hiKillChild, hiLaunchBrowser, hiLayerDispMainForm, hiLayerMatchCyclicStr, hiLayerStringToLPP, hiListBoxCancel, hiListBoxDone, hiListView, hiLogDragEvents, hiLookupHRFtag, hiLowerWindow, hiMakeLPChoiceList, hiMapWindow, hiMarkNestable, hiMarkNonNestable, hiMatchColor, hiMatchColorByName, hiModalDBoxHelp, hiModelessDBoxHelp, hiMoveField, hiMoveFixedMenu, hiMoveInsBarToEnd, hiMoveWindow, hiNextWinView, hiOffsetField, hiOffsetFields, hiOpenWindow, hiOpenWorldView, hiPan, hiPickWindow, hiPrevWinView, hiQueryFont, hiQuit, hiRaiseWindow, hiReadChild, hiReattachField, hiRedo, hiRedraw, hiRefreshTextWindow, hiRegCloseProc, hiRegSetBindKeyProc, hiRegTimer, hiRegZoomPanProc, hiRegisterBindKeyPrefix, hiRemoveFixedMenu, hiRemovePromptLine, hiRemoveTextWordDelimiter, hiRepeat, hiReplaceAllBannerMenus, hiReplaceTextWordDelimiter, hiReplayFile, hiReportDeselectAllItems, hiReportDeselectItem, hiReportDeselectItems, hiReportGetSelectedItems, hiReportSelectAllItems, hiReportSelectItem, hiReportSelectItems, hiResetAbort, hiResetProgressBox, hiResizeField, hiResizeWindow, hiRestoreView, hiSaveAsViewfile, hiSaveForRestore, hiSaveView, hiSaveViewfile, hiScaleBox, hiScrollWindowBottom, hiScrollWindowDown, hiScrollWindowLeft, hiScrollWindowRight, hiScrollWindowToCurrentIndex, hiScrollWindowToIndex, hiScrollWindowToName, hiScrollWindowTop, hiScrollWindowUp, hiSelectTextByLoc, hiSetAbort, hiSetBeepVolume, hiSetBindKey, hiSetBindKeys, hiSetButtonLabel, hiSetCallbackStatus, hiSetCurrentField, hiSetCurrentIndex, hiSetCurrentWindow, hiSetCursor, hiSetDBoxDefaultLocation, hiSetEncapHistory, hiSetEncapPrompt, hiSetEncapSkillCmd, hiSetFieldEditable, hiSetFieldEnabled, hiSetFilter, hiSetFilterOptions, hiSetFixedMenuSize, hiSetFont, hiSetFormBlock, hiSetFormButtonEnabled, hiSetFormHighlights, hiSetFormMinMaxSize, hiSetFormName, hiSetFormPosition, hiSetFormSize, hiSetFormToDefaults, hiSetIconName, hiSetInsertionPosition, hiSetLayerCyclicValue, hiSetListItemCenter, hiSetListItemVisible, hiSetMenu, hiSetMenuItemCallback, hiSetMenuItemText, hiSetMultiClickTime, hiSetPriority, hiSetProgress, hiSetProgressAndText, hiSetProgressBannerText, hiSetProgressButtonText, hiSetProgressText, hiSetProgressTotalSteps, hiSetScrollBarValue, hiSetShadowMode, hiSetShadowPercent, hiSetSkillDesc, hiSetSplashBackground, hiSetSplashDefaultBackground, hiSetSplashFamily, hiSetSplashIcon, hiSetSplashLicense, hiSetSplashProduct, hiSetTextClass, hiSetTextHighlightColor, hiSetTextSelectAll, hiSetTextSelection, hiSetTopListItem, hiSetTreeItemDescription, hiSetUndoLimit, hiSetUserPreferences, hiSetViewfile, hiSetWinStyle, hiSetWindowAtts, hiSetWindowIcon, hiSetWindowMenu, hiSetWindowName, hiShowBindKeys, hiShowBindKeysByAppType, hiShowBindKeysByWindow, hiShowFieldBorders, hiShowListBox, hiSkillProcess, hiSleep, hiSoftInterrupt, hiStartGenTextIndex, hiStartLog, hiStopChild, hiStoreFormLocation, hiStringToIcon, hiStroke, hiSwitchWindowType, hiSynchronize, hiTextDisplayString, hiTextWidth, hiToggleEnterForm, hiTreeAppendItem, hiTreePrependItem, hiTreeRemoveAllItems, hiTreeRemoveItem, hiTreeRemoveItems, hiTreeTableDeselectAllItems, hiTreeTableDeselectItem, hiTreeTableDeselectItems, hiTreeTableGetExpandedItemCount, hiTreeTableGetExpandedItems, hiTreeTableGetItemCount, hiTreeTableGetItems, hiTreeTableGetSelectedItemCount, hiTreeTableGetSelectedItems, hiTreeTableSelectAllItems, hiTreeTableSelectItem, hiTreeTableSelectItems, hiUndo, hiUnmapWindow, hiUnregCloseProc, hiUnregZoomPanProc, hiUnselectText, hiUnselectTextAll, hiUnselectTextByLoc, hiUnselectTextClass, hiUpdate, hiUpdateFormBlock, hiUpdateTextSelectionColors, hiVectorPan, hiViewfile, hiWait, hiWaitForChild, hiWriteChild, hiZoomAbsoluteScale, hiZoomIn, hiZoomInAtMouse, hiZoomOut, hiZoomOutAtMouse, hiZoomRelativeScale, hiZoomWindowAtMouse, hiZoomWindowAtPoint, histogram, history, hlcheck, hnlAbortNetlist, hnlCDLPrintBJTElement, hnlCDLPrintBSIM3SOIElement, hnlCDLPrintCapElement, hnlCDLPrintCapacitorElement, hnlCDLPrintCds_Thru, hnlCDLPrintDiodeElement, hnlCDLPrintGeneralElement, hnlCDLPrintICIsrcElement, hnlCDLPrintICVsrcElement, hnlCDLPrintInductorElement, hnlCDLPrintIsrcElement, hnlCDLPrintJfetElement, hnlCDLPrintMultiCNPNElement, hnlCDLPrintMultiCPNPElement, hnlCDLPrintMultiENPNElement, hnlCDLPrintMultiEPNPElement, hnlCDLPrintNMOSfetElement, hnlCDLPrintNPNElement, hnlCDLPrintPMOSfetElement, hnlCDLPrintPNPElement, hnlCDLPrintResElement, hnlCDLPrintResistorElement, hnlCDLPrintSchottkyTranElement, hnlCDLPrintTlineElement, hnlCDLPrintVCIsrcElement, hnlCDLPrintVCVsrcElement, hnlCDLPrintVsrcElement, hnlCatIncrementalNetlistFiles, hnlCellExtracted, hnlCellInAllCells, hnlCloseCellFiles, hnlCloseMasterList, hnlCompletePrint, hnlDoInstBased, hnlDoNetBased, hnlFindAllCells, hnlFindAllInstInCell, hnlGenIncludeFile, hnlGetGlobalModelMappedName, hnlGetGlobalNetMappedName, hnlGetMappedInstNames, hnlGetMappedModelNames, hnlGetMappedNames, hnlGetMappedNetNames, hnlGetMasterCells, hnlGetPropVal, hnlGetRoundProp, hnlGetScaleCapacitance, hnlGetScaleMarginalDelay, hnlGetScaleTimeUnit, hnlGetSymbolPropVal, hnlGetTermByName, hnlGetTermNameOfSig, hnlHspiceInstPropVal, hnlHspicePrintInstPropEqVal, hnlHspicePrintInstPropVal, hnlHspicePrintMOSfetModel, hnlHspicePrintNMOSfetElement, hnlIfNoProcedure, hnlIgnoreTerm, hnlInitMap, hnlInitPrint, hnlIsAPatchCord, hnlIsAStoppingCell, hnlIsCellNetlistable, hnlIsCurrentInstStopping, hnlMakeNetlistFileName, hnlMapCellModuleName, hnlMapCellName, hnlMapInstName, hnlMapModelName, hnlMapName, hnlMapNetName, hnlMapTermName, hnlMultipleCells, hnlNameOfSignal, hnlNetNameOnTerm, hnlOpenTopCell, hnlPrintDevices, hnlPrintMessage, hnlPrintNetlist, hnlPrintSignal, hnlPrintString, hnlRunNetlister, hnlScaleCapacitance, hnlScaleMarginalDelay, hnlScaleTimeUnit, hnlSetCellFiles, hnlSetDef, hnlSetPseudoTermDir, hnlSetVars, hnlSortTerms, hnlSortTermsToNets, hnlStartNetlist, hnlStopNetlist, hnlStringToList, hnlVerilogPrintCmos, hnlVerilogPrintNmosPmos, hnlWriteBlockControlFile, hnlWriteMap, hostMode, hostName, i, ic, icChangeEditMode, icCloseWindow, icConnectivityCB, icCreateArraySiteInstances, icCreateBus, icCreateSiteInst, icDbGetCellBndry, icDbInvertTransform, icDiscardEditsMC, icGE2OpusOrien, icGE2OpusinstLoc, icGetAllSiteInstFromSpat, icGetArraySiteInstances, icGetBasicSites, icGetCVCellTypeProp, icGetCellType, icGetGrpByName, icGetGrpByRgn, icGetInstByRgn, icGetRgnName, icGetSiteInstFromSpatbyLoc, icGetSiteInstFromSpatbySiteName, icGetSiteSize, icGetSiteSymmetryProp, icGetTrackInfoByLayer, icHiChangeLayer, icHiCreateBlockage, icHiCreateWire, icHiReShape, icHiRespace, icHiSplit, icHiStretch, icHiStretchRail, icHiWireEditForm, icHierSave, icIsBaseArrayCV, icOpus2GEOrien, icOpus2GEinstLoc, icPathDeselectByAll, icPathSelectByAll, icReDoReDraw, icReturn, icReturnToLevel, icSaveAsMC, icSaveMC, icSetBaseArrayCellTypeProp, icSetDesignWindow, icSetPRengine, icSetSiteCellTypeProp, icSetSitePattCellTypeProp, icSetSiteSize, icSetSiteSymmetryProp, icTfRelease, icTpSaveTemplateFile, icUnDoReDraw, iccConnectToICC, iccDisconnectFromICC, iccDisplayExportForm, iccDisplayImportForm, iccExportCellview, iccImportCellview, iccIsConnected, iccNewRules, iccOpenCurrentCellviewRules, iccOpenRules, iccRegPostExportTrigger, iccRegPostImportTrigger, iccSendCommand, iccSendSkillCommand, iccStartCCAR, iccStartICC, iccStartVCR, iccUnregPostExportTrigger, iccUnregPostImportTrigger, ictMkDir, if, if, iim, iinteg, ilToolBox, im, imag, impHdlDisplay, importSkillVar, inScheme, inSkill, include, includeFile, index, infile, infile, inportp, installDebugger, instring, int, intToChar, integ, integerp, intersect, ip, ip3Plot, ipcActivateBatch, ipcActivateMessages, ipcBatchProcess, ipcBeginProcess, ipcCloseProcess, ipcContProcess, ipcGetExitStatus, ipcGetPid, ipcGetPriority, ipcIsActiveProcess, ipcIsAliveProcess, ipcKillAllProcesses, ipcKillProcess, ipcReadProcess, ipcSetPriority, ipcSkillProcess, ipcSleep, ipcSoftInterrupt, ipcStopProcess, ipcWait, ipcWaitForProcess, ipcWriteProcess, ipn, ipnVRI, ipnVRICurves, iqCreateDummyCell, ir, isBBox, isCallable, isClass, isContextLoaded, isDir, isExecutable, isFile, isFileEncrypted, isFileName, isInfinity, isLargeFile, isLink, isLocation, isMacro, isNaN, isReadable, isWritable, iseCloseSchWindow, iseCloseSimWindow, iseCloseWaveWindow, iseCompleteInteractive, iseEnterNodeNamesList, iseExitSimulator, iseGetExtName, iseGetInputFromEncapWindow, iseGetMappedProbeList, iseGetProbeList, iseInitDefaultWaveWindow, iseInitSchematicWindow, iseInitSimWindow, iseInitWaveWindow, iseNetExtNameCdsName, iseOpenWindows, isePrintName, isePrintNameCB, isePrintSimulatorCommand, iseReadWaves, iseReleaseNodeFrom, iseSearchForASchWindow, iseSendOutputToEncapHistory, iseSetEncapBindKeys, iseSetNodeTo, iseSimulate, iseStartInteractive, iseStartSimulator, iseUpdateNetlist, iseUpdateStimulus, iterVsValue, ivCallProc, ivCompilingFor, ivConcICe, ivCreatePCells, ivDRC, ivERC, ivExtract, ivGetCellView, ivLVS, ivVerify, kf, killJob, lambda, last, layerDefinitions, layerFunctions, layerMfgResolutions, layerRoutingGrids, layerRules, lconc, ldtrDefRead, ldtrDefWrite, ldtrGenLefoutList, ldtrLefRead, ldtrLefWrite, leAttachFig, leAutoRoute, leChopShape, leClearAllRuler, leCloseWindow, leConvertShapeToPolygon, leCreateAutoInstPin, leCreateAutoPin, leCreateContact, leCreatePath, leCreatePin, leCreateRuler, leDRDUpdateRuleOptions, leDefineExternalPins, leDefineInternalPins, leDefineMPPTemplate, leDefinePPPins, leDefineWeaklyConnectedPins, leDeiconifyLSW, leDescend, leDesignSummary, leDrdDefineDerivedLayers, leDrdSetRuleOptions, leEIPZoomAbsoluteScale, leEditDesignProperties, leEditInPlace, leEnvLoad, leFlattenInst, leFullSelectFigOfSelSet, leGetAllCdsViaDefaultParam, leGetContactDefaultParam, leGetContactNameArray, leGetContactParam, leGetContactRule, leGetEntryLayer, leGetEnv, leGetLSWBBox, leGetObjectSelectable, leGetObjectVisible, leGetRefPoint, leGetValidLayerList, leHiAbout, leHiAddShapeToNet, leHiAlign, leHiAttach, leHiChop, leHiClearRuler, leHiConvertShapeToPolygon, leHiCopy, leHiCreateBend, leHiCreateChoiceOfPin, leHiCreateCircle, leHiCreateContact, leHiCreateDevice, leHiCreateDonut, leHiCreateEllipse, leHiCreateGuardRing, leHiCreateInst, leHiCreateLabel, leHiCreatePath, leHiCreatePin, leHiCreatePinsFromLabels, leHiCreatePolygon, leHiCreateRect, leHiCreateRuler, leHiCreateSymDev, leHiCreateSymPin, leHiCreateTaper, leHiCreateTrl, leHiDelete, leHiDeleteAllAreaViewLevel, leHiDeleteAreaViewLevel, leHiDeleteShapeFromNet, leHiDescend, leHiEditDRDOptions, leHiEditDRDRuleOptions, leHiEditDisplayOptions, leHiEditDynamicMeasurementOptions, leHiEditEditorOptions, leHiEditInPlace, leHiEditProp, leHiEditTbOptions, leHiFlatten, leHiLayerGen, leHiLayerTap, leHiMakeCell, leHiMarkNet, leHiMerge, leHiModifyCorner, leHiMousePopUp, leHiMove, leHiMoveOrigin, leHiOptionLayer, leHiPaste, leHiPlotQueueStatus, leHiPropagateNets, leHiQuery, leHiReShape, leHiRemasterInstances, leHiRotate, leHiSearch, leHiSetAreaViewLevel, leHiSetRefPoint, leHiSetValidLayer, leHiShowSelSet, leHiSize, leHiSplit, leHiStretch, leHiSubmitPlot, leHiSummary, leHiTree, leHiUnmarkNet, leHiUnmarkNetAll, leHiYank, leIconifyLSW, leIsAnyContact, leIsAnyContactMaster, leIsContact, leIsContactMaster, leIsContactName, leIsFigSelectable, leIsInstSelectable, leIsLSWIconified, leIsLayerSelectable, leIsLayerValid, leIsLayerVisible, leIsPinSelectable, leIsRefPointActive, leLSWShowLayers, leLSWShowObjects, leLayerAnd, leLayerAndNot, leLayerOr, leLayerSize, leLayerXor, leLswLayers, leMakeCell, leMergeShapes, leModifyCorner, leMoveCellViewOrigin, leMoveCursor, leMoveCursorToRefPoint, lePasteFigs, lePlot, leRaiseLSW, leRegUserLayerSelectionFilter, leRegUserObjectSelectionFilter, leRemapLSW, leRemasterInstances, leReplace, leReplaceAnyInstMaster, leResizeLSW, leRules, leSearchHierarchy, leSelBoxOrStretch, leSetAllLayerSelectable, leSetAllLayerValid, leSetAllLayerVisible, leSetAllObjectsSelectable, leSetAllObjectsVisible, leSetContactParam, leSetEntryLayer, leSetEnv, leSetFormSnapMode, leSetInstSelectable, leSetLSWBBox, leSetLSWFilter, leSetLayerAttributes, leSetLayerSelectable, leSetLayerValid, leSetLayerVisible, leSetObjectSelectable, leSetObjectVisible, leSetPinSelectable, leSetRefPoint, leSetRefPointInactive, leSizeShape, leSpaceBarFunc, leSplitShape, leStretchShape, leToggleDrdMode, leToggleGravity, leToggleMaintainConnections, leUnmapLSW, leUnregUserLayerSelectionFilter, leUnregUserObjectSelectionFilter, leWeAddRotatedVia, leWeAddVia, leWeCheckRoutes, leWeDisableAsManyAsFitMenuItem, leWeDisableCopyRouteMirrorMenuItem, leWeFinishRoute, leWeHiCheckRoutes, leWeHiCopyRoute, leWeHiCreateWire, leWeHiCriticWire, leWeHiPull, leWeHiRouteOptions, leWeHiSplit, leWeHiStretch, leWePathWidthMenuCB, leWePickUpDroppedWires, leWeRefresh, leWeReports, leWeSetGatherBusWiresMenuItemText, leWeSetInteractiveCheckingText, leYankFigs, leZoomToPoint, leZoomToSelSet, leftEdge, leftshift, leiDiscardEdits, length, leqp, lessp, let, let, letrec, letseq, linRg, lineread, linereadstring, list, listFunctions, listToVector, listVariables, listp, lmCheckTerm, lmCheckView, lmCloseLib, lmDefCell, lmDefTermProp, lmDefViewProp, lmDeleteTermProp, lmDeleteViewProp, lmGetValue, lmLoadData, lmOpenLib, lmPrintLibTermProp, lmPrintLibViewProp, lmPrintTerm, lmPrintTermProp, lmPrintViewProp, lmReset, lmSimView, lmgrAddMenuItems, lmgrCreateMenu, lmgrCreateMenuItem, lmgrDefineInits, lmgrDeleteMenuItems, lmgrDisplayMessage, lmgrGetObject, lmgrInsertMenuItems, lmgrManageMenuItems, lmgrQueryNamedObjects, lmgrSensitizeMenuItems, lmgrSetObject, lmgrVerbose, ln, loCellSubstitutionMap, loLayerMap, loLibraryMap, loMapParameters, loSchematicParams, loSetAnnotateLPP, loSplitViaFunctionMap, load, load, loadContext, loadDcf, loadPcf, loadi, loadstring, log, log, log10, log10, logRg, lowerCase, lowerLeft, lpCreatePin, lsb, lshift, lxCmdOptions, lxCmdShiftOptions, lxEditPartitioning, lxEditPinPlacement, lxEditPlacementStyle, lxExtractLayers, lxGenFromSource, lxGenerateFinish, lxGenerateGetAvailableBoundaryLPPs, lxGenerateGetAvailablePinLPPs, lxGenerateGetPinNets, lxGenerateSetAreaEstimationOptions, lxGenerateSetBoundaryOptions, lxGenerateSetGenerateOptions, lxGenerateSetNetPinSpecs, lxGenerateStart, lxGetLXInfo, lxHiAbout, lxHiAlign, lxHiAutoPlace, lxHiChain, lxHiClone, lxHiConnectInstPin, lxHiCreateInstFromSch, lxHiCreateMPP, lxHiEditComponentTypes, lxHiExtract, lxHiHideIncNets, lxHiLockSelected, lxHiMoveAutomatically, lxHiPlcDisableCongestionDisplay, lxHiPlcEnableCongestionDisplay, lxHiProbe, lxHiReInitDesign, lxHiSetCorrespondence, lxHiSetOptions, lxHiShowIncNets, lxHiStack, lxHiSwapComps, lxHiUnlockSelected, lxHiUpdateCellViewPair, lxHiUpdateComponentsAndNets, lxHiUpdateDeviceCorr, lxHiUpdateLayoutParameters, lxHiUpdateSchematicParameters, lxHiUpdateSwitchViews, lxHiVerifyDesign, lxHiVerifyStatus, lxHierExtract, lxMPPTemplates, lxMovePorts, lxNoOverlapLayers, lxNoOverlapLayers, lxPermPermutePins, lxPlcAppendPlaceSetupFieldValue, lxPlcGetPlaceSetupFieldValue, lxPlcIsPlaceSetupFieldEnabled, lxPlcReplacePlaceSetupField, lxPlcSetPlaceSetupFieldEnableState, lxPlcSetPlaceSetupFieldValue, lxProbeRemoveAll, lxRules, lxSetConnRef, lxToggleIncNets, lxToggleLocalAbutment, lxTpSaveTemplate, lxTpTemplateLoadAll, mag, makeInstance, makeTable, makeTempFileName, makeVector, make_, map, mapc, mapcan, mapcar, maplist, max, max, measureTime, member, menuTrigger, mfgGridResolution, min, min, minus, minusp, mod, mod, modelFile, modulo, monitor, monteCarlo, monteCorrelate, monteDisplay, monteExpr, monteOutputs, monteResults, monteRun, monteSelectResults, mouseAddPt, mouseAddSelectPt, mouseFinishPt, mousePopUp, mouseSingleSelectPt, mouseSubSelectPt, mouseUndoPt, mprocedure, msgHelp, mspDisplaySetPartSetupForm, mspEditIEProps, nc, nconc, ncons, needNCells, negativep, neq, nequal, netlist, newWindow, newline, newline, next, nextMethodp, nindex, nlDisplayOption, nlError, nlGetCdf, nlGetCellName, nlGetCurrentSwitchMaster, nlGetDesign, nlGetFormatter, nlGetGlobalNets, nlGetId, nlGetLibName, nlGetModelName, nlGetNetlistDir, nlGetNetlister, nlGetNumberOfBits, nlGetOption, nlGetOptionNameList, nlGetParamList, nlGetParamStringValue, nlGetSignalList, nlGetSimName, nlGetSimTerminalNets, nlGetSwitchViewList, nlGetTerminalList, nlGetTerminalNets, nlGetTerminalSignalName, nlGetToolName, nlGetTopCellName, nlGetTopLibName, nlGetTopViewName, nlIncludeSrcFile, nlIncludeVerilogaFile, nlInfo, nlInitialize, nlIsModelNameInherited, nlMapGlobalNet, nlObjError, nlPrintComment, nlPrintComments, nlPrintFooter, nlPrintHeader, nlPrintHeaderComments, nlPrintIndentString, nlPrintInst, nlPrintInstComments, nlPrintInstEnd, nlPrintInstName, nlPrintInstParameters, nlPrintInstSignals, nlPrintModelName, nlPrintString, nlPrintStringNoFold, nlPrintSubcktBegin, nlPrintSubcktEnd, nlPrintSubcktFooter, nlPrintSubcktFooterComments, nlPrintSubcktHeader, nlPrintSubcktHeaderComments, nlPrintSubcktName, nlPrintSubcktParameters, nlPrintSubcktTerminalList, nlPrintTopCellFooter, nlPrintTopCellFooterComments, nlPrintTopCellHeader, nlPrintTopCellHeaderComments, nlSetOption, nlSetPcellName, nlSocketAdjustNameCase, nlSocketGetPropValue, nlSocketNetlistProc, nlSocketPrintComments, nlTranslateFlatIEPathName, nlWarning, nlambda, nmp, nmp, nmpCDBAFlatToCDBA, nmpCDBAToCDBAFlat, nmpCDBAToConcept, nmpCDBAToDef, nmpCDBAToGcf, nmpCDBAToGenesis, nmpCDBAToIdent, nmpCDBAToLef, nmpCDBAToPrint, nmpCDBAToSdf, nmpCDBAToSpectre, nmpCDBAToSpectreHDL, nmpCDBAToSpef, nmpCDBAToSpf, nmpCDBAToSpice, nmpCDBAToVHDL, nmpCDBAToVHDL87, nmpCDBAToVHDLAMS, nmpCDBAToVerilog, nmpCDBAToVerilogA, nmpCDBAToVerilogAMS, nmpConceptToCDBA, nmpConceptToDef, nmpConceptToGcf, nmpConceptToGenesis, nmpConceptToIdent, nmpConceptToLef, nmpConceptToPrint, nmpConceptToSdf, nmpConceptToSpectre, nmpConceptToSpectreHDL, nmpConceptToSpef, nmpConceptToSpf, nmpConceptToSpice, nmpConceptToVHDL, nmpConceptToVHDL87, nmpConceptToVHDLAMS, nmpConceptToVerilog, nmpConceptToVerilogA, nmpConceptToVerilogAMS, nmpDefToCDBA, nmpDefToConcept, nmpDefToGcf, nmpDefToGenesis, nmpDefToIdent, nmpDefToLef, nmpDefToPrint, nmpDefToSdf, nmpDefToSpectre, nmpDefToSpectreHDL, nmpDefToSpef, nmpDefToSpf, nmpDefToSpice, nmpDefToVHDL, nmpDefToVHDL87, nmpDefToVHDLAMS, nmpDefToVerilog, nmpDefToVerilogA, nmpDefToVerilogAMS, nmpGcfToCDBA, nmpGcfToConcept, nmpGcfToDef, nmpGcfToGenesis, nmpGcfToIdent, nmpGcfToLef, nmpGcfToPrint, nmpGcfToSdf, nmpGcfToSpectre, nmpGcfToSpectreHDL, nmpGcfToSpef, nmpGcfToSpf, nmpGcfToSpice, nmpGcfToVHDL, nmpGcfToVHDL87, nmpGcfToVHDLAMS, nmpGcfToVerilog, nmpGcfToVerilogA, nmpGcfToVerilogAMS, nmpGenesisToCDBA, nmpGenesisToConcept, nmpGenesisToDef, nmpGenesisToGcf, nmpGenesisToIdent, nmpGenesisToLef, nmpGenesisToPrint, nmpGenesisToSdf, nmpGenesisToSpectre, nmpGenesisToSpectreHDL, nmpGenesisToSpef, nmpGenesisToSpf, nmpGenesisToSpice, nmpGenesisToVHDL, nmpGenesisToVHDL87, nmpGenesisToVHDLAMS, nmpGenesisToVerilog, nmpGenesisToVerilogA, nmpGenesisToVerilogAMS, nmpGetSpaceNames, nmpIdentTo, nmpIdentToCDBA, nmpIdentToConcept, nmpIdentToDef, nmpIdentToGcf, nmpIdentToGenesis, nmpIdentToLef, nmpIdentToPrint, nmpIdentToSdf, nmpIdentToSpectre, nmpIdentToSpectreHDL, nmpIdentToSpef, nmpIdentToSpf, nmpIdentToSpice, nmpIdentToVHDL, nmpIdentToVHDL87, nmpIdentToVHDLAMS, nmpIdentToVerilog, nmpIdentToVerilogA, nmpIdentToVerilogAMS, nmpIsLegal, nmpIsLegalCDBA, nmpIsLegalConcept, nmpIsLegalDef, nmpIsLegalGcf, nmpIsLegalGenesis, nmpIsLegalLef, nmpIsLegalPrint, nmpIsLegalSdf, nmpIsLegalSpectre, nmpIsLegalSpectreHDL, nmpIsLegalSpef, nmpIsLegalSpf, nmpIsLegalSpice, nmpIsLegalVHDL, nmpIsLegalVHDL87, nmpIsLegalVHDLAMS, nmpIsLegalVerilog, nmpIsLegalVerilogA, nmpIsLegalVerilogAMS, nmpLefToCDBA, nmpLefToConcept, nmpLefToDef, nmpLefToGcf, nmpLefToGenesis, nmpLefToIdent, nmpLefToPrint, nmpLefToSdf, nmpLefToSpectre, nmpLefToSpectreHDL, nmpLefToSpef, nmpLefToSpf, nmpLefToSpice, nmpLefToVHDL, nmpLefToVHDL87, nmpLefToVHDLAMS, nmpLefToVerilog, nmpLefToVerilogA, nmpLefToVerilogAMS, nmpPrintToCDBA, nmpPrintToConcept, nmpPrintToDef, nmpPrintToGcf, nmpPrintToGenesis, nmpPrintToIdent, nmpPrintToLef, nmpPrintToSdf, nmpPrintToSpectre, nmpPrintToSpectreHDL, nmpPrintToSpef, nmpPrintToSpf, nmpPrintToSpice, nmpPrintToVHDL, nmpPrintToVHDL87, nmpPrintToVHDLAMS, nmpPrintToVerilog, nmpPrintToVerilogA, nmpPrintToVerilogAMS, nmpSdfToCDBA, nmpSdfToConcept, nmpSdfToDef, nmpSdfToGcf, nmpSdfToGenesis, nmpSdfToIdent, nmpSdfToLef, nmpSdfToPrint, nmpSdfToSpectre, nmpSdfToSpectreHDL, nmpSdfToSpef, nmpSdfToSpf, nmpSdfToSpice, nmpSdfToVHDL, nmpSdfToVHDL87, nmpSdfToVHDLAMS, nmpSdfToVerilog, nmpSdfToVerilogA, nmpSdfToVerilogAMS, nmpSpectreHDLToCDBA, nmpSpectreHDLToConcept, nmpSpectreHDLToDef, nmpSpectreHDLToGcf, nmpSpectreHDLToGenesis, nmpSpectreHDLToIdent, nmpSpectreHDLToLef, nmpSpectreHDLToPrint, nmpSpectreHDLToSdf, nmpSpectreHDLToSpectre, nmpSpectreHDLToSpef, nmpSpectreHDLToSpf, nmpSpectreHDLToSpice, nmpSpectreHDLToVHDL, nmpSpectreHDLToVHDL87, nmpSpectreHDLToVHDLAMS, nmpSpectreHDLToVerilog, nmpSpectreHDLToVerilogA, nmpSpectreHDLToVerilogAMS, nmpSpectreToCDBA, nmpSpectreToConcept, nmpSpectreToDef, nmpSpectreToGcf, nmpSpectreToGenesis, nmpSpectreToIdent, nmpSpectreToLef, nmpSpectreToPrint, nmpSpectreToSdf, nmpSpectreToSpectreHDL, nmpSpectreToSpef, nmpSpectreToSpf, nmpSpectreToSpice, nmpSpectreToVHDL, nmpSpectreToVHDL87, nmpSpectreToVHDLAMS, nmpSpectreToVerilog, nmpSpectreToVerilogA, nmpSpectreToVerilogAMS, nmpSpefToCDBA, nmpSpefToConcept, nmpSpefToDef, nmpSpefToGcf, nmpSpefToGenesis, nmpSpefToIdent, nmpSpefToLef, nmpSpefToPrint, nmpSpefToSdf, nmpSpefToSpectre, nmpSpefToSpectreHDL, nmpSpefToSpf, nmpSpefToSpice, nmpSpefToVHDL, nmpSpefToVHDL87, nmpSpefToVHDLAMS, nmpSpefToVerilog, nmpSpefToVerilogA, nmpSpefToVerilogAMS, nmpSpfToCDBA, nmpSpfToConcept, nmpSpfToDef, nmpSpfToGcf, nmpSpfToGenesis, nmpSpfToIdent, nmpSpfToLef, nmpSpfToPrint, nmpSpfToSdf, nmpSpfToSpectre, nmpSpfToSpectreHDL, nmpSpfToSpef, nmpSpfToSpice, nmpSpfToVHDL, nmpSpfToVHDL87, nmpSpfToVHDLAMS, nmpSpfToVerilog, nmpSpfToVerilogA, nmpSpfToVerilogAMS, nmpSpiceToCDBA, nmpSpiceToConcept, nmpSpiceToDef, nmpSpiceToGcf, nmpSpiceToGenesis, nmpSpiceToIdent, nmpSpiceToLef, nmpSpiceToPrint, nmpSpiceToSdf, nmpSpiceToSpectre, nmpSpiceToSpectreHDL, nmpSpiceToSpef, nmpSpiceToSpf, nmpSpiceToVHDL, nmpSpiceToVHDL87, nmpSpiceToVHDLAMS, nmpSpiceToVerilog, nmpSpiceToVerilogA, nmpSpiceToVerilogAMS, nmpVHDL87ToCDBA, nmpVHDL87ToConcept, nmpVHDL87ToDef, nmpVHDL87ToGcf, nmpVHDL87ToGenesis, nmpVHDL87ToIdent, nmpVHDL87ToLef, nmpVHDL87ToPrint, nmpVHDL87ToSdf, nmpVHDL87ToSpectre, nmpVHDL87ToSpectreHDL, nmpVHDL87ToSpef, nmpVHDL87ToSpf, nmpVHDL87ToSpice, nmpVHDL87ToVHDL, nmpVHDL87ToVHDLAMS, nmpVHDL87ToVerilog, nmpVHDL87ToVerilogA, nmpVHDL87ToVerilogAMS, nmpVHDLAMSToCDBA, nmpVHDLAMSToConcept, nmpVHDLAMSToDef, nmpVHDLAMSToGcf, nmpVHDLAMSToGenesis, nmpVHDLAMSToIdent, nmpVHDLAMSToLef, nmpVHDLAMSToPrint, nmpVHDLAMSToSdf, nmpVHDLAMSToSpectre, nmpVHDLAMSToSpectreHDL, nmpVHDLAMSToSpef, nmpVHDLAMSToSpf, nmpVHDLAMSToSpice, nmpVHDLAMSToVHDL, nmpVHDLAMSToVHDL87, nmpVHDLAMSToVerilog, nmpVHDLAMSToVerilogA, nmpVHDLAMSToVerilogAMS, nmpVHDLToCDBA, nmpVHDLToConcept, nmpVHDLToDef, nmpVHDLToGcf, nmpVHDLToGenesis, nmpVHDLToIdent, nmpVHDLToLef, nmpVHDLToPrint, nmpVHDLToSdf, nmpVHDLToSpectre, nmpVHDLToSpectreHDL, nmpVHDLToSpef, nmpVHDLToSpf, nmpVHDLToSpice, nmpVHDLToVHDL87, nmpVHDLToVHDLAMS, nmpVHDLToVerilog, nmpVHDLToVerilogA, nmpVHDLToVerilogAMS, nmpVerilogAMSToCDBA, nmpVerilogAMSToConcept, nmpVerilogAMSToDef, nmpVerilogAMSToGcf, nmpVerilogAMSToGenesis, nmpVerilogAMSToIdent, nmpVerilogAMSToLef, nmpVerilogAMSToPrint, nmpVerilogAMSToSdf, nmpVerilogAMSToSpectre, nmpVerilogAMSToSpectreHDL, nmpVerilogAMSToSpef, nmpVerilogAMSToSpf, nmpVerilogAMSToSpice, nmpVerilogAMSToVHDL, nmpVerilogAMSToVHDL87, nmpVerilogAMSToVHDLAMS, nmpVerilogAMSToVerilog, nmpVerilogAMSToVerilogA, nmpVerilogAToCDBA, nmpVerilogAToConcept, nmpVerilogAToDef, nmpVerilogAToGcf, nmpVerilogAToGenesis, nmpVerilogAToIdent, nmpVerilogAToLef, nmpVerilogAToPrint, nmpVerilogAToSdf, nmpVerilogAToSpectre, nmpVerilogAToSpectreHDL, nmpVerilogAToSpef, nmpVerilogAToSpf, nmpVerilogAToSpice, nmpVerilogAToVHDL, nmpVerilogAToVHDL87, nmpVerilogAToVHDLAMS, nmpVerilogAToVerilog, nmpVerilogAToVerilogAMS, nmpVerilogToCDBA, nmpVerilogToConcept, nmpVerilogToDef, nmpVerilogToGcf, nmpVerilogToGenesis, nmpVerilogToIdent, nmpVerilogToLef, nmpVerilogToPrint, nmpVerilogToSdf, nmpVerilogToSpectre, nmpVerilogToSpectreHDL, nmpVerilogToSpef, nmpVerilogToSpf, nmpVerilogToSpice, nmpVerilogToVHDL, nmpVerilogToVHDL87, nmpVerilogToVHDLAMS, nmpVerilogToVerilogA, nmpVerilogToVerilogAMS, nodeset, noise, noiseSummary, not, nprocedure, nth, nthcdr, nthelem, null, numOpenFiles, numberp, ocnDisplay, ocnGetWaveformTool, ocnHelp, ocnMeasBandwidth, ocnMeasCompare, ocnMeasConv, ocnMeasDelay, ocnMeasDerivative, ocnMeasDeviceNoise, ocnMeasDnl, ocnMeasEval, ocnMeasFFT, ocnMeasFindWhen, ocnMeasFreq, ocnMeasFreqSpan, ocnMeasGBWProd, ocnMeasGain, ocnMeasGainMargin, ocnMeasIdft, ocnMeasInl, ocnMeasIntegral, ocnMeasInterval, ocnMeasOPBip, ocnMeasOPMOS, ocnMeasOutputSwing, ocnMeasOvershoot, ocnMeasPeriod, ocnMeasPhaseMargin, ocnMeasProbe, ocnMeasSample, ocnMeasSettlingTime, ocnMeasShift, ocnMeasSlewRate, ocnMeasSpectrum, ocnMeasUnityGainFreq, ocnPrint, ocnWaveformTool, ocnYvsYplot, ocnm, ocnmArgCheck, ocnmCheckErrMsg, ocnmCheckMeasFile, ocnmClearErrMsg, ocnmClip, ocnmCloseMeasFile, ocnmDumpErr, ocnmDumpRes, ocnmFindDomain, ocnmGetArgs, ocnmGetErrMsg, ocnmGetInfoString, ocnmGetKeepDomain, ocnmGetMeasFile, ocnmGetMeasKeys, ocnmGetMultipleResultType, ocnmMakeOccArg, ocnmMakeRangeArg, ocnmOcc, ocnmOpenMeasFile, ocnmReturnResult, ocnmSetDebug, ocnmSetErrMsg, ocnmSetInfoString, ocnmSetKeepDomain, ocnmSetMeasKeys, ocnmSetMultipleResultType, oddp, off, onep, openResults, openportp, optimizeAlgoControl, optimizeGoal, optimizePlotOption, optimizeRun, optimizeVar, option, or, orderedCharacterizationRules, orderedSpacingRules, otherp, outfile, outfile, outportp, outputParams, outputs, overshoot, pairp, paramAnalysis, paramRun, parseString, path, pcColinearPoints, pcConcatOrient, pcDefineCondition, pcDefineInheritParam, pcDefinePCell, pcDefineParamCell, pcDefineParamLabel, pcDefineParamLayer, pcDefineParamPath, pcDefineParamPolygon, pcDefineParamProp, pcDefineParamRect, pcDefineParamRefPointObject, pcDefinePathRefPointObject, pcDefineRepeat, pcDefineSteppedObject, pcDefineStretchLine, pcDeleteCondition, pcDeleteParam, pcDeleteParamLayer, pcDeleteParamProp, pcDeleteParamShape, pcDeleteRefPoint, pcDeleteRepeat, pcDeleteSteppedObject, pcExprToProp, pcExprToString, pcFilterPoints, pcFix, pcGetBendAngle, pcGetConditions, pcGetInheritParamDefn, pcGetInheritParams, pcGetOffsetPath, pcGetOffsetPolygon, pcGetParamLabelDefn, pcGetParamLabels, pcGetParamLayerDefn, pcGetParamLayers, pcGetParamProps, pcGetParamShapeDefn, pcGetParamShapes, pcGetParameters, pcGetPathRefPoint, pcGetRefPointDefn, pcGetRefPoints, pcGetRepeatDefn, pcGetRepeats, pcGetStepDirection, pcGetSteppedObjectDefn, pcGetSteppedObjects, pcGetStretchDefn, pcGetStretchSummary, pcGetStretches, pcGrowBox, pcGrowPoints, pcHICompileToSkill, pcHIDefineCondition, pcHIDefineInheritedParameter, pcHIDefineLabel, pcHIDefineLayer, pcHIDefineParamCell, pcHIDefineParamRefPointObject, pcHIDefineParameterizedShape, pcHIDefinePathRefPointObject, pcHIDefineProp, pcHIDefineRepeat, pcHIDefineSteppedObject, pcHIDefineStretch, pcHIDeleteCondition, pcHIDeleteLayer, pcHIDeleteParameterizedShape, pcHIDeleteProp, pcHIDeleteRefPointObject, pcHIDeleteRepeat, pcHIDeleteSteppedObject, pcHIDisplayCondition, pcHIDisplayInheritedParameter, pcHIDisplayLayer, pcHIDisplayParameterizedShape, pcHIDisplayParams, pcHIDisplayProp, pcHIDisplayRefPointObject, pcHIDisplayRepeat, pcHIDisplaySteppedObject, pcHIEditParameters, pcHIModifyCondition, pcHIModifyLabel, pcHIModifyLayer, pcHIModifyParams, pcHIModifyRefPointObject, pcHIModifyRepeat, pcHIModifySteppedObject, pcHIModifyStretchLine, pcHIQualifyStretchLine, pcHIRedefineStretchLine, pcHISummarizeParams, pcModifyParam, pcRedefineStretchLine, pcRestrictStretchToObjects, pcRound, pcSetFTermWidth, pcSkillGen, pcStepAlongShape, pcTechFile, pcUserAdjustParameters, pcUserGenerateArray, pcUserGenerateInstance, pcUserGenerateInstancesOfMaster, pcUserGenerateLPP, pcUserGeneratePin, pcUserGenerateProperty, pcUserGenerateShape, pcUserGenerateTerminal, pcUserInitRepeat, pcUserPostProcessCellView, pcUserPostProcessObject, pcUserPreProcessCellView, pcUserSetTermNetName, pcdbClose, pcdbGetIgnoredInstMasterGen, pcdbGetInstGen, pcdbGetInstMasterGen, pcdbInstMasterCell, pcdbInstMasterCellFixed, pcdbInstMasterLib, pcdbInstMasterLibFixed, pcdbInstMasterView, pcdbInstMasterViewFixed, pcdbInstName, pcdbInstWidth, pcdbNextInst, pcdbNextInstMaster, pcdbOpen, pcdbSave, pcdbViewHasPcdbInfo, peakToPeak, pfile, phase, phaseDeg, phaseDegUnwrapped, phaseMargin, phaseNoise, phaseRad, phaseRadUnwrapped, physicalRules, piCellNameMap, piCifLayerMap, piLayerMap, piTextMap, pipoDisplay, pipoErrShapesHandler, plist, plot, plotStyle, plus, plusp, poCellNameMap, poCifLayerMap, poParamCellNameMap, poTextMap, portp, postInstallTrigger, postdecrement, postincrement, pow, pp, pprint, prGenViaRules, prMastersliceLayers, prMaxStackVias, prNonDefaultRules, prOverlapLayer, prRoutingLayers, prRoutingOffset, prRoutingPitch, prRules, prStackVias, prViaRules, prViaTypes, praBlockPlace, praGEBlockPlace, praOptimizeIO, prasrGetInstTransform, preXY, predecrement, preincrement, prependInstallPath, prependPath, print, printFunctions, printObject, printVariables, printViolations, printf, printf, printlev, println, println, printstruct, probeTrigger, procedure, procedurep, profile, profileReset, profileSummary, prog, prog1, prog2, progn, prsAlignPins, prsDeselectAll, prsDeselectAllTypes, prsDeselectBusByAll, prsDeselectPinByAll, prsFindBusByNamePat, prsFindInstByNamePat, prsFindInstPinByNamePat, prsFindNetByNamePat, prsMouseSelBox, prsMouseSelPt, prsPlcAlignCol, prsPlcAlignRow, prsSelGetBusPropMatch, prsSelGetInstPropMatch, prsSelGetNetPropMatch, prsSelGetPinPropMatch, prsSelectAll, prsSelectArea, prsSelectBusByAll, prsSelectByConnTo, prsSelectPinByAll, prsSelectPoint, prsSetofPin, prsShowSelSet, prsSwitchView, prsSwitchViewMC, prsdbGetBusConnectedTo, prsdbGetInstConnectedTo, prsdbGetNetConnectedTo, prshiSelectByName, prtDEFInCB, prtDEFOutCB, prtLEFInCB, prtLEFOutCB, psConfigLoaded, psLoadCdsPlotInit, psQueryPaperSize, psQueryPaperSizes, psQueryPlotterElement, psQueryPlotters, psQueueStatus, psQueueStatus, psd, psdbb, putd, putprop, putpropq, putpropqq, pv, pwd, pzPlot, pzSummary, quote, quotient, random, random, range, read, readTable, readstring, real, realp, regExitAfter, regExitBefore, remExitProc, remainder, remd, remdq, remoteDir, remove, removeLabel, remprop, remq, renameFile, report, residual, restore, resultParam, results, resultsDir, resume, resumeJob, return, reverse, rexCompile, rexExecute, rexMagic, rexMatchAssocList, rexMatchList, rexMatchp, rexReplace, rexSubstitute, rightEdge, rightshift, rindex, riseTime, rms, rmsNoise, rodAddMPPChopHole, rodAddPoints, rodAddToX, rodAddToY, rodAlign, rodAssignHandleToParameter, rodCheck, rodCheckAlignData, rodCreateHandle, rodCreatePath, rodCreatePolygon, rodCreateRect, rodDeleteHandle, rodFillBBoxWithRects, rodGetHandle, rodGetNamedShapes, rodGetObj, rodIsFigNameUnused, rodIsHandle, rodIsObj, rodNameShape, rodPointX, rodPointY, rodSubPoints, rodUnAlign, rodUnNameShape, root, round, round, routeSpecs, rplaca, rplacd, rshift, ruleContactDevice, run, runsim, sample, save, saveContext, saveOption, scAddCell, scAddColumn, scAddMaster, scAddRow, scBias, scCellHeight, scCellWidth, scCellsPerColumn, scCellsPerRow, scChangeBias, scChangeCell, scChangeIcon, scChangeMaster, scChangeOrient, scCompileArray, scDeleteCell, scDeleteColumn, scDeletePlane, scDeleteRow, scDrawPlane, scGetPlanePoint, scIcon, scMaster, scOrient, scPlane, scPromotePin, scResizePlane, scStretchEmptyCell, scSwapCells, scSwapColumns, scSwapRows, scUnusedIcon, scatterplot, schAddSelectPt, schCheck, schCheckHier, schCheckHierConfig, schCloneSymbol, schCmdOption, schComputePinRef, schCopy, schCreateInst, schCreateInstBox, schCreateNetExpression, schCreateNoteLabel, schCreateNoteShape, schCreatePin, schCreateSheet, schCreateSymbolLabel, schCreateSymbolPin, schCreateSymbolShape, schCreateWire, schCreateWireLabel, schDelete, schDeleteIndex, schDeleteSheet, schDeselectAllFig, schDirectEdit, schDrawSymbolPin, schEditPinOrder, schEditSheetSize, schExistsEditCap, schExtendSelSet, schExtendSelectPt, schExtractConn, schExtractStatus, schGetBundleDisplayMode, schGetCellViewListInSearchScope, schGetEnv, schGetMatchingObjects, schGetPinOrder, schGetPostCheckTriggers, schGlueLabel, schHDLReturn, schHdlPrintFile, schHdlPrintVars, schHiAbout, schHiAlternateView, schHiCellViewProperty, schHiChangeEditMode, schHiCheck, schHiCheckAndSave, schHiCheckHier, schHiCloneSymbol, schHiComputePinRef, schHiCopy, schHiCreateBlockInst, schHiCreateInst, schHiCreateInstBox, schHiCreateKanjiSymbol, schHiCreateNetExpression, schHiCreateNoteLabel, schHiCreateNoteShape, schHiCreatePatchcord, schHiCreatePin, schHiCreateSheet, schHiCreateSimBox, schHiCreateSymbolLabel, schHiCreateSymbolPin, schHiCreateSymbolShape, schHiCreateWire, schHiCreateWireLabel, schHiDelete, schHiDeleteIndex, schHiDeleteSheet, schHiDescend, schHiDescendEdit, schHiDescendRead, schHiDisplayOptions, schHiDrawSymbolPin, schHiEditInPlace, schHiEditPinOrder, schHiEditSheetSize, schHiEditTitleBlock, schHiEditorOptions, schHiEnvSaveLoad, schHiExtractConn, schHiFind, schHiFindMarker, schHiFollowPin, schHiGotoSheet, schHiGridOptions, schHiHiliteLabel, schHiInstToView, schHiMousePopUp, schHiMove, schHiNetExprAvailProps, schHiNetExprEvalNames, schHiNewCellView, schHiObjectProperty, schHiOpenCellView, schHiOpenOtherView, schHiOpenSymbolOrSchematicView, schHiPinListToView, schHiPlot, schHiPlotQueueStatus, schHiRenumberAllSheet, schHiRenumberInstances, schHiRenumberSheet, schHiReplace, schHiResetInvisibleLabels, schHiReturn, schHiReturnToTop, schHiRotate, schHiRouteFlightLine, schHiSRC, schHiSaveCellView, schHiSelectAll, schHiSelectByProperty, schHiSetSymbolOrigin, schHiShowScope, schHiSolder, schHiStretch, schHiSymStretch, schHiUpdatePinOrder, schHiUpdatePinsFromView, schHiVHDLProperty, schHiVIC, schHiVICAndSave, schHiViewToView, schInhConFind, schInhConSet, schInstToView, schInstallHDL, schIsHDLCapEnabled, schIsIndexCV, schIsSchEditOk, schIsSheetCV, schIsSymEditOk, schIsViewCapEnabled, schIsWireLabel, schLayoutToPinList, schMouseApplyOrFinish, schMove, schNetExprAvailProps, schNetExprEvalNames, schPinListToSchem, schPinListToSchemGen, schPinListToSymbol, schPinListToSymbolGen, schPinListToVerilog, schPinListToView, schPlot, schRegPostCheckTrigger, schRegisterFixedMenu, schRegisterPopUpMenu, schRenumberAllSheet, schRenumberInstances, schRenumberSheet, schReplaceProperty, schSRC, schSaveCurrentPlotOptions, schSchemToPinList, schSelectAllFig, schSelectPoint, schSetAndLoadTsgTemplateType, schSetBundleDisplayMode, schSetCmdOption, schSetEnv, schSetSelectOptions, schSetSymbolOrigin, schSetTextDisplayBBox, schShiftCmdOption, schSingleSelectBox, schSingleSelectPt, schSnapToConn, schSolder, schSubSelectBox, schSymbolToPinList, schSync, schUnregPostCheckTrigger, schUnregisterFixedMenu, schUnregisterPopUpMenu, schUpdateUserSRCErrorAndWarn, schVIC, schVerilogToPinList, schViewToView, schZoomFit, schemeTopLevelEnv, selectProcess, selectResult, set, setContext, setCurrentWindow, setDevCheckOptions, setFnWriteProtect, setPrompts, setShellEnvVar, setSkillPath, setSlotValue, setVarWriteProtect, setarray, setcar, setcdr, setof, setplist, setq, setqbitfield, setqbitfield1, settlingTime, setup, sevAboutTool, sevActivateSelectedAnas, sevActiveSelectedAna, sevAddMenuItemLists, sevAnnotateResults, sevBALMTool, sevBPMTool, sevChangeOutsOnSchematic, sevChooseDesign, sevChooseEnvironmentOptions, sevChooseSimulator, sevChooseTemperature, sevCircuitCond, sevConvergence, sevCopyCellViewVariables, sevCopyVariablesToCellView, sevCornersTool, sevCreateMainWindow, sevDebug, sevDeleteResults, sevDeleteSelectedAnas, sevDeleteSelectedOuts, sevDeleteSelectedVars, sevDeleteSelections, sevDevChecking, sevDirectPlotMenu, sevDisplayViolations, sevEditModels, sevEditOptions, sevEditPlottingOptions, sevEditSelectedAnas, sevEditSelectedOuts, sevEditSelectedVars, sevEditSimulationFile, sevEditStimulus, sevEnvironment, sevEvaluateAndPlotExpressions, sevFindSelectedVars, sevIcon, sevInvokeNCBrowse, sevInvokeSimvision, sevInvokeSimvisionDebugger, sevIsContinuable, sevKmodelTool, sevLMGTool, sevLoadState, sevMenuItems, sevMonteCarloTool, sevMpuTool, sevNetlistAndDebug, sevNetlistAndRun, sevNetlistFile, sevNoAnaSelections, sevNoDesign, sevNoEnvironment, sevNoLog, sevNoOutSelections, sevNoOutputLog, sevNoPlottableExpressions, sevNoPlottableOutputs, sevNoPlottableSignals, sevNoResults, sevNoVarSelections, sevNoViolationsFound, sevNonActiveSelectedAna, sevNonMixedSignal, sevOpenCalculator, sevOpenDRLBrowser, sevOpenEncap, sevOpenJobMonitor, sevOpenPlotWindow, sevOpenPrintWindow, sevOpenSchematic, sevOptimizationTool, sevPCMTool, sevPKGTool, sevParametricTool, sevParasiticsDisplayed, sevParasiticsNotDisplayed, sevPlotAllOutputs, sevPlotSignals, sevPrintResults, sevQuit, sevReset, sevRun, sevRunEngine, sevSaveOceanScript, sevSaveOptions, sevSaveResults, sevSaveState, sevSelectResults, sevSession, sevSetConnectModules, sevSetEngineOptions, sevSetMainWindowPulldownMenus, sevSetMenuItemLists, sevSetPropertyForSelectedOuts, sevSetSchematicPulldownMenus, sevSetSimDataDir, sevSetSolver, sevSetTopSaveDir, sevSetTypeInWindowPulldownMenus, sevSetupSimulationFiles, sevSetupStimuli, sevSimulator, sevStartSession, sevStopEngine, sevTopSaveDir, sevViewCompilerLog, sevViewElabLog, sevViewNcverilogLog, sevViewNetlisterLog, sevViewSimLog, sevViewSimulatorOutput, sevWhatsNew, sh, simAddProbeCapByName, simAddProbeCapByScreen, simAddProbeCapForBusBit, simCheckExist, simCheckHeader, simCheckVariables, simCheckViewConfig, simCleanRun, simDateStamp, simDeleteRunDirFile, simDrain, simEditFileWithName, simExecute, simFindFile, simFlattenWithArgs, simGetLoginName, simGetTermList, simIfNoProcedure, simIlSleep, simInWithArgs, simInitControl, simInitEnv, simInitEnvWithArgs, simInitRaw, simInitRunDir, simInitSimulator, simInstCdsNameExtName, simInstExtNameCdsName, simJobMonitor, simLoadNetlisterFiles, simLoadSimulatorFiles, simNetCdsNameExtName, simNetExtNameCdsName, simNetlistWithArgs, simNoNetlist, simOutWithArgs, simPostNameConvert, simPreNameConvert, simPrintEnvironment, simPrintError, simPrintErrorLine, simPrintMessage, simPrintTermList, simReadNetCapFile, simRep, simRunDirInfile, simRunDirLoad, simRunDirOutfile, simRunNetAndSim, simRunNetAndSimWithArgs, simRunNetAndSimWithCmd, simSetDef, simSetDefWithNoWarn, simStringsToList, simSubProbeCapByName, simSubProbeCapByScreen, simVertToHoriz, simViewFileWithArgs, simWaveOpen, simin, simplifyFilename, simulator, sin, sin, siteDefs, skDisableMessage, skDisableMessageBlock, skEnableMessageBlock, skIgnoreMessage, skTabulateSKILL, skUnignoreMessage, skillDebugger, sklint, slewRate, slotValue, solver, sort, sortcar, sp, spacingRules, specLimits, spectralPower, sprintf, sqrt, sqrt, srandom, srandom, ssb, sstatus, stacktrace, standardViaDefs, status, stddev, step, stepend, stepout, stimulusFile, store, strcat, strcmp, streamLayers, stringToFunction, stringToSymbol, stringToTime, stringp, strlen, strncat, strncmp, sub1, sub1, subclassesOf, subclassp, subst, substring, superclassesOf, suspendJob, sweepNames, sweepValues, sweepVarValues, sxtd, syActivateConstraint, syCompactView, syCreateAlignmentConstraint, syCreateFence, syCreateJogConstraint, syCreateMagneticConstraint, syCreatePathWidthConstraint, syCreateSeparationConstraint, syCreateWireWidthConstraint, syCrossProbe, syDRC, syDeactivateConstraint, syDeleteConstraint, syERC, syExtractView, syFrameView, syGeBackground, syGetAlignmentConstraint, syGetAutoJog, syGetBottomAbutted, syGetDimensions, syGetEnv, syGetFixLowLeftBoundary, syGetForceChildCellCompact, syGetGeoOptions, syGetHardCell, syGetInitialDirection, syGetLeftAbutted, syGetMagneticConstraint, syGetMaxAutoJogs, syGetMaxIterations, syGetPFrame, syGetPeelbackDistance, syGetPostProcess, syGetPreserveCellRow, syGetPreserveWireWidth, syGetRightAbutted, syGetSeparationConstraint, syGetTopAbutted, syGetUseChildBdy, syGetUserConstraint, syGetWireJogConstraint, syGetWireWidthConstraint, syGetXAnchorBoundary, syGetXPinGrid, syGetXPinGridOffset, syGetYAnchorBoundary, syGetYPinGrid, syGetYPinGridOffset, syMenuAlignCreate, syMenuAlignEdit, syMenuAutoJog, syMenuAutoJogByNet, syMenuCompact, syMenuCompareCellView, syMenuConstrGrp, syMenuConvertToGeometric, syMenuCritPath, syMenuCrossProbe, syMenuExplain, syMenuExtract, syMenuFenceCreate, syMenuFenceEdit, syMenuFrame, syMenuHardCell, syMenuJogCreate, syMenuJogEdit, syMenuMagCreate, syMenuMagEdit, syMenuObjEdit, syMenuOverconst, syMenuReInitProbe, syMenuRecompact, syMenuReferenceCell, syMenuReferenceEdge, syMenuSepCreate, syMenuSepEdit, syMenuShowCellSize, syMenuShowFixIO, syMenuShowPG, syMenuShowRelaxed, syMenuSwitchView, syMenuWWidthCreate, syMenuWWidthEdit, syMenuWireMerge, syReadAlignmentConstraint, syReadMagneticConstraint, syReadSeparationConstraint, syReadWireJogConstraint, syReadWireWidthConstraint, syRefCellConstraint, syRefEdgeConstraint, syRemoveProbes, sySetAutoJog, sySetAutoJogByLayer, sySetAutoJogByNet, sySetAutoJogByWire, sySetBottomAbutted, sySetEnv, sySetFixLowLeftBoundary, sySetForceChildCellCompact, sySetGeoOptions, sySetHardCell, sySetInitialDirection, sySetLeftAbutted, sySetMaxAutoJogs, sySetMaxIterations, sySetPFrame, sySetPeelbackDistance, sySetPostProcess, sySetPreserveCellRow, sySetPreservePathWidth, sySetPreserveWireWidth, sySetRightAbutted, sySetTopAbutted, sySetUseChildBdy, sySetXAnchorBoundary, sySetXPinGrid, sySetXPinGridOffset, sySetYAnchorBoundary, sySetYFirst, sySetYPinGrid, sySetYPinGridOffset, sySymToGeo, sySymToPolygon, syUserAllDel, syUserAllOff, syUserAllOn, syUserAllXDel, syUserAllXOff, syUserAllXOn, syUserAllYDel, syUserAllYOff, syUserAllYOn, syWireMergeMovePins, symContactDevice, symDepletionDevice, symEnhancementDevice, symPinDevice, symRectPinDevice, symRules, symWires, symbolToString, symbolp, symckClearDisplay, symckCreateSwitchView, symckDRCSetUp, symckDisplayOptions, symckERCSetUp, symckExplain, symckFitCurrent, symckFitNext, symckFitVisible, symckReset, symckRunDRC, symckRunERC, symckStatus, symckView, symeval, symstrp, system, tableCharacterizationRules, tableSpacingRules, tableToList, tablep, tailp, tan, tan, tangent, tcCreateCDSDeviceClass, tcCreateDeviceClass, tcDeclareDevice, tcDeleteDevice, tcDisplayAttachTechForm, tcDisplayCheckTechForm, tcDisplayCompTechForm, tcDisplayDiscardTechForm, tcDisplayDumpTechForm, tcDisplayLoadTechForm, tcDisplayNewTechForm, tcDisplaySaveTechForm, tcDisplaySetupTechForm, tcDumpTechFile, tcLoadTechFile, tcQcInstallDevices, tcRegPostDumpTrigger, tcRegPostLoadTrigger, tcRegPreDumpTrigger, tcRegPreLoadTrigger, tcUnregPostDumpTrigger, tcUnregPostLoadTrigger, tcUnregPreDumpTrigger, tcUnregPreLoadTrigger, tconc, tddfReadISPFileMC, techBindTechFile, techCloseTechFile, techCopyTechFile, techCreateCustomViaDef, techCreateCustomViaDefByName, techCreateElectricalRuleTable, techCreateLP, techCreateLayer, techCreatePurpose, techCreateRouteSpec, techCreateSiteDef, techCreateSpacingRuleTable, techCreateStdViaDef, techCreateViaSpec, techDeleteDeviceClass, techDeleteLP, techDeleteLayer, techDeletePurpose, techDeleteRouteSpec, techDeleteSiteDef, techDeleteTechFile, techDeleteTechFileName, techDeleteTechLibName, techDeleteTwoLayerProp, techDeleteViaDef, techDeleteViaSpec, techDisplayDFMTechFormCB, techDisplays, techDumpLayerMapping, techFindCustomViaDefByMasterID, techFindRouteSpecByName, techFindSiteDefByName, techFindViaDefByName, techFindViaSpecByName, techGetCellViewSiteDefs, techGetCompactorLayers, techGetCompactorUsage, techGetDBUPerUU, techGetDefaultRouteSpec, techGetDefaultTechName, techGetDeviceCParam, techGetDeviceCellView, techGetDeviceClass, techGetDeviceClassProp, techGetDeviceClassViewList, techGetDeviceFParam, techGetDeviceInClass, techGetDeviceProp, techGetDeviceTechFile, techGetElectricalRule, techGetElectricalRuleTable, techGetElectricalRuleTableEntry, techGetElectricalRuleTables, techGetElectricalRules, techGetEquivLayers, techGetInstDeviceClass, techGetLP, techGetLPAttr, techGetLPPacketName, techGetLayerAbbrev, techGetLayerFunction, techGetLayerFunctions, techGetLayerMaskNumber, techGetLayerMfgResolution, techGetLayerMfgResolutions, techGetLayerName, techGetLayerNum, techGetLayerProp, techGetLayerRoutingGrid, techGetLayerRoutingGrids, techGetLeLswLayers, techGetLxExtractLayers, techGetLxNoOverlapLayers, techGetMPPTemplateByName, techGetMPPTemplateNames, techGetMfgGridResolution, techGetOpenTechFiles, techGetOrderedElectricalRule, techGetOrderedElectricalRules, techGetOrderedSpacingRule, techGetOrderedSpacingRules, techGetOuterViaLayers, techGetParam, techGetParams, techGetPermission, techGetPermissions, techGetPrGenViaParams, techGetPrGenViaParams, techGetPrGenViaProps, techGetPrGenViaRules, techGetPrGenViaRules, techGetPrMastersliceLayers, techGetPrMastersliceLayers, techGetPrMaxStackVias, techGetPrMaxStackVias, techGetPrNonDefaultParams, techGetPrNonDefaultParams, techGetPrNonDefaultProps, techGetPrNonDefaultRules, techGetPrNonDefaultRules, techGetPrOverlapLayers, techGetPrRoutingDirection, techGetPrRoutingDirection, techGetPrRoutingLayers, techGetPrRoutingLayers, techGetPrRoutingOffset, techGetPrRoutingOffset, techGetPrRoutingPitch, techGetPrRoutingPitch, techGetPrStackVias, techGetPrStackVias, techGetPrViaParams, techGetPrViaParams, techGetPrViaProps, techGetPrViaRules, techGetPrViaRules, techGetPrViaType, techGetPrViaType, techGetPrViaTypes, techGetPrViaTypes, techGetPurposeAbbrev, techGetPurposeName, techGetPurposeNum, techGetSpacingRule, techGetSpacingRuleTable, techGetSpacingRuleTableEntry, techGetSpacingRuleTables, techGetSpacingRules, techGetStreamDatatype, techGetStreamLayer, techGetStreamLayers, techGetStreamNumber, techGetStreamTranslate, techGetSymRules, techGetSymWireParams, techGetSymWires, techGetTechFile, techGetTechFileDdId, techGetTechFileName, techGetTechLibName, techGetTimeStamp, techGetTwoLayerProp, techGetUserUnit, techGetViaLayers, techGetViaStackLimit, techGetViaStackLimits, techGetViewTypeUnits, techIsCompactorLayer, techIsDevice, techIsInstQcell, techIsLeLswLayer, techIsLxExtractLayer, techIsLxNoOverlapLayer, techIsPrMastersliceLayer, techIsPrMastersliceLayer, techIsPrRoutingLayer, techIsPrRoutingLayer, techIsPrStackVia, techIsPrStackVia, techIsPrViaDevice, techIsPrViaDevice, techIsReadPermission, techIsViaLayer, techLayerProperties, techLayerPurposePriorities, techLayers, techMakeVirtuosoIccInfo, techManagerOpenDisplayToolBox, techManagerOpenTechToolBox, techManagerToolBox, techOpenDefaultTechFile, techOpenTechFile, techParam, techParams, techPermissions, techPurgeTechFile, techPurposes, techRefreshTechFile, techReopenTechFile, techSaveTechFile, techSetCompactorLayer, techSetCompactorLayers, techSetDBUPerUU, techSetDeviceClassProp, techSetDeviceProp, techSetElectricalRule, techSetElectricalRuleTableEntry, techSetEquivLayer, techSetEquivLayers, techSetEvaluate, techSetLPAttr, techSetLPPacketName, techSetLayerAbbrev, techSetLayerFunction, techSetLayerFunctions, techSetLayerMaskNumber, techSetLayerMfgResolution, techSetLayerMfgResolutions, techSetLayerName, techSetLayerProp, techSetLayerRoutingGrid, techSetLayerRoutingGrids, techSetLeLswLayer, techSetLeLswLayers, techSetLxExtractLayer, techSetLxExtractLayers, techSetLxNoOverlapLayer, techSetLxNoOverlapLayers, techSetMPPTemplate, techSetMfgGridResolution, techSetOrderedElectricalRule, techSetOrderedSpacingRule, techSetParam, techSetPrGenViaProp, techSetPrGenViaRule, techSetPrGenViaRule, techSetPrMastersliceLayer, techSetPrMastersliceLayer, techSetPrMastersliceLayers, techSetPrMastersliceLayers, techSetPrMaxStackVia, techSetPrMaxStackVia, techSetPrMaxStackVias, techSetPrMaxStackVias, techSetPrNonDefaultProp, techSetPrNonDefaultRule, techSetPrNonDefaultRule, techSetPrOverlapLayer, techSetPrOverlapLayers, techSetPrRoutingLayer, techSetPrRoutingLayer, techSetPrRoutingLayers, techSetPrRoutingLayers, techSetPrRoutingOffset, techSetPrRoutingOffset, techSetPrRoutingPitch, techSetPrRoutingPitch, techSetPrStackVia, techSetPrStackVia, techSetPrStackVias, techSetPrStackVias, techSetPrViaProp, techSetPrViaRule, techSetPrViaRule, techSetPrViaType, techSetPrViaType, techSetPrViaTypes, techSetPrViaTypes, techSetPrecision, techSetPurposeAbbrev, techSetPurposeName, techSetReadPermission, techSetReadWritePermission, techSetSpacingRule, techSetSpacingRuleTableEntry, techSetStreamDatatype, techSetStreamLayer, techSetStreamNumber, techSetStreamTranslate, techSetSymRules, techSetSymWire, techSetTechFileName, techSetTechLibName, techSetTimeStamp, techSetTwoLayerProp, techSetUserUnit, techSetViaDefResistancePerCut, techSetViaLayer, techSetViaLayers, techSetViaStackLimit, techSetViaStackLimits, techTruncateTechFile, techUnattachTechFile, techVerifyTechFileId, temp, textFontMap, tfcDeleteCDSDeviceClass, tfcDeleteDevice, tfcDeleteDeviceClass, tfcDeleteDeviceProp, tfcDeleteRuleContactDevice, tfcDeleteSymContactDevice, tfcDeleteSymDepletionDevice, tfcDeleteSymEnhancementDevice, tfcDeleteSymPinDevice, tfcDeleteSymRectPinDevice, thd, theEnvironment, timeToString, timeToTm, times, tmToTime, topEdge, toplevel, tracef, tracep, tracev, tran, transCdlOutDisplay, truncate, tsg, type, unalias, unbreakpt, uncount, undrawEnterFun, uninstallDebugger, unless, unless, unprofile, untrace, untracep, untracev, upperCase, upperRight, userAppTrigger, userMenuTrigger, userPostInstallTriggers, v, value, vcdFile, vcdInfoFile, vceEnableConnectivityAware, vceExtractConnectivity, vcePrepareLibrary, vceVerifyArea, vceVerifyCellView, vceVerifyMarkers, vceVerifyNets, vdb, vecFile, vector, vectorToList, vectorp, vhdlHiImport, vhdlHiInvokeToolBox, vhdlImport, vhdlPinListToVHDL, vhdlRegisterSimulator, vhdlToPinList, vhmsCompilationFailure, vhmsDefaultEdit, vhmsGetCellParameters, vhmsPinListToVHDLAMS, vhmsSaveFile, vhmsSymbolToPinListGen, vhmsToPinList, vhmsUpdateCellCDFParams, vi, viaDefs, viaLayers, viaSpecs, viaStackLimits, vicOpenVlogCallBack, view, viewTypeUnits, vim, vm, vmsUpdateCellViews, vp, vr, vswr, wait, warn, waveVsWave, when, when, where, whereIs, which, while, while, window, windowp, write, writeTable, wtypep, xCoord, xLimit, xcamCameraPS, xcamCameraRas, xcons, xdifference, xmax, xmin, xor, xplus, xquotient, xtimes, xval, yCoord, yLimit, yield, ymax, ymin, zerop, zm, zref, zxtd%
	},
	morekeywords=[1]{% SKILL Keywords from Project
		iasAOCreateForm, iasAOMapToInstance, iasAPXApplyRecursively, iasAPXApplyToCurrentSchematic, iasAPXLog, iasApplyOutputs, iasApplyParametersToCell, iasApplyParametersToInstance, iasBackupIQ, iasDescendInstances, iasGenerateCoefficients, iasHandleReadError, iasHandleWriteError, iasInitADE, iasLogParam, iasMFFTCreateForm, iasMFFTOpenMatlab, iasMatlabFFT, iasOpenNetlistShell, iasOpenVerilogOfCell, iasParseParameters, iasReloadHelpers, iasRotateSpectreLog, iasWriteOutPins %
	},
	sensitive=true,
	morestring=[b]",
	morecomment=[l]{;},
	morecomment=[s]{/*}{*/}
}[keywords,comments,strings] % for SKILL Language
\lstdefinelanguage{VerilogAMS}%
	{
	morekeywords=[0]={ %Verilog XL reserved keywords
		specify, endspecify, follow, invert, unknown, latchhigh, latchlow, pulselow, pulsehigh, ifnone, femtosecs, picosecs, nanosecs, microsecs, millisecs, seconds, picofarads, volts, megahertz, centigrade, package, endpackage, module, macromodule, endmodule, primitive, endprimitive, parameter, input, output, inout, reg, integer, time, real, event, task, endtask, function, endfunction, table, endtable, defparam, when, wire, wand, wor, tri, triand, trior, trireg, tri0, tri1, supply0, supply1, buf, not, and, nand, or, nor, xor, xnor, bufif0, bufif1, notif0, notif1, nmos, rnmos, pmos, rpmos, cmos, rcmos, pulldown, pullup, rtran, rtranif0, rtranif1, tran, tranif0, tranif1, strong0, strong1, pull0, pull1, weak0, weak1, highz0, highz1, small, medium, large, scalared, vectored, signed, assign, deassign, force, release, initial, always, begin, end, fork, join, if, else, case, casez, casex, endcase, default, forever, repeat, while, for, wait, posedge, negedge, edge, disable, specparam, realtime, strength, attribute, endattribute, const, use},%
	morekeywords=[1]{% Verilog AMS reserved keywords
		above, abs, absdelay, abstol, acos, acosh, ac_stim, aliasparam, always, analog, analysis, and, asin, asinh, assign, atan, atan2, atanh, begin, bound_step, branch, buf, bufif0, bufif1, case, casex, casez, ceil, cmos, connectrules, cos, cosh, cross, ddt, ddx, deassign, default, defparam, delay, disable, discipline, discontinuity, driver_update, edge, else, end, endcase, endconnectrules, enddiscipline, endfunction, endmodule, endnature, endparamset, endprimitive, endspecify, endtable, endtask, event, exclude, exp, final_step, flicker_noise, floor, flow, for, force, forever, fork, from, function, generate, genvar, ground, highz0, highz1, hypot, idt, idtmod, if, ifnone, inf, initial, initial_step, inout, input, integer, join, laplace_nd, laplace_np, laplace_zd, laplace_zp, large, last_crossing, limexp, ln, localparam, log, macromodule, max, medium, min, module, nand, nature, negedge, net_resolution, nmos, noise_table, nor, not, notif0, notif1, or, output, parameter, pmos, posedge, potential, pow, primitive, pull0, pull1, pullup, pulldown, pwr, rcmos, real, realtime, reg, release, repeat, rnmos, rpmos, rtran, rtranif0, rtranif1, scalared, sin, sinh, slew, small, specify, specparam, sqrt, strobe, strong0, strong1, supply0, supply1, table, table_model, tan, tanh, task, temperature, time, timer, tran, tranif0, tranif1, transition, tri, tri0, tri1, triand, trior, trireg, vectored, vt, wait, wand, weak0, weak1, while, white_noise, wire, wor, wreal, xnor, xor, zi_nd, zi_np, zi_zd, zi_zp, abstol, access, bound_step, ddt_nature, delay, discontinuity, idt_nature, temperature, abstol, access, continuous, ddt_nature, discrete, domain, idt_nature, units, connect, merged, resolveto, split, wreal, electrical, V, I},%
	morekeywords=[2]{% system tasks and functions
		$bitstoreal,$countdrivers,$display,$fclose,$fdisplay,$fmonitor,$fopen,$fstrobe,$fwrite,$finish,$getpattern,$history,$incsave,$input,$itor,$key,$list,$log,$monitor,$monitoroff,$monitoron,$nokey,$strobe,$display,$random,$abstime,$dist,$dist_uniform,$dist_normal,$dist_exponential,$dist_poisson,$dist_chi_square,$dist_t,$dist_erlang,$stop,$bound_step},%
	morekeywords=[3]{`define, `endmacro, `macro, `optimize_data, `nooptimize_data, `remove_gatenames, `noremove_gatenames, `remove_netnames, `noremove_netnames, `accelerate, `noaccelerate, `default_nettyp, `e, `unconnected_drive, `nounconnected_drive, `pull0, `pull1, `autoexpand_vectornets, `expand_vectornets, `noexpand_vectornets, `automacmods, `noautomacmods, `timescale, `s, `ms, `us, `ns, `ps, `fs, `typdelays, `mindelays, `maxdelays, `celldefine, `endcelldefine, `nocheck_behavior, `check_behavior, `suppress_faults, `nosuppress_faults, `delay_mode_path, `delay_mode_distributed, `delay_mode_unit, `delay_mode_zero, `enable_portfaults, `disable_portfaults, `default_rswitch_strength, `default_switch_strength, `default_trireg_strength, `ifdef, `else, `endif, `include, `default_decay_time, `infinite, `undef, `pre_16a_paths, `end_pre_16a_paths, `switch, `default, `XL, `resistive, `rs_technology, `name, `lowthresh, `highthresh, `resistance, `mapres, `mapcap, `slope, `dynamic_low, `dynamic_high, `n_channel, `rn_channel, `p_channel, `rp_channel, `depletion, `rdepletion, `pullup, `pulldown, `strong1, `strong0, `weak1, `weak0, `large1, `large0, `medium1, `medium0, `small1, `small0, `cox, `deltal, `deltaw, `ldiff, `cgo, `xa, `cdiff, `uselib, `bpi, `nobpi, `default_discipline,`default_transition},%
	alsoletter=\`,%
	sensitive,%
	morecomment=[s]{/*}{*/},%
	morecomment=[l]//,% nonstandard
	morestring=[b]"%
	}[keywords,comments,strings]% 
 % Verilog-A/AMS Language
% define colors for source code list
\ifbool{PrintVersion} %
{% Print Version - Use no colors
	\definecolor{colKeys}{gray}{0}
	\definecolor{colIdentifier}{gray}{0}
	\definecolor{colSystem}{gray}{0}
	\definecolor{colDirectives}{gray}{0}
	\definecolor{colComments}{gray}{0}
	\definecolor{colString}{gray}{0}
	\definecolor{colNumbers}{gray}{0}
	%Latex lisitngs
	\definecolor{latexListBlue}{gray}{0}
	\definecolor{latexListGreen}{gray}{0}
	\definecolor{latexListString}{gray}{0}
}%
{%  Online Version - Use colors and links
	\definecolor{colKeys}{RGB}{0,46,184}
	\definecolor{colIdentifier}{RGB}{0,0,0}
	\definecolor{colSystem}{RGB}{0,46,184}
	\definecolor{colDirectives}{RGB}{0,33,133}
	\definecolor{colComments}{RGB}{0,153,0}
	\definecolor{colString}{RGB}{148,0,209}
	\definecolor{colNumbers}{RGB}{0,0,0}
	%Latex lisitngs
	\definecolor{latexListBlue}{rgb}{0,0,1}	
	\definecolor{latexListGreen}{rgb}{0.133,0.545,0.133}	
	\definecolor{latexListString}{rgb}{0.627,0.126,0.941}			
}

% solve problem with linespacing in list
\newcommand{\lstsetsinglespacing}{%
	\setstretch{1}% no correction afterwards
}

\lstset{%
	language={[LaTeX]TeX},
	abovecaptionskip=\medskipamount,
	numbers=left,
	numberstyle=\tiny\ttfamily\color{colNumbers},
	firstnumber=1,
	frame=single,
	frameround=fttt,
	tabsize=4,
	basicstyle={\small\ttfamily\lstsetsinglespacing},
	numbersep=5mm, 
	breaklines=true,
	prebreak=\raisebox{0ex}[0ex][0ex]{\ensuremath{\hookleftarrow}},
	postbreak=\raisebox{0ex}[0ex][0ex]{\ensuremath{\hookrightarrow\space}}, %,
	frame=single, framexleftmargin=8mm, xleftmargin=8mm,
	frameround=fttt, escapeinside=??,
	otherkeywords={includegraphics,setboolean,setbool,newglossaryentry,newacronym,includestandalone,subcaptionbox},
	keywordstyle=\color{latexListBlue},                    % keywords
	commentstyle=\color{latexListGreen},    % comments
	stringstyle=\color{latexListString},  % strings
	%columns=fullflexible   
	extendedchars=true,  % use literates to support utf8
	%inputencoding=utf8, % set UTF-8 as encoding, not setting this changes nothing
}%
% By default, listings does not support multi-byte encoding for source code. 
% The extendedchar option only works for 8-bits encodings such as latin1.
% To handle UTF-8, you should tell listings how to interpret the special characters by defining them
% The below table will cover most characters in latin languages. 
% For a more detailed explanation of the usage of the literate option check section 6.4 in the Listings Documentation.
\lstset{literate=%
	{á}{{\'a}}1 {é}{{\'e}}1 {í}{{\'i}}1 {ó}{{\'o}}1 {ú}{{\'u}}1%
	{Á}{{\'A}}1 {É}{{\'E}}1 {Í}{{\'I}}1 {Ó}{{\'O}}1 {Ú}{{\'U}}1%
	{à}{{\`a}}1 {è}{{\`e}}1 {ì}{{\`i}}1 {ò}{{\`o}}1 {ù}{{\`u}}1%
	{À}{{\`A}}1 {È}{{\'E}}1 {Ì}{{\`I}}1 {Ò}{{\`O}}1 {Ù}{{\`U}}1%
	{ä}{{\"a}}1 {ë}{{\"e}}1 {ï}{{\"i}}1 {ö}{{\"o}}1 {ü}{{\"u}}1%
	{Ä}{{\"A}}1 {Ë}{{\"E}}1 {Ï}{{\"I}}1 {Ö}{{\"O}}1 {Ü}{{\"U}}1%
	{â}{{\^a}}1 {ê}{{\^e}}1 {î}{{\^i}}1 {ô}{{\^o}}1 {û}{{\^u}}1%
	{Â}{{\^A}}1 {Ê}{{\^E}}1 {Î}{{\^I}}1 {Ô}{{\^O}}1 {Û}{{\^U}}1%
	{œ}{{\oe}}1 {Œ}{{\OE}}1 {æ}{{\ae}}1 {Æ}{{\AE}}1 {ß}{{\ss}}1%
	{ű}{{\H{u}}}1 {Ű}{{\H{U}}}1 {ő}{{\H{o}}}1 {Ő}{{\H{O}}}1%
	{ç}{{\c c}}1 {Ç}{{\c C}}1 {ø}{{\o}}1 {å}{{\r a}}1 {Å}{{\r A}}1%
	{€}{{\EUR}}1 {£}{{\pounds}}1%
	%TODO LISTING add if quotation in Lisitngs does not work when option upquote=true, is used in listing. To use the right quote signs
	%{"}{{"{}}}1 {'}{{'{}}}1 
	%To add if ? or other chars in listings are used.	
	{?}{{?}}1%
}%
%

% ################################################################
% # Configuration of PDF creation and Data in PDF  #
% ################################################################
\usepackage[
pdfpagelayout=TwoPageRight,
bookmarks={true},
bookmarksopen=true,
bookmarksopenlevel=2, 
plainpages=false, 
pdfpagelabels,
]{hyperref}

\ifbool{PrintVersion}
{   % disable colors and hide the links so all text is in the same color.
	\hypersetup{bookmarksnumbered=true , linktoc=all, colorlinks=false, breaklinks=true, final=true, hidelinks = true} 
}
{  
	\hypersetup{bookmarksnumbered=true , linktoc=all, colorlinks=true, linkcolor=blueLink, citecolor=lightGray, urlcolor=darkBlueLink, breaklinks=true, final=true} 
}

\ifbool{IsEnglish}
{
	\author{ {\IASAuthor \ on behalf of the laboratory for Integrated Analog Circuits and RF Systems of the RWTH Aachen University.} }
	\hypersetup{
		pdfauthor= {{{\IASAuthor} on behalf of the laboratory for Integrated Analog Circuits and RF Systems of the RWTH Aachen University.}}, 
		pdftitle={{\ThesisTitleEnglish}}, 
		pdfsubject={{\IASSubject}}
	}
	
}
{
	\author{ {{\IASAuthor} im Auftrag des Lehrstuhl fuer Integrierte Analogschaltungen der RWTH Aachen.} } 
	\hypersetup{
		pdfauthor={{{\IASAuthor} im Auftrag des Lehrstuhl fuer Integrierte Analogschaltungen der RWTH Aachen.}},
		pdftitle={{\ThesisTitleGerman}},
		pdfsubject={{\IASSubject}}
	}
}

%Package bookmark automatically disables hyperref’s bookmarks.  However the package uses some of the code of hyperref . For example, the package redefines the hook \Hy@writebookmark that hyperref inserts in \addcontentsline and other commands. Therefore hyperref’s bookmarks should not be disabled.
%Bookmarks of Hyperref does not work properly so i included Bookmark package
\usepackage{bookmark}
\bookmarksetup{
	numbered, 
	open,
}

% ##########################################################################
% # ACRONYMNS AND GLOSSARIES  # has to be loaded after hyperref
% ##########################################################################
%For the usage of Glossaries with make index put following command in the command line for the make glossaries section. found at Options/ Configure Textstudio/commands/ makeglossaries.
%
% makeindex -s %.ist -t %.glg -o %.gls %.glo 
%
% 1.) Build the Dokument with PDFLatex
% 2.) run make Glossaries
% 3.) run PDFLatex
%
% If you also want Acronyms you need to define a second makeindex call 
%  makeindex -s %.ist -t %.alg -o %.acr %.acn 
%
% For that you have to define a user command at Options/ Configure Textstudio/commands/build
% for that you have to klick show advanced options in the left bottom corner of the Configure Textstudio Window 
%
% For more glossaries you have to add more, look in the Manual 
%
% There is a Profile for TexStudio in the folder which configures the user Macros in Tools/User/makeGlsAcr to build both.
% The "Build & View" Button is defined to run userCommand0 which makes both everytime you compile.
% The "Build & View" Button is also configured to run biber everytime you compile.
% You have to load the profil under Options/Load Profile.

\usepackage[%
automake,% recompile automatically glossaries and acronyms
toc,% put to table of contents
%section=chapter,% 
acronym,% also make acronyms not only Glossary
nonumberlist,% do not list the pages where it is used after the acronyms
nogroupskip,% no gap in list when next acr word first letter changes
nopostdot, % no dot after explanation
style=super% automatically adjust the list in Tabular to get spacing between the first and second column to longest word
%TODO ACRONYMS use this style to adjust the width of the Acronyms Tabular when you are finished with writing
% style = superborder
]{glossaries}
% TODO ACRONYMS set the with of the second row of the tabular of the acronym list to fit you needs
\setlength{\glsdescwidth}{0.8\textwidth}

% Set acronoms font to bold 
\renewcommand{\glsnamefont}[1]{\textbf{#1}} 
% Add the glossaries and acronymns in the file GlossariesAndAcronyms.tex so you can use them.
% Follow the examples to add glossaries and the acronmyns 

\newglossaryentry{gls:latex}
{
	name=Latex,
	description={Is a mark up language specially suited for 
		scientific documents}
}

\newglossaryentry{gls:maths}
{
	name=mathematics,
	description={Mathematics is what mathematicians do}
}

\newglossaryentry{gls:formula}
{
	name=formula,
	description={A mathematical expression}
}


% \begin{acronym}[Text with the largest label]
		
\newacronym{acr:gcd}{GCD}{Greatest Common Divisor}

\newacronym{acr:lcm}{LCM}{Least Common Multiple}

\newacronym{acr:longest}{LAITL}{Longest Acronym in this List}   
\makeglossaries

% #######################################
% # Page Layout  has to be loaded after hyperref
% #######################################
% Geometry sets all needed values to generate the Text Area. It takes the needed values from the dokument class.
%This package provides a flexible and easy interface to page dimensions. You can change the page layout with intuitive parameters. For instance, if you want to set a margin to 2cm from each edge of the paper, you can type just \usepackage[margin=2cm]{geometry}. The page layout can be changed in the middle of the document with \newgeometry command.

\ifbool{PrintVersion} %
{% Print Version - needs Bcor and twoside true
	\KOMAoptions{twoside=true}
	%TODO set binding correction depending on the page count.
	% 1. Binding glue can take up as much as 1mm by itself
	% 2. With a good binding the correction should be no more than half the thickness of all pages, i.e. half the height of the stack of paper you want to bind.
	% 75 Pages are aprox 5mm with our glued binding a value of 1+2.5=3.5 cm or 4cm is a good value
	\KOMAoptions{BCOR=4mm}
	\usepackage[a4paper,
	textwidth = 455.244095pt, % 16cm
	%textheight= 625.960629pt, % 22cm
	includehead,
	includefoot,
	%showframe 
	]{geometry} % 455.24945pt=160.00000mm not cm because of accuracy
}%
{%  Online Version - reading pdf is nicer without bcor and twoside semi
	\KOMAoptions{twoside=true}
	\KOMAoptions{BCOR=0mm}
	\usepackage[a4paper,
	textwidth = 455.244095pt, % 16cm
	%textheight= 625.960629pt, % 22cm
	includehead,
	includefoot,
	centering,
	%showframe 
	]{geometry} % 455.24945pt=160.00000mm not cm because of accuracy
	
	%{%  Online Version - no bcor, sett text in the middle of the page
	\setlength{\oddsidemargin}{-0.4mm} 
	\setlength{\evensidemargin}{-0.4mm}
}

%GADGET Can be used for comments by supervisor or you
% comment out when finisched
% encrease page to have some space for todonotes
%\addtolength{\marginparwidth}{3cm}
%\addtolength{\paperwidth}{4cm}
%%\reversemarginpar
%\usepackage[ textwidth=\marginparwidth, prependcaption]{todonotes}


%TODO COMMENT OUT this Packages in your document
% ##########################################################################
% # for this template only, remove in working document  #
% ##########################################################################
\usepackage{filecontents} % to set file contents in this file.
\usepackage{lipsum} % for dummy text only
\usepackage{printlen} % used to print textwidth and height in cm
\uselengthunit{mm} % set the unit for printlenght


% to debug page layout
%\usepackage{layouts}
% print out the penalty values to log
%\typeout{Current Penalties:^^J
%	\string\clubpenalty=\the\clubpenalty^^J
%	\string\widowpenalty=\the\widowpenalty^^J
%	\string\displaywidowpenalty=\the\displaywidowpenalty}%


\begin{document}
	
% TODO comment in to get Layout information	
% to debug page layout display page information
%	\currentpage
%	\pagediagram
%	\pagevalues
%	
%	{\the\oddsidemargin} = {\printlength\oddsidemargin} \\
%	{\the\evensidemargin} = {\printlength\evensidemargin}
%	\newpage
% End to debug page layout display page information

\pagestyle{empty}

% !TeX spellcheck = en_US
% !TeX encoding = UTF-8
% remove page anker so warning: "pdfTeX warning (ext4): destination with the same identifier (name{page.1}) has been already used, duplicate ignored" does not occur.
\hypersetup{pageanchor=false}

\begin{titlepage}

\pagestyle{myheadings}
% set linespacing to single line.
{\setstretch{1}%
	

\ifbool{IsEnglish}
{
\pdfbookmark[0]{Front Page}{Front Page} %Titelseite in die pdf-Bookmarks aufnehmen        
}
{
\pdfbookmark[0]{Titelseite}{Titelseite} %Titelseite in die pdf-Bookmarks aufnehmen    
}

% set textheight so the date also fits when Tittles get bigger
\enlargethispage{2.3in}
% set topmargin so picture fits in Page
\setlength{\voffset}{-1in}

\ifbool{NoLogo}
{
	
	
\vspace*{1.85cm}
}
{% else logo
	\ifbool{PrintVersion}
	{   
		\begin{flushright}
		 	\includegraphics[trim= 14pt 20pt 15pt 20pt , clip ]{{Logos_2015/rwth_ias_bild_schwarz_grau_cmyk.eps}}
		\end{flushright}
	}
	{
	  \begin{flushright}
	  \includegraphics[trim= 14pt 20pt 15pt 20pt , clip ]{{Logos_2015/rwth_ias_bild_rgb.eps}}
	  \end{flushright}
	}%end
}%end else logo
\setcounter{subfigure}{0}
\ifbool{IsEnglish}
{
	{\small This work was submitted at the Laboratory of Integrated Analog Circuit and RF Systems of the institute for semiconductor technology of the RWTH Aachen University}        
}
{
	{\small Diese Arbeit wurde vorgelegt am Lehrstuhl für Integrierte Analogschaltungen
		der RWTH Aachen.}   
}
\bigskip \\ \bigskip

{\Large \bfseries \ThesisTitleGerman\\} \bigskip 

{\Large \bfseries \ThesisTitleEnglish\\} \bigskip
 
% bring the next text down to the bottom so variation of title lenght don´t change the position of this text
\par\vspace*{\fill}

{\Large \IASSubject\\}\bigskip

{\Large \ifbool{IsEnglish}{by}{von} \\}~{\Large \IASAuthor\\}
{\Large Matr.-Nr. \IASMatNr\\} \bigskip \bigskip

\begin{tabbing}
\ifbool{IsEnglish}
{ \Large Supervised by:}
{ \Large Betreut von:} \= \Large \IASSupervisor \\
\ifbool{secondSupervisor}
{ 	\> \Large \IASSupervisortwo \\   
}{ \> \\ }
 \> \\
 \> \\
\Large 1.~\ifbool{IsEnglish}{Examiner}{Prüfer}: \> \Large Prof.~Dr.~-Ing.~Stefan Heinen \\
\Large 2.~\ifbool{IsEnglish}{Examiner}{Prüfer}: \> \Large Prof.~Dr.~sc.~techn.~Renato Negra \\
\end{tabbing}

\confidential{
\ifbool{IsEnglish}
  { \centering \textbf{CONFIDENTIAL}}
  { \centering \textbf{VERTRAULICH}}
  }{ \smallskip \medskip \bigskip  }

\centering{ \Large \IASNumber}
\begin{flushright}
	 {\Large Aachen, \IASSubmissionDate} \\~\\
\end{flushright}
\newpage

% reset the voffset so the following pages have the Normal Layout
\setlength{\voffset}{0in}

\thispagestyle{empty}
\rule{0pt}{0pt}
\cleardoublepage

}%END set linespacing to single line.
\end{titlepage}

% set page anker so hyperlinks work
\hypersetup{pageanchor=true}
\cleardoublepage

\cleardoublepage
\pagestyle{plain}
% Start frontmatter, with roman Numbering
\frontmatter 
%\pagenumbering{Roman}  % if you want CAPAITAL Roman Letters for Pagenumbering
\pagestyle{empty}
% set textheight so the date also fits when Tittles get bigger
\enlargethispage{3cm}
% set linespacing to single line.
{

\setstretch{1}%
	
\begin{center}
	\textbf{Eidesstattliche Versicherung}
\end{center}
\begin{tabular}{cp{5.2cm}c}
\hspace{5cm}	& & \hspace{4cm} \\ 
\IASAuthor 		& & \IASMatNr \\\cline{1-1} \cline{3-3}
Name   & & Matrikelnummer \\
				& &(freiwillige Angabe) \\
\end{tabular}
\par
Ich versichere hiermit an Eides Statt, \hfill \newline 
dass ich die vorliegende \IASSubject \- mit dem Titel \hfill 
\par
\ThesisTitleGerman 
\par
\ThesisTitleEnglish 
\par
selbständig und ohne unzulässige fremde Hilfe erbracht habe. Ich habe keine anderen als die angegebenen Quellen und Hilfsmittel benutzt. Für den Fall, dass die Arbeit zusätzlich auf einem Datenträger eingereicht wird, erkläre ich, dass die schriftliche und die elektronische Form vollständig übereinstimmen. Die Arbeit hat in gleicher oder ähnlicher Form noch keiner 
Prüfungsbehörde vorgelegen.
\par
% bring the next text down to the bottom so variation of title lenght don´t change the position of this text
\begin{tabular}{cp{5.2cm}c}
	\hspace{5cm}				& & \hspace{4cm} 	\\ 
	Aachen, \IASSubmissionDate 	& &   				\\ \cline{1-1} \cline{3-3} 
	Ort, Datum     				& & Unterschrift 	\\
\end{tabular}
% bring the next text down to the bottom so variation of title lenght don´t change the position of this text
\par 
\vspace*{\fill} 
\par
\textbf{Belehrung:} 

\footnotesize 

\textbf{§~156~StGB: Falsche Versicherung an Eides Statt}

\setstretch{1.2}%
Wer vor einer zur Abnahme einer Versicherung an Eides Statt zuständigen Behörde eine solche Versicherung falsch abgibt oder unter Berufung auf eine solche Versicherung falsch aussagt, wird mit Freiheitsstrafe bis zu drei Jahren oder mit Geldstrafe bestraft.

\textbf{§~161~StGB: Fahrlässiger Falscheid; fahrlässige falsche Versicherung an Eides Statt }

(1) Wenn eine der in den §§~154 bis 156 bezeichneten Handlungen aus Fahrlässigkeit begangen worden ist, so tritt Freiheitsstrafe bis zu einem Jahr oder Geldstrafe ein. \newline
(2) Straflosigkeit tritt ein, wenn der Täter die falsche Angabe rechtzeitig berichtigt. Die Vorschriften des §~158 Abs. 2 und 3 gelten entsprechend.

\normalsize 
Die vorstehende Belehrung habe ich zur Kenntnis genommen:\\
\begin{tabular}{cp{5cm}c}
	\hspace{5cm}				& & \hspace{4cm} 	\\ 
	Aachen, \IASSubmissionDate 	& &   				\\ \cline{1-1} \cline{3-3} 
	Ort, Datum     				& & Unterschrift 	\\
\end{tabular}
}%END set linespacing to single line.
\cleardoublepage
\ifbool{IsEnglish}
{ \chapter{Abstract} 		}
{ \chapter{Inhaltsangabe}	}
In this exquisite work we will show you how to use the \LaTeX \,Template for your thesis at the Institute for Integrated Analog Circuits and RF Systems. Please refer to the chapters which explain the loaded packages and their functions. We tried to give as many examples as possible so you can find anything you need for writing a great thesis.


\ifbool{IsEnglish}
{ \chapter{Acknowledgement} }
{ \chapter{Danksagung}		}
Ich bedanke mich bei den Mitarbeitern des Lehrstuhls für Integrierte Analogschaltungen für deren Erfahrungen verbunden mit der Bereitschaft, diese zu teilen und mir so zu dieser Studienarbeit Unterstützung zukommen zu lassen.
\tableofcontents
\listoffigures
\listoftables
\lstlistoflistings
\printglossaries

%flusch all floatings
\cleardoublepage
% save the Roman Number Page Values to Continue with them later
\newcounter{page_tmp_roman}
\setcounter{page_tmp_roman}{\value{page}}

% Set Header and Chapter font Ecactly as in the text body (i.e., not bold)
%\setkomafont{disposition}{ \normalfont\bfseries }
% \KOMAoption{chapterprefix}{true}
%\renewcommand*\raggedchapter{\centering}
% Make uppercase makes an error \addtokomafont{chapter}{\MakeUppercase}
%\setkomafont{chapterprefix}{\normalsize\mdseries}

\pagestyle{scrheadings}
%\lehead{\textsc{\leftmark}}  % linke aussenseite Kapitel mit numer in Header
%\rohead{\rightmark} % Rechte Aussenseite Section mit nummer in header
%\renewcommand\chaptermarkformat{} % entfern Kapitel nummer in header
%\renewcommand\sectionmarkformat{} % entfernsection nummer in header
%\ofoot{ {\pagemark} } %seitenzahlen an der Aussenseite der Fußzeile 

\mainmatter 
         
\ifbool{IsEnglish}
{ \confidential {\cfoot{\textbf{CONFIDENTIAL}}	}{}}
{ \confidential {\cfoot{\textbf{VERTRAULICH} }	}{}}

%%%%%%%%%%%%%%%%%%%%%%%%%%%%%%%%%%%%%%%%%%%%%%%%%%%%%%%%%%%%%%%%%%%%%%%%%%%%%%%%%
% Here you can add your Chapters												%
%%%%%%%%%%%%%%%%%%%%%%%%%%%%%%%%%%%%%%%%%%%%%%%%%%%%%%%%%%%%%%%%%%%%%%%%%%%%%%%%%
%TODO INCLUDE your Chapters in the file AllChapter than a update of the template can be done by just change the main file.
% !TeX encoding = UTF-8
% !TeX spellcheck = en_US
%TOdo add all you chapter or Sections in a seperated file 
% \include{filename} essentially does a \clearpage before and after \input{filename}, together with some magic to switch to another .aux file.
% It omits the inclusion at all if you have an \includeonly without the filename in the argument. This is primarily useful when you have a big project on a slow computer; 
% changing one of the include targets won't force you to regenerate the outputs of all the rest.

%\includeonly introduction % would only build the selected section

\chapter{Introduction}
Application-specific \cite{latexcompanion}  integrated circuit, or ASIC, is a widely used solution to many electronic application. In the Chair of Integrated Analog Circuits, a variety of ASIC chips are developed. Those chips consist of different functional blocks. Given the architectural design of the chip, the different system blocks are fairly independent of each other, so the design and implementation of each system block can be performed independently by a corresponding designer.

In practice, the functional blocks must be configurable individually, so that the chips are flexible enough to meet the various requirements and work in different situations. The system blocks are configured via predefined signals, which are stored in registers. To configure the chip, we write specific values to corresponding registers via a digital interface such as SPI. The status of the system blocks are also buffered in registers and are accessible via the digital interface.

\begin{figure}[htbp]
\centering
\includegraphics[width = \textwidth]{asic}
\caption{Structure of an ASIC chip\label{fig:Structure of an ASIC Chip}}
\end{figure}

To implement such a digital configuration interface, the designers are required to
\begin{itemize}
\item Define the control and information signals in each system block
\item Define registers and map them to corresponding signals
\item Implement read and write logic for the registers
\item Write documentation of the digital configuration interface
\end{itemize}

The digital configuration interface is implemented in a Hardware Description Language such as VHDL and the documentation is written in LaTeX. Each part of the interface source code and documentation are written by a corresponding system block designer. Finally, there has to be a single designer, let’s say a chip owner, that collects all these VHDL and LaTeX source code and make up a complete SPI interface and documentation.

There are, however, several shortcomings in this manual design approach. First, implementation of the documentation and the VHDL source code can be very tedious and labor intensive. The designers have to write everything completely by hand. However, much of the source code has a strong pattern and can actually be generated by software.

Then, manual implementation of the digital configuration interface can be very error-prone. A good example would be generation of initial values for registers. For each register bit, users have to find out which bit of which signal it is mapped to, and then convert the initial value of the signal from hexadecimal to binary, and get the value of the corresponding signal bit, that is the initial value of this register bit. Some register bits might not be mapped to any signals, then its initial value could be just 0. Finally, concatenate initial value of each register bit and convert it to hexadecimal. The designers have to be very careful not to map a register bit to a wrong signal bit, or write down a wrong value for a register bit.

Also, special attention must be paid to synchronization of the VHDL source code and the LaTeX documentation. For example, when the initial value of a register is updated, the designers have to update the same number everywhere.

Finally, the chip owner has to communicate well with each system block designer. This requires additional efforts, and might cause errors resulted from poor synchronization or communication.

In practice, the chips designed by the Chair can contain decades of system blocks with hundreds of registers, and even more signals. As the complexity of the chips grows, the shortcomings become a serious problem. 

In this background, we propose a software that automates the implementation of the digital configuration interface. The basic idea is simple: the system block designers define the registers and signals and signal-register mappings using the software. They put in documents for each system block, signal or register using a friendly GUI. When the definition is finished, the software can collect all information about the chip, put everything together and generate the VHDL source code and a complete LaTeX documentation automatically.

\begin{figure}[htbp]
\centering
\includegraphics[width = \textwidth]{softwaresolution}
\caption{Proposed Software Solution\label{fig:Proposed Software Solution}}
\end{figure}

Every designer works independently on a single system block, but to generate a complete VHDL source code and documentation, information about all system blocks has to be collected. Due to this fact the software must use a shared data storage which is accessible to all users. We propose to use an open source SQL database system as our data container. With an appropriate database design, we are able to manage the data in a highly secure and efficient way.

The proposed solution has the following great advantages
\begin{itemize}
\item Designers only have to do minimum necessary work, such as defining the name and type of the registers and signals, defining initial values of the signals, mapping signal and register partitions and so on. The rest of work, such as inferring initial values for registers, will be done by the software. By developing a highly friendly user interface, human labor can be reduced to a minimum level. Finally, when definitions of all system blocks are complete, the chip owner do not have to manually collect information. The software shall retrieve all required information from the database and generate the VHDL source code and LaTeX documentation. What the chip owner must do is just some mouse clicks.
\item Synchronization of different outputs is intrinsically ensured. The shared database is the unique data source. Whatever is changed, will be directly reflected on all generations.
\item Errors are less likely to occur, since the most error-prone parts are taken care of by the software. In fact, if all information given by the designers is correct, errors can ideally be eliminated. This requires a highly reliable software, though.
\end{itemize}
\chapter{Software Development Model}
We identify our project as a software development task. To deliver a high quality software product in time, we have to study different development process models, and adopt an appropriate one. The software development model is an abstract process of how software is developed. As software engineering evolves, many software development models have been introduced. Despite of their differences, basically all these models involve the process of requirement analysis, architecture design, implementation, testing etc.

\section{Waterfall Model}
The waterfall model is a preliminary software development model. It breaks down the software development process into linear sequential phases. Each phase only depends on the output of the previous step. The developers have to work step by step, and there is no way to turn back.

\begin{figure}[htbp]
\centering
\includegraphics[width = \textwidth]{WaterfallModel}
\caption{Waterfall Model\label{fig:Waterfall Model}}
\end{figure}

The first step is requirement analysis. The developers use different methods to abstract user requirements, and make the system specifications. Requirement analysis is not as easy as it seems to be. There are many practical issues. In many cases, the developers do not have good access to users, making it difficult to analyze users' requirements. In fact, even the users themselves might not really know what they want. In some cases, due to lack of software engineering practice, users or clients might raise unrealistic requirements. When the user requirements are available, the developers have to make the system specifications, aka system requirements, which is a more detailed description of the software system. Finally, the output is a requirement documentation either in natural language or formulated language.

With the requirement documentation available, the developers have to design the software architecture. The design process is crucial in that the software architecture is difficult and costly to change in the future steps. A bad architecture can lead to poor testability, maintainability and quality. It also makes implementation more difficult.

When the architecture design is complete, we then implement the software. Then, we test the software and verify whether all requirements are satisfied. Finally, we deploy the software and maintain it in everyday use.

The waterfall model breaks down the software development process in a highly structured way. It pays attention to the early stages of the software development process, which tends to be neglected by many developers. According to studies, time spent at the early stages can lead to a great cost reduction in the later stages. A major drawback is inflexibility. The waterfall development process does not provide any way to adapt to changes in requirements. Also, there is no way to look back and make changes on the previous steps.

\section{V-Model}
The V-model is an extension of the waterfall model. It breaks down the development process into the project definition and the test and integration phases. Each development phase has an associated testing phase.

\begin{figure}[htbp]
\centering
\includegraphics[width = \textwidth]{VModel}
\caption{V Model\label{fig:V Model}}
\end{figure}

The V-model places great importance on testing. However, it received similar criticisms as the waterfall model. The process is rigid and unable to well respond to change.

\section{Iterative Model}
To tackle problems with the waterfall model and the V-model and to introduce more flexibility, the iterative model was proposed. The idea is to repeat the whole development process and develop the software iteratively. After each iteration, the software is delivered to the users for evaluation and temporary use. Then, the developer works with the users and update the requirements. After that, the developer repeats the design, implementation and testing phases, and deliver an updated software to the users. Developers repeat this process until the users are satisfied, and then make the final deployable version.

\begin{figure}[htbp]
\centering
\includegraphics[width = \textwidth]{IterativeModel}
\caption{Iterative Model\label{fig:Iterative Model}}
\end{figure}

The iterative model differs to the waterfall model in many important ways. In the waterfall model, requirements, software design and code are fixed once they are finished. Users do not participate in the development process once requirements are fixed. The software is delivered only when the whole development process is complete. Finally, the waterfall model is only suitable for large projects. The iterative model, however, assumes everything is variable in the first place. After every iteration, users are involved in evaluation of the software. The software is developed incrementally, and can be delivered at early iterations. The iterative model is especially suitable for small projects, however, it can also benefit large projects.

\section{Prototype Model}
The waterfall model assumes requirements can be completely abstracted at the early stage of software development, and remain unchanged afterwards. In practice, however, initial requirements might not be complete and they can change over time. By building prototypes, the software developer can review the prototypes with users so as to update the requirements. In this way, users are involved early in the project and provide valuable feedback to the software developers. Prototyping is also a technique to verify whether the designed technical specifications are fulfilled. Based on existing prototypes, current implementation can be improved and prototypes evolve into the stereotype. Additionally, prototyping is also a good way for developers to get an initial idea of the project and estimate timelines of the project.

\begin{figure}[htbp]
\centering
\includegraphics[width = \textwidth]{PrototypeModel}
\caption{Prototype Model\label{fig:Prototype Model}}
\end{figure}

In the field of software engineering, rapid prototyping is especially good for designing the user interface. It is not easy to know what the UI should exactly look like without building one. The developers first write preliminary requirements and a draft of the UI. Then, a prototype is built against the requirements. The prototype is reviewed by the users, and if necessary, the requirements are updated and the prototype is revised. Finally, the requirements are fixed.

Rapid prototyping is not concerned about implementation of complete functionalities, but only the some certain aspects like UI. Also, prototypes are made for review, and are thrown away afterwards. Sometimes rapid prototyping is also referred to as throwaway prototyping.

\section{Conclusion}
We investigated different classical software development models. Recent years, agile development and its variations have become increasingly popular. It emphasizes adaptive planning, face-to-face communication, incremental development and rapid delivery, adaptation to change, and collaboration between the development teams and users etc. It is a brand new development philosophy compared to the classical development models.

In our software development practice, we did not mechanically adopt a certain development model. To better understand the project, make the final version of the requirements and design the UI, we did rapid prototyping. In the following stages, we adopted an iterative and incremental strategy. We set up a weekly meeting to discuss the newly implemented functionality and the current development status. After the first version was complete, we delivered a beta version to the users. According to user feedbacks, we removed errors and made minor changes to the software. And finally, we made a product version and delivered it to the user.

In the whole development process, users actively participated. We responded to user feedbacks rapidly and adapted very well to changes in requirements. A scientific development process ensured that the software product was satisfactory, high quality, and delivered timely. 
\chapter{Requirements Analysis}
In this chapter, we will discuss the process of requirements analysis. Requirements have two aspects: user requirements and system requirements. User requirements reflect what kind of service the software is expected to provide to its users. They can usually be an abstract description of the software. System requirements are more detailed descriptions of the functionalities, services and constraints of the software system. Requirements analysis is usually the very first step of software engineering. It is that important because it will be very difficult to change the requirements in the next stages of software development. Failure in requirements analysis can lead to the consequence that requirements cannot be satisfied or the software cannot be delivered in time.

There are various approaches to requirements analysis \cite{kotonya1998requirements}. Since the requirements for our software are comparably simple and users are well accessible, we decided to adopt the interview approach. The interviewee was the supervisor of the thesis, who would also be a user and admin of the software. The interviews were either in form of emails or face-to-face talks. We had several rounds of interviews. After each interview, we updated the requirements and raised new questions. We did this in an iterative way until we confirmed the requirements were complete. We might want to add more functionalities in the later development stages. But no major modifications would be done on the requirements. We formulated the requirements with natural language.

\section{User Requirements}

To make an easy-to-use software, we have to design a very friendly graphical user interface. Different UI components have to be well organized. They should be simple and intuitive, with underlying logic hidden from the users. One of our goals is to minimize users' operations on the GUI.

The goal of the project is to develop a graphical user interface based software that aids in implementation of the digital configuration interface for AISC chips. It shall allow chip designers to define the system blocks, different types of registers and signals, initial values of signals, and signal-register mappings, and to write different document items of different kinds, including pure text, mathematics equations, tables and images, on different hierarchy levels using a friendly dialog. The software shall be able to display defined system blocks, registers, signals, signal-register mappings, chip designers, register pages and document items etc. in a navigable way. After definition of the chip is complete, the software shall be able to generate working VHDL and LaTeX source code. The user interface shall be carefully designed and hide as much details as possible from the chip designers. To implement the GUI, an appropriate framework should be selected.

The software shall allow chip designers to work simultaneously on different system blocks of the same chip. All inputs given from different users shall be stored in the same well-designed SQL database running on an appropriate database management system.

For data security, user access control and a login mechanism shall be implemented. A basic role model shall be designed, such that standard users are distinguished from the admins in both database or software level and project level permissions. To make the software more reliable, sanity checks must be implemented especially before database operations. A logger shall be implemented, and user activities and database operations shall be logged.

\section{System Requirements}
\subsection{Main Window}
\begin{itemize}
\item The central component of the GUI is a main window.
\item The main window shall contain two components: the navigator and the working area.
\item The navigator shall be a tree view in a CHIP-BLOCK-REGISTER-SIGNAL structure. It shall allow users to search for components of a given pattern and highlight the matching components. The working area shall contain a chip editor on the left and a document editor on the right. The two editors can be switched on or off.
\item The chip editor view shall be a stacked widget containing two pages: the chip level page and the block level page. The chip level page shall display read-only information about the basics of the chip, system blocks, chip designers and register pages. The block level page of the chip editor shall contain two tabs: one for registers and the other for signals. On the signal tab there should be two tables showing signals of the current system block and signal-register mappings of the current signal. The register tab is organized in a similar way.
\item The document editor view shall contain a table showing the list of document items, and a document editing area. The editing area is hidden. When users add a document item or edit an existing one, the editing area shall be displayed.
\item The working area shall respond to changes in the current item in the navigator. When the current item changes, the chip editor shall switch to either the chip level or the block level page. If in the block level page, either the register or signal tab is set as current. Then, the corresponding information shall be displayed on the chip editor. At the same time, the documents regarding this item shall be displayed, if the document editor page is switched on.
\item Unless otherwise noted, below each table in the chip editor and the document editor, there shall be an \textbf{Add} and a \textbf{Remove} button. The users shall be allowed to add a new item, remove or edit an existing item. 
\item At the top of the main window should be a menu bar containing actions on users, chips, views, exports and so on.
\end{itemize}
\subsection{Chips and Documents}
\begin{itemize}
\item The software shall allow users to edit chips and documents in dialogs. Each dialog is designed for a specific use case or functionality, such as adding a new system block. The dialogs must provide appropriate components for users to put in information.
\item The software shall allow users to add new chips or edit exiting chips. It should allow users to put in the chip name, register width, address width, and select whether the chip is MSB first.
\item The software shall allow users to add new system blocks or edit existing ones. It should allow users to put in the block name, block abbreviation, start address in hexadecimal format and select the responsible designer. The start address should fit the chip's address width.
\item The software shall allow project admins to add new chip designers or edit existing ones. It should allow admins to select a designer out of all users in the database, and select the designer's project role.
\item The software shall allow users to add new registers to the current system block or edit existing ones. It should allow users to put in the register name and select a register type out of all register types in the database.
\item The software shall allow users to add new signals to the current system block or edit existing ones. It should allow users to put in the signal name, width, select the signal type, and check whether to add a port for this signal. If the signal is a register signal, i.e. it shall be mapped to registers, the dialog shall allow users to select a register type. When the register type is writable, it shall also allow users to define the initial value in hexadecimal format.
\item The software shall allow users to map signal bits and register bits in an efficient and convenient way. A signal bit can only be mapped to one register bit, and the other way around. 
\item The software shall allow users to add documents to the currently selected item in the navigator or edit existing ones. It shall allow users to select a document type. For a text document, it shall allow users to put in text data. For an image document, it shall allow users to select an image and edit its caption and width. For a table document, it shall allow users to edit the table contents in a table widget and put in its caption. The number of rows and columns of the table shall be adjustable.
\item The software shall allow users to add register pages or edit existing ones. It shall allow users to select an available control signal and specify the page name and a valid page count. It shall allow users to select available registers in a convenient way. The control signal must not be mapped to any registers in the page it controls. A signal can control at most one register page block.
\end{itemize}
\subsection{User Access and Authentication}
\begin{itemize}
\item The software shall allow database admins to create users conveniently.
\item A login mechanism shall be implemented. The main window of the software shall open only when the users provide correct username and password through the login dialog.
\item The software shall allow users to select an existing chip to open. The chips shall be listed in a table containing basic information about them.
\item Each user shall be assigned a database role. The database roles shall be defined in the database, determining whether the users can or cannot add or remove projects or users. 
\item Each project shall have at least one chip designers. Chip designers shall be assigned a project role. The project roles are defined in the database, determining whether the users can add or remove system blocks, chip designers, register pages, registers, signals, signal-register mappings, and documents belonging to a specific chip item (i.e. a system block, a register etc). 
\item The UI shall be adjusted according to the the database role and the project role. For example, some widgets might be disabled if the user is not eligible for it.
\item When the design of a certain chip is complete, the chip owner can freeze the chip. In this case, everything about the chip, except documentation, is frozen, and cannot be edited anymore. In such case, the UI shall disable relevant UI components. It shall allow the chip owner to unfreeze the frozen chip.
\end{itemize}
\subsection{SPI Interface and Documentation generation}
\begin{itemize}
\item The software shall be able to generate a VHDL SPI interface. It shall allow users to select SPI interface templates and configure the output before generation.
\item The software shall be able to generate a LaTeX documentation. It shall allow users to configure the documentation, for example, to select which system blocks to include in the documentation.
\end{itemize}
\subsection{Reliability}
\begin{itemize}
\item Error handling for database operations shall be implemented.
\item All write operations on the database and certain user activities shall be logged for error diagnosis purpose.
\item The software must check if each chip editing or document editing operation, such as adding a new system block or removing a register, is legal. The software shall check e.g. names of the signals/registers/system blocks/users etc, the initial value of a signal, the start address of a system block and so on.
\item Before generating the VHDL source code, the software shall check whether the configuration file contains all necessary predefined variables, whether address blocks of different system blocks overlap or exceed the address width, and whether the signal-register mappings are valid.
\end{itemize}

\section{Conclusion}
Requirement engineering is the very first step of software development. It is of crucial importance, since it is very costly to change the requirements in the following stages. In case of a comparatively small software project, we might be able to define all requirements before the following steps. However, when the project becomes more complex, it will be difficult to write a complete requirement list. There are cases where new requirements have to be added and existing requirements have to be modified.

In our practice, we first made a draft containing the most fundamental requirements. We then made a prototype to help us determine requirements in particular regarding the UI. Other requirements regarding data security and reliability can be comparably easily determined without a prototype.
\chapter{Selection of Infrastructure\label{ch:Selection of Infrastructure}}
In the software engineering practice, it is a common case that we depend on external tools and infrastructure, because we cannot build everything from scratch. To develop a software with a graphical user interface, we have to use a GUI framework. In our project, we also have to use a SQL database management system. They are the infrastructure of our software solution. Before implementation, we have to decide what infrastructure we want to use. To select the best infrastructure a few aspects must be taken into account. 

We should consider the license, which stipulates the conditions on which we could use the software. There are many open source GUI frameworks and database systems. Their license could be tricky, however. Some open source licenses require that the users shall open source their software if the source code of the toolkit is used. This practice is referred to by the community as \textit{copyleft} \cite{copyleft}, since it is against copyright. We want to preserve copyright of our software as much as possible, so we would prefer those toolkits bearing a looser license. 

Another aspect is supported operation systems or platforms. When we select the infrastructure, we would prefer those cross-platform toolkits.

Resources such as learning materials and documentations are also of major importance. Some toolkits have really good official documentations on use cases, sample code and descriptions of the APIs, but others not. These documentations can greatly help us learn the toolkit and use it in our software. Apart from the official resources, the community also plays an important role. It is much easier to get a solution from a larger community when we run into problems.

For the GUI framework, we should consider the supported programming languages. Although programming language is not a high priority, we would prefer to user a modern, friendly language. Some GUI frameworks might offer a specialized designer which makes UI design much easier. The API should also be taken into account. We would definitely prefer a friendly and concise API.

\section{GUI Framework}
Many GUI frameworks are developed and used to build software applications. We investigated and selected the most popular listed Table \ref{tab:Popular GUI Frameworks}.

\begin{table}[htb]
\centering
\caption {Popular GUI Frameworks\label{tab:Popular GUI Frameworks}} 
\begin{tabular}{p{0.15\textwidth}p{0.125\textwidth}p{0.15\textwidth}p{0.125\textwidth}p{0.25\textwidth}}
\hline
Toolkit Name & Cross-platform & Programming Language & Specialized IDE &  License\\
\hline
Qt & Yes & Multiple & Yes & LGPL \\
MFC & No & C++ & No & Proprietary \\
GTK & Yes & Multiple & Yes & LGPL \\
Swing & Yes & Java & Yes & - \\
Tk & Yes & Multiple & No & BSD \\
wxWidgets & Yes & Multiple & No & WxWindows license \\
\hline
\end{tabular}
\end{table}

The first framework we will discuss is Qt. It is also the one we finally chose. It is one of the most popular open source GUI frameworks. The Qt framework is written in C++. However, the Qt company and community developed wrappers for various programming languages including the popular Python. Qt is distributed under the LGPL license. It is a fairly loose license under which we can use Qt for free and do not have to open source our software. Because of this, it is possible to commercialize our software. The Qt company developed a highly customized IDE, the Qt Creator, for developing Qt applications. With Qt Creator, we can design GUI in a what-you-see-is-what-you-get (WYSIWYG) manner.

GTK, Tk and wxWidgets are also very popular cross-platform GUI frameworks. They are all written in C or C++. However, like Qt, they also provide wrappers for other languages. Despite of some differences, their licenses are all pretty loose. However, one shortcoming compared to Qt is they do not have a specialized IDE, although we might use a third-party one. Another shortcoming is they have less documentation. Also, unlike Qt, which is supported by the Qt company and is partially commercial, they are totally maintained by the community. The developer community is also smaller than Qt's.

Microsoft Foundation Class (MFC) is a C++ library for developing desktop applications on Windows operating system. It is essentially an object-oriented wrapper for the Windows API. Due to this fact, it only supports Windows operating system and is not cross-platform. Microsoft offers a community version of Visual Studio, which is recognized by many software engineers as the best IDE for general C++ programming, and it supports MFC very well. Users can design the GUI in a WYSIWYG manner. But MFC is proprietary. Also, it is a little bit out of date. In these considerations, MFC is not a good choice for us.

In the end, we decided to select Qt as our GUI framework. Although PyQt \cite{pyqt}, the Python wrapper, is being very popular nowadays, we still decided to use the C++ API for several reasons. First, Qt Creator does not support PyQt very well. Second, we can build a standalone application in Qt C++. If we use PyQt, however, we have to either depend on Python, or compile it into an executable, which requires additional effort and might bring us more uncertainty or unreliability. Finally, C++ can be much faster than Python, which is beneficial for a GUI application. In fact, the Qt C++ API is very concise and elegant, so we could still develop the software pretty fast. Also, most documentations of Qt are about the original C++ API. Considering all these points, we chose C++ as our programming language.

\section{Database Management System}
The database is crucial for our software. A shared database ensures data consistency and makes cooperation on the same chip project easier. The database we use shall fulfill the following requirements:
\begin{itemize}
\item Accessible to all users.
\item Can be concurrently read or written by multiple users.
\item Data must be reliably stored.
\item Data should be well organized.
\item Data access should be efficient.
\end{itemize}

A relational database is comprised of tables with a certain number of columns. Each column is a data field. It is called relational because a table column can refer to a column in another table. This is one of the core ideas of the relational database. With a relational database, relationship between data can be easily and efficiently handled. The database held on a relational database system is managed with SQL, the Structured Query Language. Thus, such a relational database is often referred to as an SQL database.

The SQL database has many benefits. First of all, it supports highly efficient queries. Usually each table has at least one index field. With the index, the performance of queries can be greatly sped up. Another benefit is concurrency. The database management system is designed for concurrent queries from the very beginning. It can also take care of concurrent writing operations. The database system usually implements an account mechanism, and the data is not directly exposed to users. In this way, data can be securely and reliably stored. Besides that, many database management systems adopt the client/server model. The database can be held remotely on a host machine and users can access it from their local machines. Our software requires a shared database. Although we can put the database on a shared file system, a database server would provide much more flexibility.

\begin{table}[htb]
\centering
\caption {Popular Database Systems\label{tab:Popular Database Systems}} 
\begin{tabular}{p{0.15\textwidth}p{0.125\textwidth}p{0.15\textwidth}p{0.125\textwidth}p{0.3\textwidth}}
\hline
Toolkit Name & Cross-platform & License & Specialized IDE &  Comment\\
\hline
MySQL & Yes & GPL & Yes & Most widely used open \newline source database system \\
Microsoft \newline Access & No & Proprietary & No & Expensive \\
SQLite & Yes & Public \newline domain & No & Limited functionalities \\
Microsoft \newline SQL Server & No & Proprietary & Yes & Expensive \\
Oracle & Yes & Proprietary & Yes & Expensive \\
MariaDB & Yes & GPL & Yes & Compatible to MySQL \\
\hline
\end{tabular}
\end{table}

Among the database systems listed in Table \ref{tab:Popular Database Systems}, Microsoft Access and Microsoft SQL Server are famous systems on Windows operating system. They are both proprietary software. Microsoft SQL Server is of the client/server model while Microsoft Access is not. Oracle is one of the most reputable commercial database management systems. It is proprietary and quite expensive. SQLite, as the name indicates, is a lightweight SQL database system. It provides limited functionalities and is designed for holding small-scale data locally. In fact, Qt itself includes the SQLite database system. 

We are especially interested in MySQL \cite{mysql}. It is the most popular open source database in the world. It is released under the GNU General Public License (GPL) version 2, or the proprietary license. MySQL is used many popular websites including Facebook, Twitter and Youtube. It is a component of the famous LAMP (Linux, Apache, MySQL, PHP) web application software stack. After MySQL's acquisition by the Oracle company, the original developers forked MariaDB \cite{mariadb} from MySQL. It is basically completely compatible to MySQL.

Besides the differences mentioned above, these SQL database systems also have some discrepancy in the SQL language. It is not a problem though. In the end, we chose MySQL or MariaDB in our project. For the software to access the database, we need a MySQL connector which can be downloaded from the official website. It is to be noted that MySQL and MariaDB use the GPL license where copyleft is present. However, the database connector is released under the LGPL license and does not require copyleft. In our software we only need to incorporate the connector, and the database management system is standalone, so we do not have to open source our software.
\chapter{Software Architecture Design\label{ch:Software Architecture Design}}
In this chapter, we will discuss the process of software architecture design. It is a crucial step of software engineering. A poor architecture might often result in low quality, reliability and maintainability. It also makes development more difficult. The architecture is hard to develop in an incremental way, and it is costly to change in the later stages. Therefore, enough attention should be paid to the architecture design.

\section{Layered Software Architecture}
There are different types of software architectures, such as layered architecture, event-driven architecture, micro-kernel architecture, micro-services architecture, cloud architecture etc. \cite{richards2015software}. The layered architecture is the most widely used, and the de facto standard software architecture.

The layered architecture contains roughly four layers: the presentation layer, the business layer, the persistence layer, and the database layer. The presentation layer is where the users interact with the software. It displays internal status of the software and receives user inputs. In case of a GUI software, the presentation layer consists of GUI widgets. The business layer is where a certain business logic is defined and implemented. The persistence layer, also known as data layer, is where the users access the database layer. The database connector API, for example, is on this layer. The bottom layer is the database layer containing SQL databases, files, and so on.

Each layer contains multiple modules that achieve a specific purpose. Within a specific layer, modules are fairly independent of each other. For example, we might want to define a dialog for adding new system blocks and another for adding registers. These two dialogs basically have nothing to do with each other. However, in practice, modules in a certain layer might also need to directly interact with non-neighboring layers. A well-designed layered architecture allows users to easily revise a certain module, or add a new module into a certain layer, without much modification on other modules on the same layer or neighboring layers.

\section{Layered Architecture Design}
With the layered architecture, what we need to do is divide the software into function modules and fill them into those layers. To design a good layered architecture, it is crucial to define good modules. We have to ask ourselves these questions:
\begin{itemize}
\item What modules should be there?
\item What is the boundary of each function module?
\end{itemize}

We found it might be a good idea to define modules by functionality. For each functionality, we create a UI module and a business logic module. We put them in the presentation layer and business layer respectively. The persistence layer and the database layer in our case is quite straightforward. Intuitively, user authentication shall be somewhere between the presentation layer and the business layer. When users do something on the GUI, the authenticator shall check whether the operation is legal, and determine whether to execute the corresponding business logic function or prompt an error message to the user. In fact, we also want the presentation layer to adjust its appearance according to the authentication. As a result, we decided we should modify the default four-layered architecture by adding an authentication layer between the presentation and the business layer.

Following what is discussed, we listed functionalities from the requirements. For many of those functionalities a certain UI component is required. Some functionalities might share a single UI component. For example, when we add or edit a system block, we need the same dialog. Finally, we designed the software architecture as in Figure \ref{fig:Layered Architecture}. Only part of the modules are included in the figure.

\begin{figure}[htb]
\centering
\includegraphics[width = \textwidth]{LayeredArchitecture}
\caption{Layered Architecture\label{fig:Layered Architecture}}
\end{figure}

\section{Object-Oriented System Design}
So far we designed the layered architecture of the software by dividing it into function modules, and assigning each module to a certain layer. The layered architecture, however, does not model the interactions between modules within a layer or across layers. Also, the models are not associated with any implementation form. In practice, they can be implemented with a class or a function. In this section, we will design the software system in more details following the object-oriented programming paradigm, and finally produce a class diagram. However, at this point we do not consider module design and implementation details. We will leave this to Chapter \ref{ch:Module Design and Implementation}.

\subsection{Object-Oriented Programming}
Object-oriented programming (OOP) is a programming paradigm based on objects. An object contains data and methods. Objects interact with each other and this usually leads to modification of the data in the objects. In the class-based OOP languages, an object is an instance of a class, a template for creating objects and providing initial values and behaviors. 

OOP has three central ideas \cite{eckel2003java}
\begin{itemize}
\item Encapsulation: the classes hide their internal status, and only expose necessary information.
\item Inheritance: it represents the \textit{is-a-type-of} relationship between classes.
\item Polymorphism: it provides a single interface to entities of different types.
\end{itemize}

The definition of classes and their relationships can be represented by the UML class diagram. The class is represented with a rectangle with three boxes as in Figure \ref{fig:Example Class}.
\begin{figure}[htb]
\centering
\includegraphics[width = 0.3\textwidth]{ExampleClass}
\caption{Example Class\label{fig:Example Class}}
\end{figure}
The top box contains the class name. The members variables (i.e. the data) are in the middle box. In the bottom box are member methods. The member variables and methods are decorated with +, \# or - representing public, protected and private respectively. Classes related to others are connected with lines of different forms as in Figure \ref{fig:Class Relationships}. Relationships can be categorized into instance-level and class-level. According to the degree of dependency, instance-level relationships are \textit{dependency}, \textit{association}, \textit{aggregation} and \textit{composition}. Dependency relationship exists when an instance references another. Association represents a \textit{has-a} relationship. Aggregation and composition are variants of the \textit{has-a} association relationship in that they represent a \textit{is-part-of} relationship. The difference is, the composition relationship entails ownership. If class A has a composition relationship with B, when B is destroyed A will be destroyed too. The class-level relationships reflect the inheritance relationships between classes. The two relationships differ in that realization involves an abstract superclass, aka an interface. 

\begin{figure}[htb]
\centering
\includegraphics[width = 0.8\textwidth]{ClassRelationships}
\caption{Class Relationships\label{fig:Class Relationships}}
\end{figure}

\subsection{System Design}
With the OOP background, our task is to design the system and produce a class diagram.  At this point, we have to answer two questions
\begin{itemize}
\item What classes should there be and what are their functionalities?
\item What are their relationships with each other?
\end{itemize}
These two questions are not independent. The answer to one can influence the other.

In Qt, a general design pattern is to first design a UI form using the Qt Creator, and then design the class incorporating the UI as one of its member variable. We can design the UI in a WYSIWYG manner. The Qt Creator will then generate a UI form in XML format containing all information about the UI and a C++ class from the UI form. Then, we will create a class with an instance of the UI class. We can then write business logic functions in this class. In this way, modules in the presentation layer and business logic layers are unified. 

Some functionalities require user inputs. For example, to add a system block the name, abbreviation and start address of the system block must be given by the user. The client requires developers to design an appropriate dialog for such use cases. However, the dialog can be reused for different functionalities. The system dialog can be used for not only adding but also editing system blocks, for example. In such cases, a single UI class can be associated with multiple functionalities. Following the same pattern, we can design dialog classes for other functionalities such as adding or editing a signal, register and so on. Since these functionalities are independent of each other, these classes basically do not have relationships with each other, and we can design and implement them in an incremental manner.

According to the requirements, the main window shall contain a chip navigator, a chip editor view and a document editor view. Intuitively, we can design a \textbf{ChipNavigator}, \textbf{ChipEditorView} and \textbf{DocumentEditorView} class respectively. The \textbf{ChipEditorView} displays information about the chip and provides access to those functional dialogs related to chip editing. Similarly, the \textbf{DocumentEditorView} displays document items and provides access to document editing dialogs. The \textbf{MainWindow} itself is also associated with global level functional dialogs such as \textbf{SPIGenerationDialog}, \textbf{DocumentationGenerationDialog}, \textbf{ChangePasswordDialog} and so on. In the end, the system can be illustrated with the class diagram as shown in Figre \ref{fig:Register Manager Class Diagram}.

\begin{figure}[htb]
\centering
\includegraphics[width = \textwidth]{RegisterManagerClassDiagram}
\caption{Register Manager Class Diagram\label{fig:Register Manager Class Diagram}}
\end{figure}

Note that this class diagram is simplified in many ways. First, it only shows classes we want to create. The classes provided by Qt are omitted. The UI classes are not in the diagram either. Instead, they appear as a member named \textit{ui} of many classes in the diagram. Since we were still in the design phase so the classes were incomplete. Only the most important member variables and methods are shown in the diagram. Also, most classes have a dependency relationship with the \textbf{DatabaseHandler} class. However, their relationships are omitted in the diagram so as to make the diagram more readable and understandable.

The center of the software is the \textbf{MainWindow}. From the diagram we see that the \textbf{MainWindow} \textit{has} a \textbf{ChipNavigator}, \textbf{ChipEditorView}, \textbf{DocumentEditor}, \textbf{Authenticator} and a \textbf{LoginDialog}. Unlike other dialog classes such as \textbf{CreateUserDialog}, the \textbf{LoginDialog} is a member of the class \textbf{MainWindow} and it exits in the whole lifetime of the software (which is exactly the lifetime of the \textbf{MainWindow}). This is because we consider that users may want to log out and log in with another account. The \textbf{MainWindow} is directly dependent of many dialogs which are not directly associated to the \textbf{DocumentEditorView} or the \textbf{ChipEditorView}. These dialogs are not members of the class \textbf{MainWindow}. When they are needed, a certain action is triggered by the user and a temporary instance will be created.

The \textbf{ChipNavigator}, \textbf{ChipEditorView} and \textbf{DocumentEditorView} are members of the \textbf{MainWindow} and they exist in the whole life cycle of the software. The \textbf{ChipEditorView} are dependent of those chip editor dialogs including \textbf{EditChipDialog}, \textbf{EditSystemBlockDialog}, \textbf{EditRegisterDialog}. The \textbf{DocumentEditorView} is dependent of the \textbf{DocumentGenerator}. Likewise, these dialogs are created only when needed.

A special case is the \textbf{Authenticator}. It is the major component of the authentication layer. It stores the database permissions and project permissions of the current user to the current project. Whenever the current status changes, for example, when the user clicks on another system block on the navigator, permissions are requested and the UI is updated by enabling or disabling certain widgets. The \textbf{MainWindow} \textit{has} a member \textbf{Authenticator}. The pointer to the member \textbf{Authenticator} is passed to the \textbf{ChipNavigator}, \textbf{ChipEditorView} and \textbf{DocumentEditorView}, such that they actually access the same \textbf{Authenticator} instance.

Dependency relationships between the \textbf{MainWindow}, \textbf{ChipEditorView} etc. and the functional dialogs are easy and clear because those functional dialogs are basically independent of each other. However, the relationships between the \textbf{ChipNavigator} and the \textbf{ChipEditorView} or \textbf{DocumentEditorView} are tricky. In our design, they are not directly related to each other, but via the \textbf{MainWindow}. Whenever the user clicks on something in the navigator, it will \textit{tell} the \textbf{MainWindow} that something is triggered by sending it a \textbf{signal} (will be discussed in the Chapter \ref{ch:Introduction to Qt and SQL Programming}). The \textbf{MainWindow} then updates the current system block, register, signal etc., and updates the status of the \textbf{ChipEditorView} and \textbf{DocumentEditorView} by function calls. Likewise, when the chip is edited, the \textbf{ChipEditorView} sends a signal to the \textbf{MainWindow}. The \textbf{MainWindow} then calls a corresponding function to update the navigator.

To generate the SPI interface and the documentation, we use the \textbf{SPIGenerationDialog} and the \textbf{DocumentGenerationDialog}. With these dialogs, we can specify necessary inputs and configure the export. The \textbf{DocumentGenerationDialog} then creates a \textbf{DocumentGenerator} and calls the corresponding methods to create documentation rather in LaTeX or HTML (see Section \ref{sec:Document Generation}) format. Likewise, to generate the SPI interface, the \textbf{SPIGenerationDialog} creates a certain type of generator given the format of the export.

The \textbf{NamingTemplate} is designed in response to the requirement that the names of registers and signals shall be formatted. For example, we may want all register names to start with the block abbreviation and end with REG. To display the current naming templates, we use the \textbf{NamingTemplateDialog}, which also provides entries to editing the naming templates using the \textbf{EditNamingTemplateDialog}.

Although the class diagram does not regard how each class is implemented but only what functionalities they provide, we cannot design a software system completely without considering the design and implementation details of each class. We will discuss this in the Chapter \ref{ch:Module Design and Implementation}.
\chapter{Introduction to Qt and SQL Programming\label{ch:Introduction to Qt and SQL Programming}}
So far we have designed the detailed software architecture. Before starting module design and implementation we have to introduce Qt GUI programming and SQL programming.

\section{Qt GUI Programming}
Qt \cite{qt} identifies itself as a cross-platform software development framework. It was originally designed for making graphical user interface applications. The Qt framework consists of the following main modules in Table \ref{tab:Qt Modules} \cite{qtmodules}. As the name indicates, Qt Core is the core of the Qt framework. It defines QObject, the base class of all Qt objects. QObject provides lower-level infrastructure for the Qt Framework and applications.

\begin{table}[htb]
\centering
\caption {Qt Modules\label{tab:Qt Modules}}
\begin{tabular}{p{0.2\textwidth}p{0.75\textwidth}}
\hline
Module & Description \\
\hline
Qt Core & Core non-graphical classes used by other modules. \\ 
Qt GUI & Base classes for graphical user interface (GUI) components. Includes OpenGL. \\
Qt Multimedia & Classes for audio, video, radio and camera functionality. \\
Qt Multimedia \newline Widgets & Widget-based classes for implementing multimedia functionality. \\
Qt Network & Classes to make network programming easier and more portable. \\
Qt QML & Classes for QML and JavaScript languages. \\
Qt Quick & A declarative framework for building highly dynamic applications with custom user interfaces. \\
Qt Quick Controls & Provides lightweight QML types for creating performant user interfaces for desktop, embedded, and mobile devices. These types employ a simple styling architecture and are very efficient. \\
Qt Quick Dialogs & Types for creating and interacting with system dialogs from a Qt Quick application. \\
Qt Quick Layouts & Layouts are items that are used to arrange Qt Quick 2 based items in the user interface. \\
Qt Quick Test & A unit test framework for QML applications, where the test cases are written as JavaScript functions. \\
Qt SQL & Classes for database integration using SQL. \\
Qt Test & Classes for unit testing Qt applications and libraries. \\
Qt Widgets & Classes to extend Qt GUI with C++ widgets. \\
\hline
\end{tabular}
\end{table}

\begin{figure}[htb]
\centering
\includegraphics[width = \textwidth]{SignalAndSlot}
\caption{Signal and Slot\label{fig:Signal and Slot}}
\end{figure}

Special attention should be paid to the \textbf{Signal and Slot} \cite{signalsandslots} mechanism supported by QObject. It is a powerful, friendly and intuitive alternative to the callback technique, and is a central feature of Qt. When an event occurs, a signal is emitted. The corresponding slots receive that signal and respond by executing a function. The signal and slot mechanism has three elements: the signal, the slot, and connection between them. They can be illustrated with Figure \ref{fig:Signal and Slot}.

Qt has predefined many useful signals for the UI widgets, and they are automatically emitted once a corresponding event occurs. For example, when we click on a push button, it will send a \textbf{click()} signal. To customize the behavior of an event, we can define slot functions and connect them to the corresponding signal. We can also subclass the standard Qt widgets and define our own signal functions.

In a widget based application, Qt GUI and Qt Widgets must be included. Qt GUI classes provide infrastructure for user interfaces. They are extended to Qt Widgets, which are primary elements for creating GUI. Widgets can display data and status information, receive user inputs, and provide a container for other widgets that should be grouped together. Qt Widgets can be categorized into the following groups
\begin{itemize}
\item Layouts are geometry managers. They automatically arrange child widgets in their containers according to their \textbf{sizeHint} and \textbf{sizePolicy} properties. According to the direction in which the child widgets are distributed, there are vertical, horizontal, grid and form layout.
\item Spacers provide blank space in a layout. There are vertical and horizontal spacers.
\item Buttons provide push button or checkable button widgets, among which the most frequently used are push buttons, check boxes and radio buttons. A special case is the button box, which contains a configurable set of predefined push buttons.
\item Item widgets/item views provide architecture to manage the relationship between data and the way it is presented to the user. There are list, tree, table view/widget etc.
\item Containers provide an area to group widgets together. There are group box, tool box, tab widget, stacked widget etc.
\item Input widgets provide a friendly way to take different types of user inputs. There are combo boxes, line edits, text edits, spin boxes etc. These text widgets can also be configured as readonly and used to display text data. Qt also provides widgets to take date and/or time.
\item Display widgets provide a friendly way to display different types of data. There are labels, text browsers, calendar widgets, progress bars and others.
\end{itemize}

Most of these basic widgets are configurable according to our requirements. We can also customize the widgets by subclassing them. With these widgets, we are able to design a powerful and friendly user interface.

The Qt Company develops an IDE called Qt Creator. It is different from any other IDEs in that it is highly customized for Qt. It provides a very friendly and convenient wizard for creating Qt applications and what it calls designer form classes, which are essentially subclasses of Qt widgets such as \textbf{QDialog}, \textbf{QMainWindow} etc. or the generic \textbf{QWidget}, depending on what kind of UI we want to create. The user interface is generated from the XML UI form and is a member variable of the designer form class. The definition a designer form class is similar to Code \ref{lst:Header of an Example Designer Form Class} and \ref{lst:Source Code  of an Example Designer Form Class}.

\begin{lstlisting}[language=C++, caption={Header of an Example Designer Form Class\label{lst:Header of an Example Designer Form Class}}]
// mainwindow.h
#include <QMainWindow>

namespace Ui {
class MainWindow;
}

class MainWindow : public QMainWindow
{
    Q_OBJECT
public:
    explicit MainWindow(QWidget *parent = nullptr);
    ~MainWindow();
private:
    Ui::MainWindow *ui; // instance of the UI class.
};
\end{lstlisting}

\begin{lstlisting}[language=C++, caption={Source Code of an Example Designer Form Class\label{lst:Source Code  of an Example Designer Form Class}}]
// mainwindow.cpp  
#include "mainwindow.h"
// header generated by Qt Creator.
// The UI class is defined here
#include "ui_mainwindow.h"  

MainWindow::MainWindow(QWidget *parent) :
    QMainWindow(parent),
    ui(new Ui::MainWindow)
{
    ui->setupUi(this);
}

MainWindow::~MainWindow()
{
    delete ui;
}
\end{lstlisting}

\section{SQL Programming}
In Chapter \ref{ch:Selection of Infrastructure}, we have already introduced the relational database management systems and had some basic knowledge about them. We selected MySQL as our database system. Before starting implementation, we have to learn about SQL operations, the MySQL connector and the designed database structure as well.

\subsection{SQL Operations}
SQL operations can be categorized into \textit{database} operations and \textit{data} operations. Database operations allow users to create, remove, update or query databases and tables. To be concluded, they allow users to define a database and manage its data structure. Data operations allow users to insert, remove, update or query data entries in the tables. To develop a database application, we usually predefine a database and its data structure. The database applications only do data operations and do not modify the data structure.  

\textbf{1. Database Operations\\}
To create a database, we first create an empty database in the database management system. After that, we create tables in which data is contained in a structured way. When a table is created, we have to specify the data columns the table should contain and their corresponding data types. Code \ref{lst:Creation of Database and Tables} shows an example of how databases and tables are created.

\begin{lstlisting}[language=SQL, caption={Creation of Database and Tables\label{lst:Creation of Database and Tables}}]
CREATE DATABASE ias;
USE DATABASE ias; -- use the database ias

-- create table staff
CREATE TABLE IF NOT EXISTS staff (
   -- data columns
   staff_id INT AUTO_INCREMENT,
   name VARCHAR(40) NOT NULL,
   phone_number VARCHAR(40) NOT NULL,
   address VARCHAR(100),
   PRIMARY KEY ( user_id )
) ENGINE=InnoDB DEFAULT CHARSET=utf8;

-- create table project
CREATE TABLE IF NOT EXISTS project (
   -- data columns
   project_id INT AUTO_INCREMENT,
   project_name VARCHAR(100) NOT NULL,
   responsible_person_id INT NOT NULL,
   PRIMARY KEY ( project_id )
) ENGINE=InnoDB DEFAULT CHARSET=utf8;
\end{lstlisting}

Special notes about the code: \\
\textbf{AUTO\_INCREMENT}: to make the integer increment by one automatically. Usually applied to primary keys. \\
\textbf{NOT NULL}: to force the data field not to be null. \\
\textbf{PRIMARY KEY}: to set the data field to be the primary key for indexing data rows. \\
\textbf{ENGINE}: to set the database engine. \\
\textbf{CHARSET}: to set character set for this table.

We can write a complete definition of the tables upon creation. However, we are also able to alter a table after it is created. With the \textbf{ALTER TABLE} command, we are able to add a new column, drop an existing column, change a column's name and its data type, change the table name, add keys and so on. See Code \ref{lst:Alter the Table}.

\begin{lstlisting}[language=SQL, caption={Alter the Table\label{lst:Alter the Table}}]
ALTER TABLE staff ADD birthday DATE [AFTER address]; -- to add a column birthday [after address]
ALTER TABLE staff DROP birthday; -- to drop the column birthday
ALTER TABLE staff CHANGE name staff_name VARCHAR(40) NOT NULL; -- to change the column name and its data type
ALTER TABLE project RENAME TO asic_project; -- to rename the table project
\end{lstlisting}

In many cases we might want to add some constraints to tables by adding keys to them. The \textbf{UNIQUE} key constraint ensure values of a certain table column are unique. Inserting a value which already exists in the table will fail. We can also add a \textbf{UNIQUE} key to multiple columns. In this case, the composition of values corresponding to the columns must be unique. A \textbf{PRIMARY} key is similar to a \textbf{UNIQUE} key which is \textbf{NOT NULL}. Each table can have at most one \textbf{PRIMARY} key. It is a good practice to always define an ID column for each table and add a \textbf{PRIMARY} key to it. A \textbf{Foreign} key is a column in table 1 whose values match the \textbf{PRIMARY} key in table 2. Values of this column in table 1 must exist in the reference column in table 2.

We still take Code \ref{lst:Creation of Database and Tables} for example. In this example, staff names must not duplicate. We can then add a \textbf{UNIQUE} key constraint to the column \textbf{name} of the table \textbf{staff}. Similarly, we might also want to add a \textbf{UNIQUE} key to the column \textbf{project\_name} of the table \textbf{project}. When we create a new project, we have to specify the responsible person by his or her staff ID. This ID must already exist in the \textbf{staff} table.
\begin{lstlisting}[language=SQL, caption={Add Key Constraints to Tables\label{lst:Add Key Constraints to Tables}}]
ALTER TABLE staff ADD UNIQUE (name);
ALTER TABLE project ADD UNIQUE (project_name);
-- ALTER TABLE staff ADD PRIMARY KEY (staff_id);
-- ALTER TABLE project ADD PRIMARY KEY (project_id);
ALTER TABLE project ADD FOREIGN KEY (responsible_person_id) REFERENCES staff (staff_id);
\end{lstlisting}
Thus, we can add a \textbf{Foreign} key to the column \textbf{responsible\_person\_id} of the table \textbf{project} referencing the column \textbf{staff\_id} of the table \textbf{staff}. The \textbf{PRIMARY} keys have been created with the tables. These keys can also be added afterwards using the \textbf{ALTER TABLE} command. See Code \ref{lst:Add Key Constraints to Tables}.

\textbf{2. Data Operations\\}
Data operations allow users to query data, insert new data entries to the database or update existing data entries. To insert a new data entry, we specify values for each data column and insert the data into a certain table using the \textbf{INSERT} command. To update a data entry, we have to specify the data columns to update and their new values. We usually have to specify which data entries to update using a \textbf{WHERE} clause. Similarly, when we delete data entries, we also need to specify a \textbf{WHERE} clause otherwise all data entries in that table will be deleted. In practice, every table usually has a ID column, which is an \textbf{AUTO\_INCREMENT} \textbf{PRIMARY} key. If we know the value, the \textbf{WHERE} clause can be \textbf{WHERE id\_column=1}. See Code \ref{lst:Insert, Update and Delete Data Entries}.

\begin{lstlisting}[language=SQL, caption={Insert, Update and Delete Data Entries\label{lst:Insert, Update and Delete Data Entries}}]
INSERT INTO staff ( staff_id, name, phone_number, address ) VALUES ( 1, "Lei", "NaN", "Sichuan" );
INSERT INTO staff ( staff_id, name, phone_number, address ) VALUES ( 2, "Bastl", "NaN", "Munich" );
INSERT INTO project ( project_id, project_name, responsible_person_id ) VALUES ( 1, "IASRegisterManager", 1 ); -- responsible person is Lei
INSERT INTO project ( project_id, project_name, responsible_person_id ) VALUES ( 2, "Best Project", 2 ); -- responsible person is Bastl
INSERT INTO project ( project_id, project_name, responsible_person_id ) VALUES ( 3, "Bad Project", 1 );   
UPDATE staff SET address="Aachen", phone_number="Unknown" WHERE staff_id=1; -- Lei's address and phone_number change
DELETE FROM project WHERE project_id=3; -- "Bad Project" is deleted
\end{lstlisting}

To query data we use the \textbf{SELECT} command. We have to specify the data columns to select, or just use \textbf{*} as a shortcut to all columns. Like the \textbf{UPDATE} and \textbf{DELETE} command we might want to specify a \textbf{WHERE} clause unless we want to query all data entries in that table. Besides data columns, we can also \textbf{SELECT} a function. The data value would be the result of that function. In many cases we might want to select data from multiple related tables. In this case, the \textbf{INNER JOIN} command might help. See Code \ref{lst:Data Query}.

\begin{lstlisting}[language=SQL, caption={Data Query\label{lst:Data Query}}]
SELECT address FROM staff WHERE name="Bastl"; -- result: Munich
SELECT name, address FROM staff WHERE staff_id=1; -- result: Lei, Aachen
SELECT COUNT(*) FROM staff; -- result: 2 
SELECT staff.name, project.project_name FROM staff INNER JOIN project WHERE staff.staff_id = project.responsible_person_id;
-- result: 
-- Lei, IASRegisterManager
-- Bastl, Best Project
\end{lstlisting}

\subsection{Database Connector and Database Handler}
MySQL uses the client/server architecture. Although we are able to connect to the server either by GUI tools such as MySQL Workbench or using the command line, to build a database application we have to use the MySQL Connector \cite{mysqlconn}. MySQL Connector is not a standard component of the MySQL package. It can be downloaded from MySQL website and installed easily.

To use the connector, we need to include the header files in our code and initialize the database connection. First, we should create a driver instance, and connect it to the database server given the host name, username and password. Then, we will create a statement instance. To execute a query statement, call the \textbf{executeQuery()} function, and parse the result set. To execute any statement that writes the database, call the \textbf{executeUpdate()} function.

We realized that the raw APIs offered by the MySQL Connector is not so friendly and reusable. One problem is that we have to write a complete SQL statement for each query or update. This is not really necessary, because the statement is highly structured. We can write functions that accept only keywords, such as table name, column name, and generate complete statements. In this way, the API would be much more friendly to users. Another problem is that we have to manually parse the results for each query. Also, we can write a function to do this. Besides these, from an object-oriented programming point of view, it would be nice to wrap all these stuffs, the driver, the connection, statement and so on, in a class. So we designed a database handler class as in Code \ref{lst:Header of the Database Handler} and \ref{lst:Source Code of the Database Handler}. In the code we omitted the database operation functions such as creating a table, querying data and so on.

\begin{lstlisting}[language=C++, caption={Header of the Database Handler\label{lst:Header of the Database Handler}}]
// database_handler.h
class DataBaseHandler
{
public:
    static bool initialize(const QString& hostname, const QString& database, const QString& username, const QString& password);
    static bool use_database(const QString &database);
    static void close();
    static void commit();
    static void rollback();

    /*
     We define database operation functions here.
    */

    static bool get_next_auto_increment_id(const QString& tablename, 
                                     const QString& id_field, 
                                     QString& id);
    static QString get_error_message();

private:
    static bool execute(const QString &statement);
    static bool execute_query(const QString &statement);

    static QString error_message_;
    static QString database_;
    static sql::Driver* driver_;
    static std::unique_ptr< sql::Connection > con_;
    static std::unique_ptr< sql::Statement > stmt_;
    static std::unique_ptr< sql::ResultSet > res_;
};
\end{lstlisting}

\begin{lstlisting}[language=C++, caption={Source Code of the Database Handler\label{lst:Source Code of the Database Handler}}]
// database_handler.cpp
bool DataBaseHandler::initialize(const QString &hostname, const QString &database, const QString& username, const QString& password)
{
    try {
        con_.reset(driver_->connect(hostname.toUtf8().constData(),
            username.toUtf8().constData(),
            password.toUtf8().constData()));
        con_->setAutoCommit(false);
        if (database != "")
        {
            con_->setSchema(database.toUtf8().constData());
            database_ = database;
        }
        stmt_.reset(con_->createStatement());
        return true;
    } catch (sql::SQLException &e) {
        error_message_ = e.what();
        return false;
    }
}

bool DataBaseHandler::execute(const QString &statement)
{
    try {
        stmt_->executeUpdate(statement.toUtf8().constData());
        return true;
    } catch (sql::SQLException &e) {
        error_message_ = e.what();
      return false;
    }
}

bool DataBaseHandler::execute_query(const QString &statement)
{
    try {
        res_.reset(stmt_->executeQuery(statement.toUtf8().constData()));
        return true;
    } catch (sql::SQLException &e) {
        error_message_ = e.what();
      return false;
    }
}
\end{lstlisting}

This initialization of the database connection is wrapped in the \textbf{initialize()} function. After initialization, we can do database operations. Every database operation function parses its parameters and generate a SQL statement string, and either calls the \textbf{execute()} or the \textbf{execute\_query()} function, depending on whether it is an update or a query operation. The return type of the database operation functions should always be bool so that we know whether they are successful. In practice we want to execute a series of SQL operations as a work unit which we call a transaction. One property of the transaction is atomicity. It requires that all SQL operations must be successful, or we have to move back to the point before this transaction. After each transaction we have to determine whether the conjunction of returned values are true, and \textbf{commit()} the transaction or \textbf{rollback()}. If a transaction fails, we can get the error message using the \textbf{get\_error\_message()} function. See Code \ref{lst:Example of the Database Transaction}.

\begin{lstlisting}[language=C++, caption={Example of the Database Transaction\label{lst:Example of the Database Transaction}}]
// database transaction example
bool success = true;
success = success && DataBaseHandler::some_operation();
success = success && DataBaseHandler::another_operation();
if (success) DataBaseHandler::commit();
else
{
    DataBaseHandler::rollback();
    QString error_message = DataBaseHandler::get_error_message();
    // do something
}
\end{lstlisting}

It should be noted that all member variables and functions inside the \textbf{DataBaseHandler} class are \textit{static}. This is because the \textbf{DataBaseHandler} is very frequently used. We found it would be very inefficient if we create a \textbf{DataBaseHandler} instance, initialize the database connector every time we want to access the database. To make everything static, we only have to initialize the database connector and connect to the database server once, and the same connector instance is used every time and everywhere. We must initialize the database connector upon initialization of the software.

\section{Database Design}
Based on an existing draft, we designed the database structure consisting of a number of tables. The tables, data fields and types of each table and relationships between tables have be been illustrated in the Figure \ref{fig:Database Structure}. 

\begin{figure}[htb]
\centering
\includegraphics[width = \textwidth]{database}
\caption{Database Structure\label{fig:Database Structure}}
\end{figure}

It is important that we do not hardcode e.g. project roles and register types in the software. Instead, they are defined in the tables starting with \textbf{def} and this provides us much flexibility. For example, we can add a new project role by simply adding a new row to the \textbf{def\_project\_role} table without modifying the software itself.
\chapter{Module Design and Implementation}
So far we have designed the software at the system level and learned background knowledge for implementation. In this chapter, we will design and implement each class.

\section{Main Window}
According to the requirements, we designed the main window as follows. The main window is not a unity as it seems to be. In fact, four classes are involved here: RegisterManager, which is the framework of the main window, ChipNavigator, ChipEditorView and DocumentEditorView. Their user interfaces are independently designed and compose the whole main window.

\begin{figure}[htbp]
\centering
\includegraphics[width = \textwidth]{mainwindow}
\caption{Main Window\label{fig:Main Window}}
\end{figure}

We have described in the chapter Architecture Design how these components work together. When the current item of the ChipNavigator changed, the ChipEditorView and the DocumentEditorView have to respond. To make this happen, we defined a slot function in the ChipNavigator so as to deal with the predefined currentItemChanged() signal of the tree widget. According to the level of the current tree widget item, we know whether it is the chip, a system block, a register, or a signal. Then, we define signal functions for each level and emit them accordingly. The logic can be described by the following pseudo code.

\begin{lstlisting}[language=C++]
void  on_treeWidgetBlock_currentItemChanged(QTreeWidgetItem *current, QTreeWidgetItem *previous)
{
    if (level(current) == LEVEL::CHIP) 
        emit(chip_clicked());
    if (level(current) == LEVEL::BLOCK) 
        emit(block_clicked(get_block_id(current)));
    if (level(current) == LEVEL::REGISTER) 
        emit(register_clicked(get_block_id(current), get_register_id(current)));
    if (level(current) == LEVEL::SIGNAL) 
        emit(signal_clicked(get_block_id(current), get_signal_id(current)));
}
\end{lstlisting}

To take care of the signal, we defined slot functions in the RegisterManager class and connect them to corresponding signals. In the slot functions, we first update the current block, register and signal ID. Then, we reset the document level. We pass the IDs to the ChipEditorView and DocumentEditorView, then let them display documents and chip information respectively.

\begin{lstlisting}[language=C++]
// register_manager.h
void on_chipNavigator_chip_clicked();
void on_chipNavigator_block_clicked(QString block_id);
void on_chipNavigator_register_clicked(QString block_id, 
        QString reg_id);
void on_chipNavigator_signal_clicked(QString block_id,
        QString sig_id);
\end{lstlisting}

\begin{lstlisting}[language=C++]
// register_manager.cpp
void RegisterManager::on_chipNavigator_register_clicked(QString block_id, QString reg_id)
{
    ui->docEditorView->set_doc_level(LEVEL::REGISTER);
    ui->docEditorView->set_block_id(block_id);
    ui->docEditorView->set_register_id(reg_id);
    if (ui->frameDoc->isVisible() && ui->actionDocEditor->isChecked()) 
        ui->docEditorView->display_documents();

    ui->chipEditorView->set_block_id(block_id);
    if (ui->frameChipEditor->isVisible()) 
        ui->chipEditorView->display_system_level_info(reg_id);
    current_block_id_ = block_id;
    current_reg_id_ = reg_id;
    current_sig_id_ = "";
}
\end{lstlisting}

We can edit the chip in the ChipEditorView by editing the chip name, adding, editing or removing system blocks, registers, signals, signal-register mappings and so on. In these cases, the chip navigator have to be updated. Similarly, we define the following signals in the ChipEditorView, and corresponding slot functions in the RegisterManager. When the chip is edited, a certain signal is emitted and information is passed to the corresponding slot function in the RegisterManager. In the slot function, a certain function of the chip navigator is called so as to update the tree widget.

\begin{lstlisting}[language=C++]
// chip_editor_view.h
void chip_basics_edited(QString chip_name, 
                  QString chip_owner, 
                  QString chip_owner_id, 
                  int register_width, 
                  int address_width, 
                  bool msb_first);
void block_added(QString block_id, 
                  QString block_name, 
                  QString block_abbr, 
                  QString responsible);
void block_removed(int row);
void block_modified(int row, 
                  QString block_name, 
                  QString block_abbr, 
                  QString responsible);
void block_order_exchanged(int from, int to);
void to_refresh_block(); 
\end{lstlisting}

The boundaries between the ChipNavigator, ChipEditorView and DocumentEditorView are adjustable. Users can also switch on or off the ChipEditorView or the DocumentEditorView. 

We design a menu bar for the main window like this. There are 4 menus as follows, each containing a number of menu entries, or actions.
\begin{itemize}
\item User Menu \\
User Management: to open a UserManagementDialog in which users can add or remove users. It is only accessible to database admins. \\
Change Password: to open a ChangePasswordDialog. For users to change their password. \\
Log Out: to log out of the current user account and to login with another one. 
\item Chip \\
New Chip: to open an EditChipDialog to create a new chip. \\
New Chip From: to create a new chip from an existing one. \\
Open Chip: to open an OpenChipDialog and open the selected chip. \\
Close Chip: to close the current chip and clear everything in the ChipNavigator, ChipEditorView and DocumentEditorView. \\
Naming: to open a NamingTemplateDialog to either edit the naming for registers or signals. \\
Freeze/Unfreeze: to freeze or unfreeze the current chip. \\
Chip Management: to open an OpenChipDialog and configure it to the management mode. Users can add or remove chips. It is only accessible to database admins. 
\item Export \\
SPI Source Code: to open an SPIGenerationDialog and generate the SPI interface. \\
Document: to open an DocumentGenerationDialog and generate the document. 
\item View \\
Chip Editor: to enable or disable the ChipEditorView. 
Document Editor: to enable the Document Editor page of the DocumentEditorView. \\
Document Preview: to enable the Document Preview page of the DocumentEditorView. The Document Editor and Document Preview are exclusively checkable. 
\end{itemize}

\subsection{ChipNavigator}
We have described the logic of the main window and its interaction with the ChipNavigator, the ChipEditorView and the DocumentEditorView. The ChipNavigator is quite simple. According to requirements, the structure of the chip is represented with a tree widget of for level: CHIP-BLOCK-REGISTER-SIGNAL. The root of the tree is the chip level, with its children being system blocks. Under each system block are registers. The children of each register are those signals mapped to it.

We implemented a search mechanism which might be very beneficial especially when the chip is complex. The users can simply type in the pattern to search, and the tree widget will highlight those components containing that pattern and make the other invisible. The search process is written in a single search() function. Whenever the input pattern changes, the search() function will be executed in response to the textChanged() signal from the tree widget.

\begin{figure}[htbp]
\centering
\includegraphics[width =0.5 \textwidth]{searchtree}
\caption{Chip Navigator: Search Example\label{fig:Chip Navigator: Search Example}}
\end{figure}

\subsection{Chip Editor View}
The ChipEditorView is one of the three major components of the main window. It displays information about the chip, and provides entrance to dialogs for editing the chip. According to the requirements, we designed the ChipEditorView as below.

\begin{figure}[htbp]
\centering
\includegraphics[width = \textwidth]{signalview}
\caption{Chip Editor: Signal Tab\label{fig:Chip Editor: Signal Tab}}
\end{figure}

\begin{figure}[htbp]
\centering
\includegraphics[width = \textwidth]{registerview}
\caption{Chip Editor: Register Tab\label{fig:Chip Editor: Register Tab}}
\end{figure}

Whether the chip level or the system level page is displayed depends on the current item of the navigator. If it is the root of the tree widget, then the RegisterManager class will call the display\_chip\_level\_info() function of the ChipEditorView, otherwise the display\_system\_level\_info(register\_id, signal\_id). If the current item is a register, then its register\_id will be passed to the function and signal\_id will be null. Likewise, if the item is a signal, then register\_id will be null. Depending on whether register\_id or signal\_id is null, the ChipEditorView will display the register page or the signal page, and set the selected register or signal in the register or the signal table as active. If the current item is a block, both register\_id and signal\_id are null, and the ChipEditorView simply displays the current page. This process can be described with the following simplified code.

\begin{lstlisting}[language=C++]
void ChipEditorView::display_system_level_info(const QString& reg_id, const QString& sig_id)
{
    // set to system level page
    ui->stackedWidgetChipEditor->setCurrentIndex(1);

    // set to register or signal page
    if (sig_id != "") ui->tabWidget->setCurrentIndex(0);
    else if (reg_id != "") ui->tabWidget->setCurrentIndex(1);

    // display either signals or registers
    if (ui->tabWidget->currentIndex() == 0) display_signals();
    else display_registers();
    
    // set current signal as active
    if (sig_id != "")
    {
        for (int row = 0; row < ui->tableSignal->rowCount(); row++)
            if (sig_id == ui->tableSignal->item(row, 0)->text())
            {
                ui->tableSignal->setCurrentCell(row, 0);
                break;
            }
    }
    // set current register as active
    else if (reg_id != "")
    {
        for (int row = 0; row < ui->tableRegister->rowCount(); row++)
            if (reg_id == ui->tableRegister->item(row, 0)->text())
            {
                ui->tableRegister->setCurrentCell(row, 0);
                break;
            }
    }
}
\end{lstlisting}

To add a system block, register, signal or a signal-to-register mappings we provide a push button under each table. After a click on the button a dialog of the corresponding class will be created and open. To remove an existing item we can click on the Remove buttons leading to execution of a corresponding slot function. The function will first remove the system block, register etc in the database, then remove it from the table. To edit an existing system block, register or signal, we can double click the table entry. The cellDoubleClicked() signal will then be emitted. We write corresponding slot functions to help us edit the items. The logic is actually very similar to adding a new item. A corresponding dialog will open, but the it will be initialized with data related to that item. The dialogs will be discussed in the later sections. 

For users' convenience, we designed a right-click context menu containing an Edit, Remove and Add action. Users do not have to click on the push buttons below each table. The menu also contains a Refresh action, which allows for reloading data from the database and refreshing the table. Once we right-click on the tables, a customContextMenuRequested() signal will be triggered. A slot function is then called to display the context menu in response to the signal.

\subsection{Document Editor View}
The DocumentEditorView is another central component of the software. It displays documents and provides entrance to editing the documents. According to the requirements, it should contain a table widget showing documents of the current item on the chip navigator. At the bottom of the DocumentEditorView there is an editing "dialog" in which we can add or edit a document. It is an instance of the EditDocumentDialog class, which is not a real dialog though, but a subclass of the generic QWidget. This is because the QDialog cannot be a child widget of another one. The dialog can be switched on or off. Like the ChipEditorView, we can add or edit documents by clicking on the Add button below or double clicking a table entry. The editing dialog will then open. If we finished editing and want to save the document, we click on the Ok button. The document will be saved to the database and the table above will be updated. Then, the editing area will be switched off. If we want to abandon the document we are editing, simply click on the Cancel button. 

\begin{figure}[htbp]
\centering
\includegraphics[width = \textwidth]{documentview}
\caption{Document Editor View\label{fig:Document Editor View}}
\end{figure}

Apart from the Add, Remove and Edit functions, we provided a copy and paste mechanism, which might be very useful for documentation editing. To achieve this, we maintain a member variable copied\_ of type QString. When we copy a document, the corresponding document type ID, document type, and the content to copy will be concatenated by a special delimiter and then be stored in this copied\_ variable. When we paste the copied content, the copied\_ variable will be parsed to fetch the document type ID, document type, and the content to paste. With this information, we can create a new entry in the document table and store this document to the database.

\section{Chip and Document Editor Dialogs}
To edit chip components and documents, we have to design different dialogs encapsulating both UI and business logics. These dialogs are named after EditSomethingDialog, such as EditChipDialog and EditSignalDialog. Their definitions follow the framework.

\begin{lstlisting}[language=C++]
class EditSomethingDialog : public QDialog
{
    Q_OBJECT

public:
    explicit EditSomethingDialog(QWidget *parent = nullptr);
    explicit EditSomethingDialog(QString something_id, bool enabled, QWidget *parent = nullptr);
    // other construction functions
    ~EditSomethingDialog();
    QString get_something_id();
    // Other public functions
    
    bool add_something();
    bool edit_something();

private:
    bool setup_ui();
    void accept(); // override QDialog::accept()
    bool sanity_check();
    /*
     Things to check
    */
    // Other private functions

    Ui::EditSomethingDialog *ui;
    const bool enabled_;
    const DIALOG_MODE mode_;
    QString something_id_;
    // Other private member variables 
};
\end{lstlisting}

The class has at least two construction functions, one for adding, the other for editing something. The DIALOG\_MODE mode\_ will correspondingly set to ADD or EDIT. Upon construction the dialog first setup the UI by calling the setup\_ui() function. If the dialog is constructed in the EDIT mode, it will pass the ID to the dialog and fill in the existing data into the widgets. Also, there is a parameter called enabled. If it is false, then the UI components will be disabled so users cannot put in anything to the dialog. Otherwise, users give their input to the dialog and click on the Ok button, and the accept() function will be executed. 

By default, if the EditSomethingDialog is executed as modal, the QDialog::accept() function will close the dialog and its exec() function will return QDialog::Accepted. However, we override the QDialog::accept() as below.

\begin{lstlisting}[language=C++]
void EditSomethingDialog::accept()
{
    if (!enabled_) return QDialog::reject();
    if (sanity_check()) return QDialog::accept();
}
\end{lstlisting}

If the dialog is not enabled, then just reject. This is equivalent to clicking on the Cancel button. Otherwise, do sanity check, which one of the system requirements. If the check is successful, then accept, else do nothing to the dialog. The purpose of the sanity check is to determine whether the input data is valid. The sanity\_check() function might be similar to

\begin{lstlisting}[language=C++]
void EditSomethingDialog::sanity_check()
{
    return check1() && check2() && check3();
}
\end{lstlisting}

We can determine what we want to check. The checking functions are defined like the example below. They basically follow the same pattern. They check something, prompt error message and return false at once if a check failed, or return true if all checks passed. In this example, we want to verify if the name is valid. The name must not empty or duplicate. If it is empty, a warning message will be prompted and the function will return false. Otherwise, if you are in the ADD mode, we have to fetch names from the database and check if the name already exists. If so, users are warned and the function returns false. If everything was fine, then return true. It is a little different in the EDIT mode. In that case, the original name actually exists in the database. If the current updated name equals the original name, the function just returns true, otherwise we check it as in the ADD mode.

\begin{lstlisting}[language=C++]
bool EditSomethingDialog::check_something_name()
{
    QString name = get_something_name();
    if (name == "")
    {
        QMessageBox::warning(this, "Add Something", "Something name must not be empty!");
        return false;
    }
    if (mode_ == DIALOG_MODE::EDIT && name == original_name_) return true;
    if (exists_in_database(name))
    {
        QMessageBox::warning(this, "Add Something", "Something" + name + "already exists!")
        return false;
    }
    return true;
} 
\end{lstlisting}

If the exec() function returns true, we then execute either the add\_something() or edit\_something() function. Then, we get what data we need from the dialog. We use the dialog to add or edit something as below.

\begin{lstlisting}[language=C++]
// add something
EditSomethingDialog add_somthing(this);
If (add_somthing.exec() == QDialog::Accepted && add_somthing.add_something())
{
    QString something_id = add_somthing.get_something_id();
    // Other data we want to get
    // Do something
}

// edit something existing
QString something_id; // got from somewhere
bool enabled; // got from somewhere
EditSomethingDialog edit_somthing(something_id, enabled, this);
If (edit_somthing.exec() == QDialog::Accepted && edit_somthing.edit_something())
{
    // Other data we want to get
    // Do something
} 
\end{lstlisting}

The most complex editor dialogs are the EditSignalDialog and EditDocumentDialog. The EditSignalDialog allows users to add or edit signals. Users have to put in the signal name, width or number of bits, signal type and whether to add a port for this signal. If the signal is writable, users have to put in its initial value. If the signal can be mapped to a register, or in other words, it is a register-signal, the dialog allows users to edit signal-register mappings. A signal can only be mapped to one type of registers.

In practice we found almost half of the signals are single-bit signals. In this case, the signal cannot be partitioned. It can be mapped to a certain bit of a register as a whole. In other cases where the signal have more bits, we might want to partition the signal, and map the signal partitions to certain registers partitions. We therefore design UI for two cases. In case of a single-bit signal, users just need to select a certain register and a certain bit of it to map the signal to. Users don't have to explicitly add a signal-register mapping to the partition list. In case of a multi-bit signal, however, users must determine the signal partition, select a register, and a register partition. Then, add the partition to the candidate list. Of course, users can also remove signal-register mapping from the candidate list. To make it more user friendly, we can add a new register by clicking on the button beside the register combo box.  

\begin{figure}[htbp]
\centering
\includegraphics[width = 0.8 \textwidth]{singlebitsignal}
\caption{Single-bit Signal\label{fig:Single-bit Signal}}
\end{figure}

\begin{figure}[htbp]
\centering
\includegraphics[width = 0.8 \textwidth]{multibitsignal}
\caption{Multi-bit Signal\label{fig:Multi-bit Signal}}
\end{figure}

Widgets in the EditSignalDialog have pretty strong interaction with each other. Whenever the signal width changed, the signal-register mappings might be invalid, thus we have to clear the candidate list. Also, we have to adjust the UI according to the signal width. If the signal type changed, we have to retrieve from the database the register types available to such signal type, and change items in the register type combo box. If the current register type changed, we have to clear the signal-register mapping candidate list, and retrieve registers of such type from the database. Whenever the candidate list changed, we also have to find available signal partitions. It is crucial that a signal bit cannot be mapped to more than one register bit, and the other way around. The register partition combo box changed according to the signal LSB and MSB combo boxes. The signal/slot technique provided by Qt makes these changes much more tractable.

The EditDocumentDialog allows users to add documents to the current item in the ChipNavigator, or edit existing ones. It is named after dialog but it is actually not one, because we want to incorporate it in the DocumentEditorView, and this is not possible for a dialog. Instead, it is a subclass of QWidget. However, we design it totally following the the pattern of the editor dialogs.

We defined three document types, text, image and table. Thus, the EditDocumentDialog is designed so that it can take in all document types. To make this possible the dialog has a stacked widget containing three pages, each handling a document type. In the previous section the DocuementEditorView example has already shown the page for text document. Below are image page and table page.

\begin{figure}[htbp]
\centering
\includegraphics[width = \textwidth]{imagedoc}
\caption{Image Document\label{fig:Image Document}}
\end{figure}

\begin{figure}[htbp]
\centering
\includegraphics[width = \textwidth]{tabledoc}
\caption{Table Document\label{fig:Table Document}}
\end{figure}

The text page has a QPlainTextEdit, in which we can edit multiple lines of text. In the image page, we can specify a caption of the image, select an image in a QFileDialog and get its path. It is also important to specify how large the image is in terms of page width. The number can vary between 0 and 1. In the table page, there is a QLineEdit holding the caption of the table. There is a QTableWidget in which we can edit the table contents. We can specify number of rows and columns for the table.

We will store the document in the database. The question is how to deal with different types of documents. The text document is simple. We just store the text in the database. For image documents, our solution is to concatenate the caption, width, and the path with a special delimiter, and store the concatenated string in the database. Similarly, for the table document, we concatenate the caption, row count, column count, and all cell contents. In this solution, of course the image caption, table caption and cell contents must not contain the delimiter. It would be a good idea to use a very rare string as the delimiter. We can easily restore the documents in a reverse way.

It's inevitable that we might have to include math equations in LaTeX format in the documents. In this case, of course we can type LaTeX code "blindly". It would be very nice if we have a preview of the LaTeX math equations. In this way, the input is more visual, and we can check if the code is correct. We also want to have a preview of the table and image. It's important that table cell contents can also contain LaTeX math equations.

To implement the preview, we might want to compile the LaTeX source code to generate a PDF document. Then, we must find a way to display the PDF document in our software. This is not a good idea. First, users have to install a LaTeX compiler on their machines and it must be accessible to the software. Second, Qt does not have a PDF render. We were able to find third-party renders, but they are either of low quality or hard to incorporate in our software.

However, we found MathJax, a JavaScript display engine for mathematics that works in web browsers. With MathJax included in the HTML web page, LaTex source code can be simply inserted in the HTML code, and browsers can render it properly. We use MathJax with the following template.

\begin{lstlisting}[language=HTML]
<!-- HTML template -->
<html>
<head>
    <script type="text/x-mathjax-config"> MathJax.Hub.Config({
        tex2jax: {inlineMath: [['$','$'], ['\\(', '\\)']]}});</script>
    <script type="text/javascript" src="MATHJAX_ROOT/MathJax.js"></script>
</head>
<body>
HTML_CONTENT    <!-- write your HTML code here -->
</body>
</html>
\end{lstlisting}

Inspired by this, we can display text, image and table documents in a web browser. For text documents it is quite simple. We just replace HTML\_CONTENT in the HTML template with the text in the text editor. For image and table documents, we created a template respectively. Using the image or table template, we generate image or table HTML source code, and replace HTML\_CONTENT in the HTML template.

\begin{lstlisting}[language=html]
<!-- HTML image template -->
<style>
    figure img {display: block; margin-left: auto; margin-right: auto;}
    figure figcaption { text-align: center;}
</style>
<center>
<figure>
   <img style='width: WIDTH; object-fit: contain' src="IMAGE" alt="CAPTION"/>
    CAPTION
</figure>
</center>
\end{lstlisting}

\begin{lstlisting}[language=html]
<!-- HTML table template -->
<style>
    table {border-top: 1px solid black; border-bottom: 1px solid black}
    th {border-bottom: 1px solid black}
    caption#tab {caption-side: CAPTION_POS}
</style>
<center>
<table>
TABLE
</table>
</center>
\end{lstlisting}

Fortunately, Qt provides a web engine based on Chromium. Thus, we are able to make a preview based on the web engine. The preview is intrinsically a web browser. To display documents of any types, we generate the HTML source code and fill it into the browser. In this way, the document editor can display the real time preview of the text we are editing.

Another effort we made is to implement an auto-completer. This is especially useful for editing text documents. The completer we designed takes all system block names, system block abbreviations, register names and signal names as candidates. Besides these, the software allows users to extend the word list by simply add txt files containing candidate words to the \textbf{completion} directory under the root directory of the software.

\begin{figure}[htbp]
\centering
\includegraphics[width = 0.5\textwidth]{completer}
\caption{Document Completer\label{fig:Document Completer}}
\end{figure} 

\section{User Account, Authentication and Password}
For data security reasons, the software is required to provide user access control and authentication. User access control has two levels, the software level or the database level, and the project level. If a user account is created, the user can log in to the software with this account. However, it does not necessarily mean he or she has access to a certain chip. This is controlled by the project level permissions.

All users must have access to the database, thus, we need to establish database accounts for them. MySQL provides a very good user account management. We can grant different permissions for different user. However, the permissions we want to implement are more complex. For example, in a certain table, a user might have access to row A but not to row B. This is not possible with MySQL. Thus, we have to manage permissions using our software. In this background, we distinguished the software account and database account. We store user accounts in the database. Users log in to the database using a shared database account, retrieve user account information and determine if they have access to the software. Similarly to the user account, chip designers, database level and project level permissions are also stored in the database.

We define different database permissions allowing users to add or remove users, add or remove projects. A special permission is reserved for the database/software admins, so they have full access to all projects. This means that they can do anything on any projects, even if they are in the designer list of those projects. We predefined two database roles, the admin and standard user. However, we can extend the database roles without modifying the software. Similarly, we defined different project permissions allowing chip designers to add system blocks, remove the system blocks they are in charge of, read all system blocks, add or remove chip designers.

To make use of the database and project permissions, we designed an Authenticator class. To make block level authentication more manageable, we also defined block permissions here. They are determined by project permissions and the responsible person of the system block. The Authenticator class allows users to set database permissions, project permissions and block permissions. The permissions are stored member variables db\_permissions\_, project\_permissions\_, block\_permissions\_ of data type int. Each bit of the permissions variable represents a permission, which is defined in the enumeration DATABASE\_PERMISSIONS, PROJECT\_PERMISSIONS, and BLOCK\_PERMISSIONS. We set or get permissions using bit operations. For example, to set ADD\_USER permission, we make the 0th bit of db\_permissions\_ be 1. To get ADD\_USER permission, we get the value of the 0th bit of db\_permissions\_.

\begin{lstlisting}[language=C++]
class Authenticator
{
public:
    enum DATABASE_PERMISSIONS
    {
        ADD_USER = 1 << 0,
        // ...
        FULL_ACCESS_TO_ALL_PROJECTS = 1 << 4
    };
    // enum PROJECT_PERMISSIONS
    // enum BLOCK_PERMISSIONS
    
    Authenticator(const QString& db_role_id, const QString& project_role_id);
    Authenticator();

    void set_database_permissions(const QString& db_role_id);
    void set_project_permissions(const QString& project_role_id);
    void set_project_permissions(bool setting);
    void set_block_permissions(bool setting);
    void freeze(bool frozen=true);

    // database permissions
    bool can_add_user() const;
    // ...
    bool can_fully_access_all_projects() const;

    // project permissions
    bool can_add_block() const;
    // ...
    bool can_fully_access_all_blocks() const;
    
    // block permissions
    bool can_add_signal() const;
    // ...
    bool can_edit_document() const;
    
    bool frozen() const;

    void clear_database_permission();
    void clear_project_permission();
    void clear_block_permission();
    void clear_all_permission();

private:
    int db_permissions_ = 0, project_permissions_ = 0, block_permissions_;
    bool frozen_;
};
\end{lstlisting}

We provide a convenient way for managing user accounts using the UserManagementDialog. It allows database admins to add or remove users. To add a user, we need the CreateUserDialog. We specify the username, initial password and select a database role. Like the editor dialogs, we have to check whether the username is valid. After creation of a user account, users can then log in to the software using the account. To remove a user, however, it will be tricky. The users to be deleted might be the owner of some chips. They might be a designer in some chips, and responsible for some system blocks. In this case, our solution is to replace the users in those chips with the database admin who is deleting these users. Users can change their passwords easily using the ChangePasswordDialog.

On the project level, designers can be added to a chip with the EditChipDesignerDialog, where a user and a project role are selected. Before adding a designer, the dialog shall check whether the selected user already exists in the current chip. We can also delete a chip designer. Like removing a user, the software has to replace the designers in this chip with a project admin.

As previously described, the software account and database account are not the same. In our solution, all software users share a single database account. The problem is that we don't want users to know the password to the database to prevent them from bypassing the software and directly work on the database. This might cause serious data security problems.

Our solution is to introduce an encryptor and decryptor. We designed an encryption class using the AES algorithm. The encryptor can encrypt a plain password with a key. With this key and encrypted password, the decryptor can generate the original plain password. Only the admins know the password. They generate encrypted password, and give users the key and the encrypted password instead of the original password. We made a tool for admins to make encrypted passwords easily.

\begin{figure}[htbp]
\centering
\includegraphics[width = 0.8\textwidth]{encryptor}
\caption{Encryptor Dialog\label{fig:Encryptor Dialog}}
\end{figure} 

\section{Signal and Register Naming}
In practice, we name signals and registers after a certain pattern. For example, a register might be named after \$\{BLOCK\_ABBREVIATION\}\_MAIN\_REG. In this example, the block abbreviation is prepended to the central name MAIN, and a constant string "REG" is appended to the name. The block abbreviation, the what we call given name, and the suffix "REG" are concatenated with an underscore.

The problem is that, if we have to change a system block's abbreviation, then all registers have to be renamed. This introduces more database operations and unreliability. Our solution is simple: we only store the given name in the database. However, the extended name is generated with the naming template and the given name. 

Thus, we designed a NamingTemplate class in which we are able to generate the extended name with the given name, and the other way around. The idea is that we maintain a naming template like \$\{BLOCK\_ABBREVIATION\}\_\$\{GIVEN\_NAME\}\_REG. The naming template contains variables like \$\{Variable\} and constant. The given name variable must always be in the template. To generate the extended name, we simply replace the variables with the abbreviation name of the current system block, and the given name of the current register, which is retrieved from the database. Similarly, we can also generate the given name given an extended name.

To change the naming of signals or registers, we don't have to do anything on the database. We simply make a new naming template and refresh the UI. We designed a very friendly dialog to change the naming template. Users just need to add keywords or constants, and specify a delimiter.

\begin{figure}[htbp]
\centering
\includegraphics[width = 0.8\textwidth]{namingtemplatedialog}
\caption{Naming Template Dialog\label{fig:Naming Template Dialog}}
\end{figure} 

\section{SPI Interface Generation}
With all the chip definitions and documents, our final goals are generating the SPI configuration interface and a LaTeX documentation. We might want to generate the interface in different languages, so far we only implement VHDL. Despite of this, during design and implement, we already leave flexibility for extensions.

The VHDL SPI interface consists of two files, the interface package, where basic information about the chip is defined, and the interface, where the VHDL interface and behavior is defined. Users can define a package and a interface template with predefined markers. The software will replace the markers with what it generated from the database.

\begin{lstlisting}[language=VHDL]
-- sample VHDL package template
LIBRARY ieee;
  USE ieee.std_logic_1164.all;
  USE ieee.numeric_std.all;
  
PACKAGE interface_package IS

  constant address_width      : integer := 12;
  constant byte_size          : integer := 8;
  constant reg_width           : integer := 8;
  -- FIFO
  constant byte_counter_width : integer := 10;
  constant fifo_size          : integer := 2**byte_counter_width;
  constant packet_size        : integer := 2**byte_size;

   -- REGISTER ADDRESSES
  @PACKAGE_ADDRESSES    -- to be completed by EDA tool

  -- REGISTER_INITIAL VALUES
  @PACKAGE_INITS    -- to be completed by EDA tool
  
END PACKAGE interface_package;
\end{lstlisting}

In the package template there are two makers, one for register address definitions, and the other for initial values definitions for writable registers. Since we allocated each system block a start address, it's easy to infer the address for each register. The initial values of each writable registers are tricky. We have to compute compute initial values for each register bit by looking for the corresponding signal bit it is mapped to, and the initial signal value. Then, the initial value of the registers are obtained by concatenating the initial values of each register bit.

\begin{lstlisting}[language=VHDL]
-- sample VHDL package template
library IEEE;
use IEEE.STD_LOGIC_1164.ALL;
use IEEE.NUMERIC_STD.ALL;

LIBRARY work;
use work.interface_package.ALL;

entity interface is
  generic
  (
    data_width : integer := 12
  );
  port
  (
    resetn                    : in  std_logic;
    clk_main                  : in  std_logic;
    enable                    : in  std_logic;
    
    -- REGISTER SIGNAL PORT DEFINITIONS      
    @PORT_DEFINITIONS
 
  );
end interface;

architecture Behavioral of interface is

  -- INTERNAL SIGNALS FOR INTERFACE FUNCTION
  signal addr_buffer            : std_logic_vector(address_width-1 downto 0);
  signal reg_read_buffer        : std_logic_vector(byte_size-1 downto 0);
  signal byte_ext2asc           : std_logic_vector(byte_size-1 downto 0);
  
  -- INTERNAL SIGNALS NOT CREATED BY THE AUTOMATIC EXPORT
  signal SM_INVERT_FIFO_CLK     : std_logic;  
  signal SM_FIFO_SPI_RESETB     : std_logic;
  signal CHIP_ID_VALUE		: std_logic_vector(16-1 downto 0);
  signal CHIP_REGPAGE_CTRL	: std_logic_vector(2-1 downto 0);

  -- REGISTER DEFINITIONS
  @REGISTER_DEFINITIONS

begin 
  
  reg_write_proc: PROCESS(clk_main, resetn, enable)
  BEGIN
    if resetn = '0' then
      -- Registers
      @REGISTER_INIT

      new_spi_reg_write       <= '0';
    else
      if rising_edge(clk_main) and enable = '1' then
        if current_state = S_DATW and byte_ext2asc_rising = '1' then
          new_spi_reg_write  <= '1';
          case addr_buffer is
            @REG_WRITE_ACCESS

            when others                         => null;
          end case;
        else
          new_spi_reg_write <= '0';
        end if;
      end if;
    end if;
  END PROCESS reg_write_proc;
  
  reg_read_proc: PROCESS(clk_main, resetn, enable)
  BEGIN
    if resetn = '0' then
      reg_read_buffer  <= (others => '0');
      new_spi_reg_read            <= '0';
    else
      if rising_edge(clk_main) and enable = '1' then
        if current_state = S_DATR then
          new_spi_reg_read        <= '1';
          case addr_buffer is
            @REG_READ_ACCESS

            when others                         => null;
          end case;
        else 
          new_spi_reg_read        <= '0';
          reg_read_buffer         <= (others => '0');
        end if;
      end if;
    end if;
  END PROCESS reg_read_proc;
  
  @SIGNAL_ASSIGNMENTS
  
end Behavioral;
\end{lstlisting}

The interface is more complex. It consists of two code blocks, the entity definition, and the behavioral definition. In the entity definition, the software shall complete the port definitions. All port signals should be included. Then, in the behavioral definition, the software first complete the register definitions. After after clock, the interface should do read or write process. It reads data from external. According to the address buffer, the data should be written to a certain register. Otherwise it reads data from a certain register, and write to external. Finally, the software should assign values. In case of a read-only register, the software shall assign to it values from the signals it is mapped to. In case of a writable register, the software shall assign values to the signals it is mapped to. Paged registers are dealt with in a special way.

We designed a VHDLGenerator class. The generation process of the interface is described as the following pseudo code. Generation of the package follows the same structure.

\begin{lstlisting}[language=C++]
QString VHDLGenerator::generate_interface() const
{
    QString ports,
            register_definitions, register_inits,
            register_write_accesses, register_read_accesses,
            readonly_register_assignments,
            control_signal_assignments,
            paged_register_assignments;
            
    QVector<QString> blocks = read_system_blocks();
    for (const auto& block : blocks)
    {
        QVector<QString> interface_block = generate_interface_block(block);
        ports += interface_block[0];
        register_definitions += interface_block[1];
        register_inits += interface_block[2];
        register_write_accesses += interface_block[3];
        register_read_accesses += interface_block[4];
        readonly_register_assignments += interface_block[5];
        control_signal_assignments += interface_block[6];
        paged_register_assignments += interface_block[7];
    }
    QString interface = get_interface_template();
    interface.replace(marker_ports, ports);
    interface.replace(marker_register_definitions, register_definitions);
    interface.replace(marker_register_init, register_inits);
    interface.replace(marker_register_write, register_write_accesses);
    interface.replace(marker_register_read, register_read_accesse);
    interface.replace(marker_signal_assignment, signal_assignments);
      
    return interface;
}

QVector<QString> VHDLGenerator::generate_interface_block(block) const
{
    QVector<QString> signals = read_signals(block);
    QVector<QString> registers = read_registers(block);
    QString ports,
            register_definitions, register_inits,
            register_write_accesses, register_read_accesses,
            readonly_register_assignments, 
            control_signal_assignments, 
            paged_register_assignments;
            
    for (const auto& register : registers)
    {
        register_definitions += generate_register_definition(register);
        register_read_accesses += generate_reading_register(register);
        if (writable(register))
        {
            register_inits += generate_initializing_register(register);
            register_write_accesses += generate_writing_register(register);
        }
        else
        {
            readonly_register_assignments += generate_assigning_value_to_readonly_register(register);
        }
        if (paged_register(register))
            paged_register_assignments += generate_assigning_value_to_paged_register(register);
    }
    
    for (const auto& signal : signals)
    {
        if (port_signal(signal)) 
            ports += generate_port_definition(signal);
        if (control_signal(signal)) 
            control_signal_assignments += generate_assigning_value_to_constrol_signal(signal);
    }
    
    return {ports,
            register_definitions, register_inits,
            register_write_accesses, register_read_accesses,
            readonly_register_assignments, 
            control_signal_assignments, 
            paged_register_assignments};
}
\end{lstlisting}

We made a SPIGenerationDialog class for users to configure the generation. The designing pattern is similar to the previous editor dialogs. Since the generated code must be directly usable, we have to make sure there are no errors. Sanity checks must be done.

\begin{figure}[htbp]
\centering
\includegraphics[width = 0.8\textwidth]{spigenerationdialog}
\caption{SPI Generation Dialog\label{fig:SPI Generation Dialog}}
\end{figure} 

\section{Document Generation}
So far chip designers have added documents to each system block, register and signal, and also the chip itself as well. The software will compose them and generate a complete documentation.

The documentation is well structured. Thus, we designed a documentation generator following this structure.

\begin{figure}[htbp]
\centering
\includegraphics[width = 0.8\textwidth]{docstructure}
\caption{Structure of the Documentation\label{fig:Structure of the Documentation}}
\end{figure} 

\begin{lstlisting}[language=C++]
// document_generator.h

// functions to generate the LaTeX documentation
// parameters are omitted
QString generate_tex_document();
QString generate_chip_level_tex_document();
QString generate_block_level_tex_document();
QString generate_register_level_tex_document();
QString generate_register_bit_table_tex();
QString generate_register_signal_bullets_tex();

// static functions to generate single documents into LaTeX
static QString generate_text_tex(const QString& text);
static QString generate_image_tex(const QString& caption, const QString& width, const QString& path);
static QString generate_table_tex(const QString& caption, const QVector<QVector<QString> >& cells);
\end{lstlisting}

In previous sections, we have introduced a document preview for the document editor.  It is basically a web browser rendering HTML web page with LaTeX source code with the help of MathJax. Since the software is already able to display a single document, why not implement a full documentation generation function in HTML format, and make a preview of the complete documentation?

Inspired by this, we implemented the HTML generation functions in a similar way to LaTeX. We can generate a full HTML documentation, and let the software display it. Users can also choose to generate the HTML document. In some cases it would be a better choice than LaTeX, for example, if the Chair of IAS wants to make an internal wiki.

\begin{figure}[htbp]
\centering
\includegraphics[width = \textwidth]{htmldocexample}
\caption{Document Preview in HTML Format\label{fig:Document Preview in HTML Format}}
\end{figure} 

Similarly, we created a DocumentGenerationDialog for users to select the document format, configure the documentation and select which system blocks they want to generate.

\begin{figure}[htbp]
\centering
\includegraphics[width = 0.8\textwidth]{docgenerationdialog}
\caption{Documentation Generation Dialog\label{fig:Documentation Generation Dialog}}
\end{figure} 

\chapter{Integration, Testing and Deployment}
So far we have completed the design and implementation process. During implementation, we adopted an incremental strategy. We tested each module right after it is implemented. Now, we want to integrate the modules and test the software as a whole.

Therefore, we made a beta version and delivered it to the users, which allows them to participate in the testing phase by using the software solution on real-world projects. During this process, some errors are reported and we were able to solve them. Users also gave suggestions on the software and we were able to improve it so as to better satisfy users' requirements. 

To test the software, users have to install a MySQL C++ connector and the whole Qt framework on their machines. Then, configure the linked library and compile the software locally. This is fine with testing. We still have to make a deployable software so that users don't have to install the dependencies.

The theory is that the external dependencies including the Qt Modules and MySQL connector are not incorporated in our source code. Instead, they are used in form of shared libraries. During compilation, the executable knows where to find the shared libraries it depends on. During execution, the executable finds the libraries and load them into the computer memory. This is called dynamic linking. If the executable cannot find the libraries it needs, it will crash.

Thus, we need to manually collect all Qt frameworks and the MySQL connector the software depends on. We package the libraries with the executable under the same directory, and make the executable find the libraries under this directory during runtime. We can either specify an rpath for the linker during compilation, or relink the executable after it is compiled.

It is not an easy thing to do. Fortunately, there are tools helping us with deployment. The Mac version of Qt itself provides a tool called macdeployqt. It is very simple to use. We just compile the software as usual. Then, we use macdeployqt to deploy the software. We only need to specify the path to the executable. The tool will do everything for us: make a software directory, find the dependency libraries, copy and paste those libraries under the directory, and relink the executable.

To deploy the software on a Linux system, although Qt does not provide a tool like macdeployqt, the community developed a tool named linuxdeployqt which basically follows macdeployqt. With the help of linuxdeployqt, we successfully deployed the software to the Chair of IAS.

% save the arabic Number Page Values to Continue with them later
\newcounter{page_tmp_arabic}
\setcounter{page_tmp_arabic}{\value{page}}

% here beginns the backmatter% 
% set numbering to roman or Roman and continue with the romnan poage number
% Start group that doesn't contain the name Chapter in their headers
% Putting in braces the \backmatter command together with the chapters that will not include
% the chapter name, avoids that other parts of the document omit this name, like the appendices.
{
\backmatter
\pagenumbering{roman}
\setcounter{page}{\value{page_tmp_roman}}
\pagestyle{scrheadings}
\printbibliography
\cleardoublepage
}

% now we have to edit the Page layout.
%\lehead{\textsc{Abbriviations}}  % linke aussenseite Kapitel mit numer in Header
%\rohead{} % Rechte Aussenseite Section mit nummer in header
%\renewcommand\chaptermarkformat{} % entfern Kapitel nummer in header
%\renewcommand\sectionmarkformat{} % entfernsection nummer in header
%\ofoot{ {\pagemark} } %seitenzahlen an der Aussenseite der Fu

\appendix
% Requested to be with continued latin numbers
\pagenumbering{arabic}
\setcounter{page}{\value{page_tmp_arabic}}

% TODO REMOVE the following 3 lines in your document
%\newbool{UseIncludeStandaloneForExamples}
%\setbool{UseIncludeStandaloneForExamples}{false}
%\input{checkAllStandaloneFiles.tex}
\pagestyle{scrheadings}
%TODO INCLUDE your presentations slides paths to the settings and remove comment of next line
\cleardoublepage % do NOT remove

\ifbool{PrintVersion} %
{% Print Version - has bcor and 
		
	% Chapter and Section Page
	\ifdefined\firstPresentation
	\includepdf[
	pagecommand={
		\ifbool{IsEnglish}{  
			\pagestyle{scrheadings}
			\chapter{Presentation Slides}
			\section{First \IASSubject \ Presentation} 
		}{  \pagestyle{scrheadings}
		\chapter{Vortragsfolien}
		\section{Erste \IASSubject s-Präsentation}
	}  
	},
	%link=false,
	noautoscale,
	fitpaper=false,
	frame=true, 
	nup=2x2,
	pages={1-4},
	width=0.48\textwidth,
	offset= -4mm -23.8mm, 
	delta= 6.4mm 6.4mm, 
	%trim=0mm 0mm 0mm 0mm,
	%clip,
	]{\firstPresentation}
	% Only PDF  	 
	\includepdf[
	pagecommand={
		\ifbool{IsEnglish}{  
			\pagestyle{scrheadings}
		}{  \pagestyle{scrheadings}
	}  
	},
	%link=false,
	noautoscale,
	fitpaper=false,
	frame=true, 
	nup=2x3,
	pages={5-last},
	width=0.48\textwidth,
	offset= -4mm 0mm, 
	delta= 6.4mm 6.4mm, 
	%trim=0mm 0mm 0mm 0mm,
	%clip,
	]{\firstPresentation}
	\fi

	\ifdefined\secondPresentation
		%  Section Page
		\includepdf[
		pagecommand={
			\ifbool{IsEnglish}{  
				\pagestyle{scrheadings}
				\section{Second \IASSubject \ Presentation} 
			}{  \pagestyle{scrheadings}
				\section{Zweite \IASSubject s-Präsentation}
			}  
		},
		%link=false,
		noautoscale,
		fitpaper=false,
		frame=true, 
		nup=2x3,
		pages={1-6},
		width=0.48\textwidth,
		offset= -4mm 0mm, 
		delta= 6.4mm 6.4mm, 
		%trim=0mm 0mm 0mm 0mm,
		%clip,
		]{\secondPresentation}
		% Only PDF  	 
		\includepdf[
		pagecommand={
			\ifbool{IsEnglish}{  
				\pagestyle{scrheadings}
			}{  \pagestyle{scrheadings}
			}  
		},
		%link=false,
		noautoscale,
		fitpaper=false,
		frame=true, 
		nup=2x3,
		pages={7-last},
		width=0.48\textwidth,
		offset= -4mm 0mm, 
		delta= 6.4mm 6.4mm, 
		%trim=0mm 0mm 0mm 0mm,
		%clip,
		]{\secondPresentation}
	\fi 	
	
	%  Section Page
	\ifdefined\finalPresentation
	\includepdf[
	pagecommand={
		\ifbool{IsEnglish}{  
			\pagestyle{scrheadings}
			\section{Final \IASSubject \ Presentation} 
		}{  \pagestyle{scrheadings}
		\section{Finale \IASSubject s-Präsentation}
	}  
	},
	%link=false,
	noautoscale,
	fitpaper=false,
	frame=true, 
	nup=2x3,
	pages={1-6},
	width=0.48\textwidth,
	offset= -4mm 0mm, 
	delta= 6.4mm 6.4mm, 
	%trim=0mm 0mm 0mm 0mm,
	%clip,
	]{\finalPresentation}
	% Only PDF  	 
	\includepdf[
	pagecommand={
		\ifbool{IsEnglish}{  
			\pagestyle{scrheadings}
		}{  \pagestyle{scrheadings}
	}  
	},
	%link=false,
	noautoscale,
	fitpaper=false,
	frame=true, 
	nup=2x3,
	pages={7-last},
	width=0.48\textwidth,
	offset= -4mm 0mm, 
	delta= 6.4mm 6.4mm, 
	%trim=0mm 0mm 0mm 0mm,
	%clip,
	]{\finalPresentation}
	\fi

}%
{%  Online Version - no bcor, sett text in the middle of the page	
	% Chapter and Section Page
	\ifdefined\firstPresentation
	\includepdf[
	pagecommand={
		\ifbool{IsEnglish}{  
			\pagestyle{scrheadings}
			\chapter{Presentation Slides}
			\section{First \IASSubject \ Presentation} 
		}{  \pagestyle{scrheadings}
		\chapter{Vortragsfolien}
		\section{Erste \IASSubject s-Präsentation}
	}  
	},
	%link=false,
	noautoscale,
	fitpaper=false,
	frame=true, 
	nup=2x2,
	pages={1-4},
	width=0.48\textwidth,
	offset= 0mm -23.8mm, 
	delta= 6.4mm 6.4mm, 
	%trim=0mm 0mm 0mm 0mm,
	%clip,
	]{\firstPresentation}
	% Only PDF  	 
	\includepdf[
	pagecommand={
		\ifbool{IsEnglish}{  
			\pagestyle{scrheadings}
		}{  \pagestyle{scrheadings}
	}  
	},
	%link=false,
	noautoscale,
	fitpaper=false,
	frame=true, 
	nup=2x3,
	pages={5-last},
	width=0.48\textwidth,
	offset= 0mm 0mm, 
	delta= 6.4mm 6.4mm, 
	%trim=0mm 0mm 0mm 0mm,
	%clip,
	]{\firstPresentation}
	\fi
	
	\ifdefined\secondPresentation
		%  Section Page
		\includepdf[
		pagecommand={
			\ifbool{IsEnglish}{  
				\pagestyle{scrheadings}
				\section{Second \IASSubject \ Presentation} 
			}{  \pagestyle{scrheadings}
			\section{Zweite \IASSubject s-Präsentation}
		}  
		},
		%link=false,
		noautoscale,
		fitpaper=false,
		frame=true, 
		nup=2x3,
		pages={1-6},
		width=0.48\textwidth,
		offset= 0mm 0mm, 
		delta= 6.4mm 6.4mm, 
		%trim=0mm 0mm 0mm 0mm,
		%clip,
		]{\secondPresentation}
		% Only PDF  	 
		\includepdf[
		pagecommand={
			\ifbool{IsEnglish}{  
				\pagestyle{scrheadings}
			}{  \pagestyle{scrheadings}
		}  
		},
		%link=false,
		noautoscale,
		fitpaper=false,
		frame=true, 
		nup=2x3,
		pages={7-last},
		width=0.48\textwidth,
		offset= 0mm 0mm, 
		delta= 6.4mm 6.4mm, 
		%trim=0mm 0mm 0mm 0mm,
		%clip,
		]{\secondPresentation}
	\fi
	
	%  Section Page
	\ifdefined\finalPresentation
	\includepdf[
	pagecommand={
		\ifbool{IsEnglish}{  
			\pagestyle{scrheadings}
			\section{Final \IASSubject \ Presentation} 
		}{  \pagestyle{scrheadings}
			\section{Finale \IASSubject s-Präsentation}
		}  
	},
	%link=false,
	noautoscale,
	fitpaper=false,
	frame=true, 
	nup=2x3,
	pages={1-6},
	width=0.48\textwidth,
	offset= 0mm 0mm, 
	delta= 6.4mm 6.4mm, 
	%trim=0mm 0mm 0mm 0mm,
	%clip,
	]{\finalPresentation}
	% Only PDF  	 
	\includepdf[
	pagecommand={
		\ifbool{IsEnglish}{  
			\pagestyle{scrheadings}
		}{  \pagestyle{scrheadings}
		}  
	},
	%link=false,
	noautoscale,
	fitpaper=false,
	frame=true, 
	nup=2x3,
	pages={7-last},
	width=0.48\textwidth,
	offset= 0mm 0mm, 
	delta= 6.4mm 6.4mm, 
	%trim=0mm 0mm 0mm 0mm,
	%clip,
	]{\finalPresentation}
	\fi
}

\cleardoublepage % do NOT remove



\end{document}